% Leçon 1 — Vocabulaire et compréhension de texte : situation de la femme (égalité homme-femme) — Chinois Tle A — Cours signé OnBuch+
% Programme officiel MINESEC. Compiler avec : tectonic ZH01-vocabulaire-situation-de-la-femme-egalite-homme-fe.tex
\newcommand{\DOCMATIERE}{Chinois}
\newcommand{\DOCNIVEAU}{Tle A}
\newcommand{\DOCMODULE}{Module 1 : Vie familiale et intégration sociale}
\newcommand{\DOCLECON}{ZH01}
\newcommand{\DOCTITRE}{Leçon 1 — Vocabulaire et compréhension de texte : situation de la femme (égalité homme-femme)}
% OnBuch+ — préambule « cours signés » ENRICHI (compilé avec tectonic / XeLaTeX)
% Système visuel « Pop » d'OnBuch+ : fond crème (jamais blanc), bordures encre
% épaisses, ombres « dures » décalées, titres Archivo Black en capitales.
% Version MATHÉMATIQUES : graphes (pgfplots/tikz), boîtes pédagogiques
% (définition, propriété, méthode, attention, le savais-tu, exercice résolu,
% exercice, TP, bilan), page de garde et signature OnBuch+ ancrée en bas.
%
% Macros attendues dans main.tex AVANT \input{preamble} :
%   \DOCMATIERE  \DOCNIVEAU  \DOCMODULE  \DOCLECON (numéro)  \DOCTITRE
\documentclass[11pt]{article}

\usepackage{fontspec}
\usepackage[french]{babel} % français = langue principale ; le chinois via \zh{..} / textebox (xeCJK)
\usepackage{xeCJK}
\usepackage{pifont} % \ding{51} \ding{55} utilisés par les modèles
\usepackage[a4paper,margin=2cm,top=3.4cm,bottom=2.8cm,headheight=1.4cm,footskip=1.3cm]{geometry}
\usepackage{amsmath,amssymb,amsfonts}
\usepackage{mathtools} % \xmapsto, \coloneqq... souvent utilisés par les modèles
\usepackage{cancel} % simplifications barrées (bilans de forces)
% \abs{z} (module, valeur absolue) et \norm{u} (norme), utilisés par les modèles
\providecommand{\abs}{}\let\abs\relax\DeclarePairedDelimiter{\abs}{\lvert}{\rvert}
\providecommand{\norm}{}\let\norm\relax\DeclarePairedDelimiter{\norm}{\lVert}{\rVert}
\usepackage[table]{xcolor} % [table] : fournit aussi \rowcolors (alternance de lignes)
\usepackage{pagecolor}
\usepackage{tikz}
\usetikzlibrary{arrows.meta,positioning,calc,shapes.geometric,shapes.misc,shapes.arrows,decorations.pathmorphing,decorations.pathreplacing,decorations.markings,patterns,shadows,mindmap,trees,fit,backgrounds,matrix,intersections,angles,quotes}
\usepackage{pgfplots}
\usepackage[version=4]{mhchem} % équations nucléaires \ce{^{235}_{92}U}
\usepackage[european,siunitx]{circuitikz} % schémas électriques
\pgfplotsset{compat=1.18}
\usepgfplotslibrary{fillbetween,groupplots} % aires entre courbes (calcul intégral)
\usepackage{enumitem}
\usepackage{fancyhdr}
% [most] charge toutes les bibliothèques tcolorbox (listings, minted, raster,
% documentation...) et gonfle inutilement la mémoire TeX sur un document long
% (~300 boîtes) : on ne charge que ce qui est réellement utilisé.
\usepackage[skins,breakable]{tcolorbox}
\usepackage{graphicx}
\usepackage{colortbl}
\usepackage{booktabs}
\usepackage{tabularx}
\usepackage{diagbox} % case d'en-tête diagonale des tableaux à double entrée
% Colonnes extensibles centrée / gauche / droite (C, L, R), souvent utilisées par les modèles
\newcolumntype{C}{>{\centering\arraybackslash}X}
\newcolumntype{L}{>{\raggedright\arraybackslash}X}
\newcolumntype{R}{>{\raggedleft\arraybackslash}X}
\usepackage{array}
\usepackage{multicol}
\usepackage{multirow}
\usepackage{siunitx}
% parse-numbers=false : accepte aussi \SI{2,5 \times 10^{-3}}{...}
\sisetup{locale=FR,output-decimal-marker={,},group-minimum-digits=4,parse-numbers=false}
\DeclareSIUnit\ohm{\ensuremath{\symup{\Omega}}} % U+2126 absent de la police
\DeclareSIUnit\division{div} % réglages d'oscilloscope (ms/div, V/div)
\DeclareSIUnit\tr{tr} % tours (vitesse de rotation, ex. \SI{1500}{\tr\per\minute})
\DeclareSIUnit\molar{M} % concentration molaire (ex. \SI{3}{\molar}, \SI{1}{\micro\molar})
\DeclareSIUnit\an{an} % année (le modèle confond parfois avec \year, un compteur TeX) — ex. \SI{5}{\kilo\meter\per\an}
\DeclareSIUnit\FCFA{FCFA} % franc CFA (ex. \SI{500}{\FCFA\per\kilo\gram})
\DeclareSIUnit\degreeBrix{\textdegree Brix} % teneur en sucre d'un sirop/jus (réfractométrie)
\DeclareSIUnit\month{mois} % le modèle utilise parfois \month, un compteur TeX, comme unité de durée
\usepackage{qrcode}
% \degree : commande fréquemment utilisée par les modèles pour un angle ou une
% température en dehors de \SI{}{} (non fournie par les paquets chargés).
\providecommand{\degree}{^{\circ}}
% \qed (fin de démonstration), utilisé par les modèles sans amsthm
\providecommand{\qed}{\hfill$\square$}
% \barre : double barre des valeurs interdites dans un tableau de variations (tkz-tab)
\providecommand{\barre}{\ensuremath{\|}}
% environnement proof (amsthm non chargé) utilisé par les modèles
\ifdefined\proof\else\newenvironment{proof}[1][Démonstration]{\par\noindent\textit{#1.}\ }{\qed\par}\fi
\usepackage{lastpage}
\usepackage{hyperref}
\hypersetup{hidelinks}

% ============================================================
% Palette « Pop » OnBuch+ — mêmes hex que la classe P (plus_ui.dart)
% ============================================================
\definecolor{poporange}{HTML}{F59321}
\definecolor{popdark}{HTML}{E07A0C}
\definecolor{popcream}{HTML}{FFF6E8}
\definecolor{popink}{HTML}{1C1714}
\definecolor{poppurple}{HTML}{7A5AE0}
\definecolor{popgreen}{HTML}{2E9E6A}
\definecolor{poppink}{HTML}{E0457B}
\definecolor{popblue}{HTML}{2F7DE1}
\definecolor{popgold}{HTML}{C98A12}
\definecolor{popmuted}{HTML}{8A7F76}
\definecolor{popline}{HTML}{EADFD2}
\definecolor{popgray}{HTML}{8A7F76}
\colorlet{popgrey}{popgray}
% Alias tolérants (le modèle nomme parfois les couleurs différemment)
\colorlet{popwhite}{white}
\colorlet{popblack}{popink}
\colorlet{popred}{poppink}
\colorlet{popyellow}{popgold}
% Teintes claires (fonds de tableaux/encadrés) — le modèle nomme parfois
% les couleurs avec un suffixe L (ex. popblueL) sans les avoir définies.
\colorlet{poporangeL}{poporange!20}
\colorlet{popdarkL}{popdark!20}
\colorlet{poppurpleL}{poppurple!20}
\colorlet{popgreenL}{popgreen!20}
\colorlet{poppinkL}{poppink!20}
\colorlet{popblueL}{popblue!20}
\colorlet{popgoldL}{popgold!20}
\colorlet{popredL}{popred!20}
\colorlet{popgrayL}{popgray!20}
% Teintes claires pour fonds de tableaux / zones
\colorlet{popblueL}{popblue!10!white}
\colorlet{poporangeL}{poporange!14!white}
\colorlet{popgreenL}{popgreen!12!white}
\colorlet{poppurpleL}{poppurple!12!white}
\colorlet{poppinkL}{poppink!12!white}
\colorlet{popgoldL}{popgold!14!white}

\pagecolor{popcream}

% ============================================================
% Typographie
% ============================================================
\defaultfontfeatures{Ligatures=TeX}
\setmainfont{PlusJakartaSans-Medium.ttf}[Path=../../../fonts/,BoldFont=PlusJakartaSans-Bold.ttf]
% Symboles, API et emojis absents de la police principale : repli sur DejaVu Sans (caractères actifs)
\newfontface\symfont{DejaVuSans.ttf}[Path=../../../fonts/]
\makeatletter
\newcommand\symactive[1]{\catcode#1=13 \begingroup\lccode`\~=#1 \lowercase{\endgroup\def~}{{\symfont\char#1}}}
\newcommand\symblank[1]{\catcode#1=13 \begingroup\lccode`\~=#1 \lowercase{\endgroup\def~}{}}
\newcommand\symhyph[1]{\catcode#1=13 \begingroup\lccode`\~=#1 \lowercase{\endgroup\def~}{\mbox{-}}}
\newcount\symc
\newcommand\symrange[2]{\symc=#1 \@whilenum\symc<#2 \do{\expandafter\symactive\expandafter{\the\symc}\advance\symc by 1 }}
\newcommand\symblankrange[2]{\symc=#1 \@whilenum\symc<#2 \do{\expandafter\symblank\expandafter{\the\symc}\advance\symc by 1 }}
\makeatother
\symrange{"0250}{"02FF}\symactive{"2713}\symactive{"2714}\symactive{"2717}\symactive{"2718}\symactive{"2610}\symactive{"2611}\symactive{"2612}
\symrange{"2600}{"27BF}
\catcode"2705=13 \begingroup\lccode`\~="2705 \lowercase{\endgroup\def~}{{\symfont\char"2713}}\symblank{"26BD}\symblank{"26A1}
\symrange{"2460}{"257F}\symactive{"223C}\symactive{"2248}\symactive{"2260}
\symhyph{"2011}\symhyph{"2E1A}\symblankrange{"1F300}{"1FAFF}\symblank{"FE0F}
% Caractères chinois : WenQuanYi Zen Hei (sous-ensemble GB2312 + pinyin), gras/italique simulés
\setCJKmainfont{WenQuanYiZenHei-subset.ttf}[Path=../../../fonts/,AutoFakeBold=3,AutoFakeSlant=0.2]
\xeCJKsetup{CJKspace=false}
% Pinyin : voyelles à caron (ǎ ǐ ǒ ǔ ǖ ǘ ǚ ǜ) absentes de la police principale -> caractères actifs
\newfontface\pyfont{WenQuanYiZenHei-subset.ttf}[Path=../../../fonts/]
\newcommand\pyactive[1]{\catcode#1=13 \begingroup\lccode`\~=#1 \lowercase{\endgroup\def~}{{\pyfont\char#1}}}
\pyactive{"01CD}\pyactive{"01CE}\pyactive{"01CF}\pyactive{"01D0}\pyactive{"01D1}\pyactive{"01D2}\pyactive{"01D3}\pyactive{"01D4}\pyactive{"01D5}\pyactive{"01D6}\pyactive{"01D7}\pyactive{"01D8}\pyactive{"01D9}\pyactive{"01DA}\pyactive{"01DB}\pyactive{"01DC}
% Maths en sans-serif (Fira Math) avec chiffres et lettres droites de Plus
% Jakarta, pour que les formules se fondent dans le texte.
\usepackage[math-style=ISO,bold-style=ISO]{unicode-math}
\setmathfont{FiraMath-Regular.otf}[Path=../../../fonts/]
\setmathfont{PlusJakartaSans-Medium.ttf}[Path=../../../fonts/,range={up/num,up/{Latin,latin},\mathup}]
\usepackage{icomma} % virgule décimale sans espace en mode math
% Caractères mathématiques Unicode absents des polices de texte (Plus Jakarta,
% Archivo Black) : rendus par leur équivalent mathématique, en texte comme en
% formule — sinon ils s'affichent en carré vide (ex. « Arithmétique dans ℤ »).
\usepackage{newunicodechar}
\newunicodechar{ℝ}{\ensuremath{\mathbb{R}}}
\newunicodechar{ℤ}{\ensuremath{\mathbb{Z}}}
\newunicodechar{ℂ}{\ensuremath{\mathbb{C}}}
\newunicodechar{ℕ}{\ensuremath{\mathbb{N}}}
\newunicodechar{ℚ}{\ensuremath{\mathbb{Q}}}
\newunicodechar{𝔻}{\ensuremath{\mathbb{D}}}
\newunicodechar{α}{\ensuremath{\alpha}}
\newunicodechar{β}{\ensuremath{\beta}}
\newunicodechar{γ}{\ensuremath{\gamma}}
\newunicodechar{δ}{\ensuremath{\delta}}
\newunicodechar{ε}{\ensuremath{\varepsilon}}
\newunicodechar{θ}{\ensuremath{\theta}}
\newunicodechar{λ}{\ensuremath{\lambda}}
\newunicodechar{μ}{\ensuremath{\mu}}
\newunicodechar{σ}{\ensuremath{\sigma}}
\newunicodechar{τ}{\ensuremath{\tau}}
\newunicodechar{φ}{\ensuremath{\varphi}}
\newunicodechar{ψ}{\ensuremath{\psi}}
\newunicodechar{ω}{\ensuremath{\omega}}
\newunicodechar{Σ}{\ensuremath{\Sigma}}
% majuscules produites par \MakeUppercase dans les titres (α-aminés -> Α)
\newunicodechar{Α}{\ensuremath{\alpha}}
\newunicodechar{Β}{\ensuremath{\beta}}
\newunicodechar{Γ}{\ensuremath{\Gamma}}
\newunicodechar{Δ}{\ensuremath{\Delta}}
\newunicodechar{Ε}{\ensuremath{\varepsilon}}
\newunicodechar{Θ}{\ensuremath{\Theta}}
\newunicodechar{Λ}{\ensuremath{\Lambda}}
\newunicodechar{Μ}{\ensuremath{\mu}}
\newunicodechar{Τ}{\ensuremath{\tau}}
\newunicodechar{Φ}{\ensuremath{\Phi}}
\newunicodechar{Ψ}{\ensuremath{\Psi}}
\newunicodechar{Ω}{\ensuremath{\Omega}}
\newunicodechar{↦}{\ensuremath{\mapsto}}
\newunicodechar{∈}{\ensuremath{\in}}
\newunicodechar{∉}{\ensuremath{\notin}}
\newunicodechar{∧}{\ensuremath{\wedge}}
\newunicodechar{⇔}{\ensuremath{\Leftrightarrow}}
\newunicodechar{⟺}{\ensuremath{\Longleftrightarrow}}
\newunicodechar{⇒}{\ensuremath{\Rightarrow}}
\newunicodechar{⟹}{\ensuremath{\Longrightarrow}}
\newunicodechar{∘}{\ensuremath{\circ}}
\newunicodechar{∪}{\ensuremath{\cup}}
\newunicodechar{∩}{\ensuremath{\cap}}
\newunicodechar{≡}{\ensuremath{\equiv}}
\newunicodechar{⊂}{\ensuremath{\subset}}
\newunicodechar{⊥}{\ensuremath{\perp}}
\newunicodechar{∃}{\ensuremath{\exists}}
\newunicodechar{∀}{\ensuremath{\forall}}
\newunicodechar{∛}{\ensuremath{\sqrt[3]{\,}}}
\newunicodechar{∜}{\ensuremath{\sqrt[4]{\,}}}
\newunicodechar{⃗}{\ensuremath{{}^{\to}}}
\newunicodechar{ⁿ}{\ifmmode^{n}\else\textsuperscript{\ensuremath{n}}\fi}
\newunicodechar{ˣ}{\ifmmode^{x}\else\textsuperscript{\ensuremath{x}}\fi}
\newunicodechar{ᵇ}{\ifmmode^{b}\else\textsuperscript{\ensuremath{b}}\fi}
\newunicodechar{ʳ}{\ifmmode^{r}\else\textsuperscript{\ensuremath{r}}\fi}
\newunicodechar{ᵘ}{\ifmmode^{u}\else\textsuperscript{\ensuremath{u}}\fi}
\newunicodechar{ʸ}{\ifmmode^{y}\else\textsuperscript{\ensuremath{y}}\fi}
\newunicodechar{ᵃ}{\ifmmode^{a}\else\textsuperscript{\ensuremath{a}}\fi}
\newunicodechar{ᵉ}{\ifmmode^{e}\else\textsuperscript{\ensuremath{e}}\fi}
\newunicodechar{ᵏ}{\ifmmode^{k}\else\textsuperscript{\ensuremath{k}}\fi}
\newunicodechar{ᵖ}{\ifmmode^{p}\else\textsuperscript{\ensuremath{p}}\fi}
\newunicodechar{ᴺ}{\ifmmode^{N}\else\textsuperscript{\ensuremath{N}}\fi}
\newunicodechar{ᶜ}{\ifmmode^{c}\else\textsuperscript{\ensuremath{c}}\fi}
\newunicodechar{ʰ}{\ifmmode^{h}\else\textsuperscript{\ensuremath{h}}\fi}
\newunicodechar{ᵠ}{\ifmmode^{q}\else\textsuperscript{\ensuremath{q}}\fi}
\newunicodechar{ₙ}{\ifmmode_{n}\else\textsubscript{\ensuremath{n}}\fi}
\newunicodechar{ₐ}{\ifmmode_{a}\else\textsubscript{\ensuremath{a}}\fi}
\newunicodechar{ᵢ}{\ifmmode_{i}\else\textsubscript{\ensuremath{i}}\fi}
\newunicodechar{ᵣ}{\ifmmode_{r}\else\textsubscript{\ensuremath{r}}\fi}
\newunicodechar{ₖ}{\ifmmode_{k}\else\textsubscript{\ensuremath{k}}\fi}
\newunicodechar{ᵦ}{\ifmmode_{\beta}\else\textsubscript{\ensuremath{\beta}}\fi}
\newunicodechar{ₓ}{\ifmmode_{x}\else\textsubscript{\ensuremath{x}}\fi}
\newfontfamily\popdisplay{ArchivoBlack-Regular.ttf}[Path=../../../fonts/]
\newfontfamily\poplogo{SpaceGrotesk-Bold.ttf}[Path=../../../fonts/]
\newfontfamily\popsemi{PlusJakartaSans-SemiBold.ttf}[Path=../../../fonts/]
\newfontfamily\popxbold{PlusJakartaSans-ExtraBold.ttf}[Path=../../../fonts/]
\newfontfamily\popxboldls{PlusJakartaSans-ExtraBold.ttf}[Path=../../../fonts/,LetterSpace=3.0]

\setlist[enumerate]{leftmargin=1.7em,itemsep=5pt,topsep=4pt}
\setlist[itemize]{leftmargin=1.7em,itemsep=4pt,topsep=4pt,label={\color{poporange}\textbullet}}
\setlength{\parindent}{0pt}
\setlength{\parskip}{5pt}
\renewcommand{\arraystretch}{1.25}

% pgfplots : style Pop par défaut
\pgfplotsset{
  every axis/.append style={axis line style={popink,line width=1pt},tick style={popink},
    grid style={popline,line width=0.5pt},label style={font=\small},tick label style={font=\footnotesize}},
  popcurve/.style={poporange,line width=1.8pt,no markers,smooth},
}
% Même style disponible dans un \draw TikZ simple (hors pgfplots)
\tikzset{popcurve/.style={poporange,line width=1.8pt}}
\tikzset{popgrid/.style={popline,very thin}}
\colorlet{popcurve}{poporange} % le nom du style est parfois utilisé comme couleur

% ============================================================
% Pill / squiggle
% ============================================================
\newcommand{\poppill}[3]{%
  \tikz[baseline=-0.55ex]\node[draw=popink,line width=1.2pt,rounded corners=2.6pt,%
    fill=#2,inner xsep=6.5pt,inner ysep=3pt]{%
    \popxboldls\fontsize{7}{7}\selectfont\color{#3}\MakeUppercase{#1}};%
}
\newcommand{\popsquiggle}[1]{%
  \tikz[baseline=-0.4ex]{
    \draw[#1,line width=2.6pt,line cap=round]
      plot [smooth] coordinates {(0,0.09) (0.34,0.01) (0.68,0.21) (1.02,0.01) (1.36,0.10)};
  }%
}
% Mot-clé surligné (marqueur orange sous le texte)
\newcommand{\cle}[1]{{\popxbold\color{popdark}#1}}

% ============================================================
% En-tête / pied
% ============================================================
\pagestyle{fancy}
\fancyhf{}
\renewcommand{\headrulewidth}{0pt}
\renewcommand{\footrulewidth}{0pt}
\lhead{\raisebox{-0.15ex}{\poplogo\fontsize{13}{13}\selectfont\color{popink}OnBuch\color{poporange}+}}
\rhead{\poppill{\DOCMATIERE\ \textbullet\ \DOCNIVEAU\ \textbullet\ Leçon \DOCLECON}{white}{popink}}
\lfoot{\popxbold\fontsize{7.6}{7.6}\selectfont\color{popmuted}\MakeUppercase{Cours signé OnBuch+}\ \textbullet\ {\popsemi\DOCTITRE}}
\rfoot{\tikz[baseline=-0.6ex]\node[rounded corners=8.5pt,fill=popink,minimum height=17pt,minimum width=17pt,inner xsep=5pt,inner sep=0]{\popxbold\fontsize{8}{8}\selectfont\color{popcream}\thepage\,/\,\pageref*{LastPage}};}

% ============================================================
% Titres
% ============================================================
\newcommand{\chaptitre}[2]{%
  \begin{tcolorbox}[enhanced,colback=poporange,colframe=popink,boxrule=2.6pt,arc=11pt,
      drop shadow=popink,left=15pt,right=15pt,top=12pt,bottom=11pt,before skip=3pt,after skip=18pt]
    \poppill{#2}{popink}{popcream}\par\vspace{9pt}
    {\raggedright\hyphenpenalty=10000\exhyphenpenalty=10000\popdisplay\fontsize{22}{24}\selectfont\color{popcream}\MakeUppercase{#1}\par}
  \end{tcolorbox}
}
% Section numérotée : pastille orange avec le numéro + titre Archivo
\newcounter{popsec}
\newcommand{\coursec}[1]{%
  \stepcounter{popsec}\setcounter{popsub}{0}%
  \par\vspace{16pt}\noindent
  \begin{minipage}{\linewidth}
  \tikz[baseline=-0.7ex]\node[draw=popink,line width=1.6pt,fill=poporange,rounded corners=4pt,minimum size=22pt,inner sep=0]{\popdisplay\fontsize{12}{12}\selectfont\color{popcream}\thepopsec};%
  \hspace{8pt}{\popdisplay\fontsize{15}{17}\selectfont\color{popink}\MakeUppercase{#1}}\par
  \vspace{4pt}\noindent{\color{poporange}\rule{36pt}{3.6pt}}
  \end{minipage}\par\nopagebreak\vspace{8pt}
}
\newcounter{popsub}
\newcommand{\courssub}[1]{%
  \stepcounter{popsub}%
  \par\vspace{9pt}\noindent{\popxbold\fontsize{11.5}{13}\selectfont\color{popink}{\color{poporange}\thepopsec.\thepopsub}\hspace{6pt}#1}\par\nopagebreak\vspace{4pt}
}

% ============================================================
% Boîtes pédagogiques — même recette : fond blanc, bordure épaisse colorée,
% ombre dure encre, badge plein. Titre optionnel [#1].
% ============================================================
\tcbset{popbase/.style={breakable,enhanced,colback=white,boxrule=2.3pt,arc=9pt,
  shadow={3pt}{-3pt}{0pt}{popink},left=14pt,right=14pt,top=11pt,bottom=11pt,before skip=11pt,after skip=15pt,
  fontupper=\popsemi\fontsize{10.6}{14.5}\selectfont\color{popink},
  fontlower=\popsemi\fontsize{10.6}{14.5}\selectfont\color{popink}}}
% Coche dessinée (le glyphe ✓ n'existe pas dans les polices du document)
\newcommand{\popcheck}[1]{\tikz[baseline=-0.5ex]\draw[#1,line width=1.6pt,line cap=round,line join=round](-0.13,0.0)--(-0.04,-0.09)--(0.13,0.1);}
\providecommand{\coche}{\popcheck{popgreen}}
\providecommand{\resultat}[1]{\par\smallskip\noindent\fcolorbox{popgreen}{popgreenL}{\parbox{\dimexpr\linewidth-2\fboxsep-2\fboxrule\relax}{\textbf{Résultat.} #1}}\par\smallskip}
\newcommand{\popboxhead}[3]{\poppill{#1}{#2}{white}\ifx\relax#3\relax\else\hspace{6pt}{\popxbold\fontsize{10.6}{12}\selectfont\color{popink}#3}\fi\par\vspace{7pt}}

\newtcolorbox{aretenir}[1][]{popbase,colframe=popgreen,before upper={\popboxhead{À retenir}{popgreen}{#1}}}
% Texte en chinois (document de lecture, dialogue, extrait) : cadre
% d'étude. Un titre optionnel (source, auteur) s'affiche dans l'en-tête.
\newtcolorbox{textebox}[1][]{popbase,colframe=popblue,before upper={\popboxhead{文本 · Texte}{popblue}{#1}},after upper={}}
% Marques ✓ / ✗ (forme correcte / fautive) souvent utilisées par les modèles
\providecommand{\cross}{\textcolor{poppink}{\ensuremath{\times}}}
\providecommand{\xmark}{\cross}\providecommand{\crossmark}{\cross}\providecommand{\cmark}{\ensuremath{\checkmark}}
% Ligne de réponse / justification à compléter (inventée par les modèles)
\providecommand{\justification}[1]{\par\noindent\textit{Justification~:} \dotfill\par}
% \textipa (package tipa non chargé) : la police ne couvre pas l'API -> simple passage
\providecommand{\textipa}[1]{#1}
% Soulignement (ulem non chargé) : \uline, \ul
\providecommand{\uline}[1]{\underline{#1}}\providecommand{\ul}[1]{\underline{#1}}
% \glossaire{\item ...} inventée par les modèles : liste de glossaire
\providecommand{\glossaire}[1]{\begin{itemize}#1\end{itemize}}
% \alert (beamer) inventée par les modèles : texte mis en évidence
\providecommand{\alert}[1]{\textbf{\textcolor{poppink}{#1}}}
% variantes ASCII d'ordinaux écrites en macro
\providecommand{\iere}{\textsuperscript{re}}\providecommand{\ieme}{\textsuperscript{e}}\providecommand{\ere}{\textsuperscript{re}}\providecommand{\eme}{\textsuperscript{e}}\providecommand{\er}{\textsuperscript{er}}\providecommand{\ie}{\textsuperscript{e}}
% Ordinaux français écrits en macro par les modèles : 1\ière, 3\ième
\newcommand{\ière}{\textsuperscript{re}}\newcommand{\ième}{\textsuperscript{e}}
% \exemple (inventée par les modèles) : étiquette « Exemple : »
\providecommand{\exemple}{\textbf{Exemple~:}\ }
% \square utilisable aussi hors mode math (cases à cocher)
\AtBeginDocument{\let\squareMath\square\renewcommand{\square}{\ensuremath{\squareMath}}}
% Mot ou phrase en chinois à l'intérieur d'un paragraphe français
\newcommand{\zh}[1]{\ifmmode\text{#1}\else{#1}\fi}
\let\eng\zh \let\esp\zh \let\all\zh % alias tolérés
% flèches d'intonation inventées par les modèles
\providecommand{\textsearrow}{}\renewcommand{\textsearrow}{\ensuremath{\searrow}}\providecommand{\textnearrow}{}\renewcommand{\textnearrow}{\ensuremath{\nearrow}}
% \fr{..} : mot français (inventé par les modèles)
\providecommand{\fr}[1]{\textit{#1}}
% zhbox : encadré de texte chinois inventé par les modèles
\newenvironment{zhbox}{\begin{quote}}{\end{quote}}
% \vs : « versus » inventé par les modèles
\providecommand{\vs}{\textit{vs}\ }
% \blank : case à compléter inventée par les modèles
\providecommand{\blank}[1][]{\underline{\hspace{1.6em}}}
\newtcolorbox{exemplebox}[1][]{popbase,colframe=poporange,before upper={\popboxhead{Exemple}{poporange}{#1}}}
\newtcolorbox{definition}[1][]{popbase,colframe=popblue,before upper={\popboxhead{Définition}{popblue}{#1}}}
\newtcolorbox{propriete}[1][]{popbase,colframe=poppurple,before upper={\popboxhead{Propriété}{poppurple}{#1}}}
\newtcolorbox{methode}[1][]{popbase,colframe=popgold,before upper={\popboxhead{Méthode}{popgold}{#1}}}
\newtcolorbox{attention}[1][]{popbase,colframe=poppink,before upper={\popboxhead{Attention}{poppink}{#1}}}
\newtcolorbox{experience}[1][]{popbase,colframe=popblue,colback=popblueL,before upper={\popboxhead{Expérience}{popblue}{#1}}}
% « Le savais-tu ? » : carte violette pleine (contexte, culture, Cameroun)
\newtcolorbox{savaistu}[1][]{popbase,colback=poppurple,colframe=popink,
  fontupper=\popsemi\fontsize{10.4}{14.2}\selectfont\color{white},
  before upper={\poppill{Le savais-tu\,?}{white}{poppurple}\ifx\relax#1\relax\else\hspace{6pt}{\popxbold\color{white}#1}\fi\par\vspace{7pt}}}
% Exercice résolu : énoncé en haut, \tcblower puis la solution sur fond crème
\newcounter{popexo}\newcounter{popexr}
\newtcolorbox{exoresolu}[1][]{popbase,colframe=poporange,colbacklower=popcream,
  segmentation style={popink,line width=1.2pt,dashed},
  before upper={\stepcounter{popexr}\popboxhead{Exercice résolu \thepopexr}{poporange}{#1}},
  before lower={\poppill{Solution}{popink}{popcream}\par\vspace{6pt}}}
% Exercice (non résolu en place) — la correction est dans la partie Corrigés
\newtcolorbox{exercice}[1][]{popbase,colframe=popink,
  before upper={\stepcounter{popexo}\popboxhead{Exercice \thepopexo}{popink}{#1}}}
% Corrigé
\newtcolorbox{corrige}[1][]{popbase,colframe=popgreen,colback=popgreenL,
  before upper={\popboxhead{Corrigé}{popgreen}{#1}}}
% Objectifs / prérequis / bilan
\newtcolorbox{objectifs}{popbase,colframe=popink,colback=poporangeL,
  before upper={\popboxhead{Objectifs de la leçon}{popink}{}}}
\newtcolorbox{prerequis}{popbase,colframe=popink,colback=popblueL,
  before upper={\popboxhead{Prérequis}{popblue}{}}}
\newtcolorbox{bilan}{popbase,colframe=popink,colback=white,boxrule=2.8pt,
  before upper={\popboxhead{Fiche bilan}{poporange}{L'essentiel à connaître pour l'examen}}}
% Figure encadrée avec légende
\newtcolorbox{popfigure}[1][]{enhanced,breakable=false,colback=white,colframe=popline,boxrule=1.4pt,arc=8pt,
  left=10pt,right=10pt,top=10pt,bottom=8pt,before skip=10pt,after skip=12pt,halign=center,
  lower separated=false,halign lower=center,
  fontlower=\popsemi\fontsize{9}{11.5}\selectfont\color{popmuted},
  #1}
\newcommand{\legende}[1]{\par\vspace{6pt}{\popsemi\fontsize{9}{11.5}\selectfont\color{popmuted}#1}}

% ============================================================
% Signature OnBuch+
% ============================================================
\newcommand{\poplogoinline}[1]{{\poplogo\fontsize{#1}{#1}\selectfont\color{popink}OnBuch\color{poporange}+}}
\newcommand{\signatureonbuch}{%
  \begin{tcolorbox}[enhanced,colback=popink,colframe=popink,boxrule=2.2pt,arc=10pt,
      drop shadow=poporange,left=14pt,right=14pt,top=10pt,bottom=10pt,before skip=0pt,after skip=0pt]
    \tikz[baseline=-0.7ex]\node[circle,fill=popgreen,draw=popcream,line width=1.3pt,minimum size=22pt,inner sep=0]{\popcheck{white}};%
    \hspace{9pt}\begin{minipage}[c]{0.62\linewidth}
      {\popxbold\fontsize{12}{13}\selectfont\color{popcream}Signé\ \textbullet\ {\poplogo OnBuch\color{poporange}+}}\par\vspace{2pt}
      {\raggedright\popsemi\fontsize{8}{10.5}\selectfont\color{popline}\MakeUppercase{Équipe pédagogique OnBuch+}\\\MakeUppercase{Conforme au programme officiel MINESEC}\par}
    \end{minipage}\hfill
    {\poplogo\fontsize{22}{22}\selectfont\color{popcream}OnBuch\color{poporange}+}
  \end{tcolorbox}%
}

% ============================================================
% Page de garde
% #1 titre · #2 matière · #3 niveau · #4 module · #5 n° leçon · #6 durée ·
% #7 description / objectifs (texte ou itemize)
% ============================================================
\newcommand{\pagegarde}[7]{%
  \thispagestyle{empty}
  \newgeometry{a4paper,margin=1.7cm,top=1.6cm,bottom=1.4cm}%
  \begin{tikzpicture}[remember picture,overlay]
    \fill[poporange!18] ([xshift=-3.2cm,yshift=-3.6cm]current page.north east) circle (4.6cm);
    \fill[poppurple!14] ([xshift=2.4cm,yshift=7.6cm]current page.south west) circle (3.8cm);
  \end{tikzpicture}%
  {\poplogo\fontsize{30}{30}\selectfont\color{popink}OnBuch\color{poporange}+}\hfill
  \raisebox{0.6ex}{\poppill{Cours signé \textbullet\ Premium}{popink}{popcream}}\par
  \vspace{1pt}\popsquiggle{poppurple}\par
  \vspace{8pt}
  \begin{tcolorbox}[enhanced,colback=poporange,colframe=popink,boxrule=2.8pt,arc=13pt,
      drop shadow=popink,left=18pt,right=18pt,top=12pt,bottom=13pt,after skip=10pt]
    \poppill{#2 \textbullet\ #3}{popink}{popcream}\hspace{5pt}\poppill{#4}{white}{popink}\par\vspace{8pt}
    {\popdisplay\fontsize{14}{14}\selectfont\color{popink}LEÇON #5}\par\vspace{4pt}
    {\raggedright\hyphenpenalty=10000\exhyphenpenalty=10000\popdisplay\fontsize{25}{27}\selectfont\color{popcream}\MakeUppercase{#1}\par}
    \vspace{8pt}
    {\popxbold\fontsize{9.5}{9.5}\selectfont\color{popink}\MakeUppercase{Durée conseillée : #6}}
  \end{tcolorbox}
  \begin{tcolorbox}[enhanced,colback=white,colframe=popink,boxrule=2.2pt,arc=10pt,
      drop shadow=popink,left=16pt,right=16pt,top=10pt,bottom=10pt,after skip=8pt]
    \poppill{Dans cette leçon}{poppurple}{white}\par\vspace{5pt}
    \popsemi\fontsize{9.3}{12.6}\selectfont\color{popink}#7
  \end{tcolorbox}
  \noindent\begin{minipage}[t]{0.487\linewidth}
    \begin{tcolorbox}[enhanced,colback=popgreen,colframe=popink,boxrule=2.2pt,arc=10pt,
        drop shadow=popink,left=11pt,right=11pt,top=10pt,bottom=10pt,height=3.1cm,valign=center]
      \tcbox[colback=white,colframe=popink,boxrule=1.3pt,arc=5pt,boxsep=0pt,nobeforeafter,
        left=3pt,right=3pt,top=3pt,bottom=3pt,valign=center]{\qrcode[height=1.9cm]{https://play.google.com/store/apps/details?id=com.luvvix.onbuch&hl=fr}}%
      \hspace{8pt}\begin{minipage}[c]{0.5\linewidth}
        {\popdisplay\fontsize{11}{12.5}\selectfont\color{white}\MakeUppercase{Télécharge OnBuch}}\par\vspace{4pt}
        {\raggedright\popsemi\fontsize{8.4}{11}\selectfont\color{white}Gratuit \textbullet\ Google Play\\Tuteur IA, annales, concours, résultats\par}
      \end{minipage}
    \end{tcolorbox}
  \end{minipage}\hfill
  \begin{minipage}[t]{0.487\linewidth}
    \begin{tcolorbox}[enhanced,colback=poppurple,colframe=popink,boxrule=2.2pt,arc=10pt,
        drop shadow=popink,left=11pt,right=11pt,top=10pt,bottom=10pt,height=3.1cm,valign=center]
      \tikz[baseline=-0.7ex]\node[circle,fill=white,draw=popink,line width=1.3pt,minimum size=1.6cm,inner sep=0]{\popdisplay\fontsize{24}{24}\selectfont\color{poppurple}+};%
      \hspace{8pt}\begin{minipage}[c]{0.6\linewidth}
        {\popdisplay\fontsize{11}{12.5}\selectfont\color{white}\MakeUppercase{Passe à OnBuch+}}\par\vspace{4pt}
        {\raggedright\popsemi\fontsize{8.4}{11}\selectfont\color{white}Tous les cours signés, Tuteur IA illimité, fiches PDF — onglet OnBuch+ de l'app\par}
      \end{minipage}
    \end{tcolorbox}
  \end{minipage}
  \vspace{8pt}
  \popcontenu
  \vfill
  \signatureonbuch
  \restoregeometry
  \newpage
}

% Rangée « Ce cours contient » (page de garde)
\newcommand{\poptile}[3]{% #1 couleur #2 gros texte #3 légende
  \begin{tcolorbox}[enhanced,colback=#1,colframe=popink,boxrule=2pt,arc=9pt,shadow={2.5pt}{-2.5pt}{0pt}{popink},
      left=6pt,right=6pt,top=9pt,bottom=9pt,halign=center,valign=center,height=1.75cm,width=\linewidth,nobeforeafter]
    {\popdisplay\fontsize{12.5}{13}\selectfont\color{white}\MakeUppercase{#2}}\par\vspace{3pt}
    {\centering\popsemi\fontsize{7.8}{9.5}\selectfont\color{white}#3\par}
  \end{tcolorbox}}
\newcommand{\popcontenu}{%
  \par\noindent{\popxboldls\fontsize{7.5}{7.5}\selectfont\color{popmuted}CE COURS SIGNÉ CONTIENT}\par\vspace{7pt}
  \noindent\begin{minipage}[t]{0.235\linewidth}\poptile{popblue}{Cours}{complet, illustré, pas à pas}\end{minipage}\hfill
  \begin{minipage}[t]{0.235\linewidth}\poptile{poporange}{Méthodes}{et exercices résolus}\end{minipage}\hfill
  \begin{minipage}[t]{0.235\linewidth}\poptile{poppink}{Exercices}{type examen + corrigés}\end{minipage}\hfill
  \begin{minipage}[t]{0.235\linewidth}\poptile{popgreen}{Fiche bilan}{l'essentiel à retenir}\end{minipage}\par}

% Sommaire
\newtcolorbox{sommaire}{enhanced,colback=white,colframe=popink,boxrule=2.2pt,arc=10pt,
  drop shadow=popink,left=18pt,right=18pt,top=13pt,bottom=13pt,before skip=6pt,after skip=20pt,
  before upper={\poppill{Sommaire}{poppurple}{white}\par\vspace{9pt}\popsemi\fontsize{10.6}{14.5}\selectfont\color{popink}}}

% Signature de fin de cours — poussée tout en bas de la dernière page
\newcommand{\findecours}{%
  \par\vfill\nopagebreak
  \begin{center}\popsquiggle{poporange}\end{center}\vspace{4pt}
  \signatureonbuch
}
\AtBeginDocument{\def\llbracket{\mathopen{[\mkern-3.5mu[}}\def\rrbracket{\mathclose{]\mkern-3.5mu]}}} % intervalles d'entiers (glyphe ⟦ absent de la police)
% Glyphes absents de Fira Math : redéfinis à partir de glyphes disponibles
\AtBeginDocument{%
  \def\setminus{\mathbin{\backslash}}\let\smallsetminus\setminus
  \def\vdots{\mathord{\vbox{\baselineskip3.2pt\lineskiplimit0pt\kern4pt\hbox{.}\hbox{.}\hbox{.}}}}%
  \def\ddots{\mathinner{\mkern1mu\raise6.5pt\hbox{.}\mkern2mu\raise3.6pt\hbox{.}\mkern2mu\raise0.7pt\hbox{.}\mkern1mu}}%
  \def\triangle{\mathord{\scalebox{0.9}{$\symup{\Delta}$}}}%
  \def\nabla{\mathord{\raisebox{\depth}{\scalebox{1}[-1]{$\symup{\Delta}$}}}}%
  \def\checkmark{\popcheck{popdark}}%
  \def\star{\mathbin{\ast}}\def\bigstar{\mathord{\ast}}%
  \def\diamond{\mathbin{\scalebox{0.6}{$\Diamond$}}}%
  \def\bigcup{\mathop{\mathchoice{\raisebox{-0.3em}{\scalebox{1.7}{$\cup$}}}{\scalebox{1.25}{$\cup$}}{\cup}{\cup}}\displaylimits}%
  \def\bigcap{\mathop{\mathchoice{\raisebox{-0.3em}{\scalebox{1.7}{$\cap$}}}{\scalebox{1.25}{$\cap$}}{\cap}{\cap}}\displaylimits}%
  \def\lesseqqgtr{\mathrel{\lessgtr}}\def\bigoplus{\mathop{\oplus}\displaylimits}%
}
\providecommand{\ssi}{si et seulement si} % abréviation fréquente
\AtBeginDocument{\providecommand{\wideparen}{\overparen}} % arc de cercle

%% FIGFIX-BEGIN v5 — correctifs globaux des figures (docs/onbuch-generation/tools/figfix)
\usepackage{adjustbox}\usepackage{etoolbox}\tcbuselibrary{raster}
% toute figure TikZ/pgfplots plus large que la ligne est réduite à la largeur disponible
\BeforeBeginEnvironment{tikzpicture}{\begin{adjustbox}{max width=\linewidth}}
\AfterEndEnvironment{tikzpicture}{\end{adjustbox}}
% légendes pgfplots lisibles par-dessus les courbes
\pgfplotsset{every axis/.append style={legend style={fill=white,fill opacity=0.92,text opacity=1,draw=black!15,inner sep=3pt}}}
% étiquettes d'arêtes (syntaxe edge["..."]) avec fond blanc
\tikzset{every edge quotes/.append style={fill=white,inner sep=1.2pt}}
%% FIGFIX-END
\begin{document}

\pagegarde{Leçon 1 — Vocabulaire et compréhension de texte : situation de la femme (égalité homme-femme)}{Chinois}{Tle A}{Module 1 : Vie familiale et intégration sociale}{ZH01}{2 h}{Cette leçon de vocabulaire et compréhension de texte porte sur la situation de la femme et l'égalité homme-femme en Chine. À travers un texte support authentique, un glossaire structuré, l'étude des radicaux et des questions de compréhension, l'élève apprend à lire, prononcer et employer le lexique du thème.
\vspace{4pt}
\begin{multicols}{2}\small
\begin{itemize}
\item À la fin de cette leçon, je sais lire à voix haute un court texte chinois sur la situation de la femme en respectant le pinyin et les tons.
\item Je sais reconnaître et traduire au moins 20 mots de vocabulaire liés à l'égalité homme-femme.
\item Je sais identifier la nature grammaticale (nom, verbe, adjectif) des mots du thème.
\item Je sais reconnaître les clés (radicaux) 女, 人, 口, 力 et expliquer leur sens.
\item Je sais écrire correctement quelques caractères du thème en respectant l'ordre des traits.
\item Je sais répondre à des questions de compréhension sur le texte, en français et en chinois simple.
\end{itemize}\end{multicols}}

\begin{sommaire}
\begin{multicols}{2}
\begin{enumerate}
\item Mise en route et découverte du thème
\item Texte support : lecture et compréhension globale
\item Glossaire du thème en tableau
\item Clés (radicaux) et ordre des traits
\item Compréhension détaillée et collocations
\item Production guidée et ouverture Cameroun-Chine
\item Activité d'intégration
\item Méthodes et exercices résolus
\item Exercices
\item Corrigés des exercices
\item Fiche bilan
\end{enumerate}
\end{multicols}
\end{sommaire}

\begin{objectifs}
\begin{itemize}
\item À la fin de cette leçon, je sais lire à voix haute un court texte chinois sur la situation de la femme en respectant le pinyin et les tons.
\item Je sais reconnaître et traduire au moins 20 mots de vocabulaire liés à l'égalité homme-femme.
\item Je sais identifier la nature grammaticale (nom, verbe, adjectif) des mots du thème.
\item Je sais reconnaître les clés (radicaux) 女, 人, 口, 力 et expliquer leur sens.
\item Je sais écrire correctement quelques caractères du thème en respectant l'ordre des traits.
\item Je sais répondre à des questions de compréhension sur le texte, en français et en chinois simple.
\item Je sais employer les collocations du thème (提高地位, 男女平等...) dans des phrases complètes.
\item Je sais comparer brièvement la situation de la femme en Chine et au Cameroun avec le vocabulaire appris.
\end{itemize}
\end{objectifs}

\begin{prerequis}
\begin{itemize}
\item Lecture du pinyin et des quatre tons (acquis de Première)
\item Vocabulaire de la famille et de la société (家, 人, 工作, 学习)
\item Structure de base de la phrase chinoise Sujet-Verbe-Complément
\item Notion de radicaux (clés) et ordre général des traits
\item Verbes courants : 是, 有, 可以, 应该
\end{itemize}
\end{prerequis}

\newpage

\coursec{Mise en route et découverte du thème}

\courssub{Situation d'accroche : la moitié du ciel}

Au Cameroun, de plus en plus de femmes sont ministres, ingénieures ou chefs d'entreprise, comme au Nigeria ou au Ghana où elles dirigent des banques et des universités. En Chine aussi, un proverbe célèbre affirme que « les femmes portent la moitié du ciel » \zh{妇女能顶半边天}. Mais l'égalité entre l'homme et la femme reste un combat quotidien, à Yaoundé comme à Pékin.

\textbf{Problématique :} comment parler de ce sujet en chinois ? Cette leçon te donne le vocabulaire et les outils pour \cle{lire}, \cle{comprendre} et \cle{discuter} un texte sur la situation de la femme en Chine.

\begin{textebox}[Proverbe célèbre]
妇女能顶半边天。\\
\emph{Fùnǚ néng dǐng bànbiāntiān.}\\
« Les femmes peuvent porter la moitié du ciel. »
\end{textebox}

Ce proverbe est traditionnellement attribué à \cle{Mao Zedong} (毛泽东, \emph{Máo Zédōng}). Il signifie que les femmes sont capables d'assumer la moitié des responsabilités de la société, au même titre que les hommes. Observons-le mot à mot :

\begin{center}
\begin{tabularx}{\linewidth}{|X|X|X|X|}
\hline
\rowcolor{popblueL} Caractères & Pinyin & Nature & Traduction \\
\hline
妇女 & fùnǚ & nom & les femmes (terme soutenu) \\
\hline
能 & néng & verbe modal & pouvoir, être capable de \\
\hline
顶 & dǐng & verbe & porter (sur la tête), soutenir \\
\hline
半边天 & bànbiāntiān & nom & la moitié du ciel \\
\hline
\end{tabularx}
\end{center}

\begin{attention}[La voyelle ü : un piège pour les francophones]
\zh{女} (nǚ, femme) se prononce avec le tréma \cle{ü} : les lèvres sont arrondies comme pour dire « ou », mais la langue est en position de « i ». Ne confonds jamais \emph{nǚ} (女, femme) et \emph{nǔ} (努, faire des efforts) : les deux points du tréma changent le sens ! Après les initiales \emph{j, q, x}, le tréma disparaît à l'écriture mais se prononce toujours : \zh{去} s'écrit \emph{qù} mais se prononce avec ü. En revanche, après \emph{n} et \emph{l}, le tréma reste obligatoire à l'écriture : \emph{nǚ} (女), \emph{lǜ} (绿, vert).
\end{attention}

\courssub{Brainstorming : quels droits les femmes ont-elles conquis ?}

Avant de lire le texte chinois, mobilisons nos connaissances en français. Au Cameroun comme dans le monde, les femmes ont conquis des droits essentiels au cours du XX\textsuperscript{e} siècle :

\begin{itemize}
\item le \cle{droit de vote} et d'éligibilité (participation politique) ;
\item le droit à l'\cle{éducation} : les filles vont à l'école au même titre que les garçons ;
\item le droit au \cle{travail} et à un salaire égal ;
\item l'accès aux \cle{fonctions publiques} : ministres, magistrates, médecins, enseignantes ;
\item le droit de choisir librement son \cle{mariage}.
\end{itemize}

\begin{experience}[Discussion de classe — 10 minutes]
En groupes de quatre, répondez en français aux questions suivantes, puis un rapporteur présente vos idées à la classe :
\begin{enumerate}
\item Citez trois métiers exercés par des femmes dans votre ville (Yaoundé, Douala, Buea…).
\item Dans votre famille, qui s'occupe des tâches ménagères ? Est-ce équitable ?
\item Connaissez-vous une femme camerounaise célèbre (politique, sport, culture) ?
\end{enumerate}
Gardez vos réponses : vous les réutiliserez en chinois à la fin de la leçon.
\end{experience}

\begin{savaistu}[La loi chinoise sur le mariage de 1950]
En Chine, la loi sur le mariage promulguée en 1950 a interdit la polygamie, les mariages forcés et arrangés, et a reconnu le divorce par consentement mutuel. Ce fut une véritable révolution pour les femmes chinoises, qui pouvaient désormais choisir librement leur époux. Au Cameroun, les droits des femmes ont aussi progressé depuis l'indépendance, notamment dans l'éducation et la vie politique.
\end{savaistu}

\courssub{Observation des caractères-clés du titre}

Le thème de la leçon s'écrit \zh{男女平等} (\emph{nánnǚ píngděng}, « égalité homme-femme »). Observons ces caractères avant de les apprendre :

\begin{center}
\begin{tabularx}{\linewidth}{|X|X|X|X|}
\hline
\rowcolor{popblueL} Caractères & Pinyin & Nature & Traduction \\
\hline
女 & nǚ & nom / clé & femme \\
\hline
男 & nán & nom & homme (sexe masculin) \\
\hline
平等 & píngděng & nom / adjectif & égalité ; égal \\
\hline
男女平等 & nánnǚ píngděng & expression & égalité homme-femme \\
\hline
\end{tabularx}
\end{center}

\begin{definition}[Le caractère 女, clé de la famille lexicale]
\zh{女} (nǚ, femme) est à la fois un caractère et une \cle{clé} (radical). On la retrouve dans de nombreux caractères liés aux femmes ou à la famille : \zh{妈} (mā, maman), \zh{姐} (jiě, grande sœur), \zh{妹} (mèi, petite sœur), \zh{好} (hǎo, bon — une femme 女 avec un enfant 子). Quand tu vois la clé 女 à gauche d'un caractère, pense « féminin ». Remarque : placée à gauche, la clé 女 s'écrit plus étroite et son troisième trait (l'horizontal) ne dépasse pas vers la droite, comme dans \zh{妈} et \zh{她}.
\end{definition}

\begin{definition}[Le caractère 男 : la force dans les champs]
\zh{男} (nán, homme) est composé de \zh{田} (tián, champ) en haut et de \zh{力} (lì, force) en bas : « celui qui met sa force dans les champs ». Cette étymologie reflète la société agricole traditionnelle chinoise. Le contraire de 男 est 女 : 男学生 (nán xuésheng, élève garçon) / 女学生 (nǚ xuésheng, élève fille).
\end{definition}

\begin{propriete}[Former des mots avec 平 et 等]
\zh{平} (píng) signifie « plat, égal, pacifique » : 和平 (hépíng, paix), 平安 (píng'ān, en sécurité). \zh{等} (děng) signifie « rang, égal ; attendre » : 等于 (děngyú, être égal à). Ensemble : 平等 (píngděng, égalité). En chinois, les mots abstraits sont souvent formés de \cle{deux caractères synonymes} accolés, comme 平等 ou 权利 (quánlì, droit).
\end{propriete}

\begin{exemplebox}[Premiers exemples avec le vocabulaire du titre]
\zh{男女平等很重要。} \emph{Nánnǚ píngděng hěn zhòngyào.} « L'égalité homme-femme est très importante. »\\
\zh{她是女学生。} \emph{Tā shì nǚ xuésheng.} « Elle est élève (fille). »\\
\zh{女人和男人都有工作的权利。} \emph{Nǚrén hé nánrén dōu yǒu gōngzuò de quánlì.} « Les femmes et les hommes ont tous le droit de travailler. »
\end{exemplebox}

\courssub{Carte mentale du thème}

Voici le réseau lexical que nous allons construire dans cette leçon. Chaque branche sera développée dans le glossaire de la section 3.

\begin{popfigure}
\begin{tikzpicture}[mindmap, grow cyclic, every node/.style={concept, align=center}, root concept/.append style={concept color=popblue, font=\bfseries}, level 1/.append style={level distance=3.2cm, sibling angle=72}, text=white]
\node[root concept, align=center]{男女平等\\ égalité\\ homme-femme}
  child[concept color=poporange]{ node[align=center]{家庭\\ jiātíng\\ famille} }
  child[concept color=popgreen]{ node[align=center]{工作\\ gōngzuò\\ travail} }
  child[concept color=poppurple]{ node[align=center]{教育\\ jiàoyù\\ éducation} }
  child[concept color=poppink]{ node[align=center]{权利\\ quánlì\\ droits} }
  child[concept color=popgold]{ node[align=center]{社会\\ shèhuì\\ société} };
\end{tikzpicture}
\legende{Carte mentale : les cinq champs lexicaux du thème « égalité homme-femme »}
\end{popfigure}

\begin{methode}[Aborder un nouveau thème lexical en chinois]
\begin{enumerate}
\item \textbf{Activer} ses connaissances en français (brainstorming) : de quoi va parler le texte ?
\item \textbf{Observer} les caractères-clés du titre : repérer les clés (女, 力…) et deviner le sens.
\item \textbf{Écouter} une première fois sans support écrit : capter le thème global, pas les détails.
\item \textbf{Relever} les mots reconnus (noms de personnes, mots déjà connus comme 工作, 学习).
\item \textbf{Formuler} une hypothèse : « Je pense que ce texte parle de… »
\end{enumerate}
\end{methode}

\courssub{Première écoute du texte, sans support écrit}

L'enseignant lit le texte support à voix haute, deux fois, à vitesse normale. Tu n'as \textbf{pas} le texte sous les yeux : ton seul outil est l'écoute. Note en français ou en pinyin ce que tu reconnais.

\begin{experience}[Écoute globale — consignes]
\begin{enumerate}
\item Première écoute : identifie le \textbf{thème général} (de quoi parle-t-on ?).
\item Deuxième écoute : relève au moins \textbf{trois mots} que tu reconnais (par exemple 女人, 工作, 孩子).
\item Compare tes notes avec ton voisin, puis formulez ensemble une phrase d'hypothèse : « Ce texte parle probablement de… »
\end{enumerate}
\end{experience}

\begin{exoresolu}[Repérage à l'écoute]
Énoncé. Pendant la première écoute, un élève a noté les mots : \zh{女人} (nǚrén), \zh{工作} (gōngzuò), \zh{家庭} (jiātíng), \zh{平等} (píngděng). Quelle hypothèse de sens peut-il formuler ?
\tcblower
Solution détaillée. Méthode : on classe les mots reconnus par champ lexical. \zh{女人} et \zh{家庭} renvoient à la sphère familiale ; \zh{工作} renvoie au travail ; \zh{平等} est le mot abstrait central. En les combinant, on obtient l'hypothèse : « Le texte parle probablement de la place des femmes dans la famille et au travail, et de l'égalité entre hommes et femmes. » C'est exactement le thème du texte support que nous lirons en section 2. À l'écoute, il suffit souvent de 3 ou 4 mots-clés pour deviner le thème global d'un texte : ne cherche jamais à tout comprendre mot par mot.
\end{exoresolu}

\begin{attention}[Écouter le chinois : les erreurs des francophones]
\begin{itemize}
\item Ne cherche pas à traduire mentalement chaque mot : tu perds le fil. Vise le \cle{sens global}.
\item Attention aux tons : \zh{买} (mǎi, acheter) et \zh{卖} (mài, vendre) ne diffèrent que par le ton. À l'écoute, le contexte t'aide.
\item Ne confonds pas \zh{他} (tā, il) et \zh{她} (tā, elle) : ils se prononcent \textbf{exactement pareil} ! Seul le contexte (ou l'écrit) permet de les distinguer. Observe d'ailleurs la clé : 她 contient la clé 女 (femme), 他 contient la clé 亻 (homme, variante de 人).
\end{itemize}
\end{attention}

\courssub{Annonce des objectifs de la leçon}

À l'issue de cette leçon, tu seras capable de :

\begin{itemize}
\item \textbf{lire et comprendre} un court texte chinois original sur la situation de la femme en Chine (section 2) ;
\item \textbf{reconnaître, prononcer et employer} une vingtaine de mots du thème : 妇女, 平等, 权利, 教育, 工作, 社会… (section 3) ;
\item \textbf{écrire} quelques caractères-clés en respectant les clés et l'ordre des traits (section 4) ;
\item \textbf{répondre} à des questions de compréhension en phrases complètes (section 5) ;
\item \textbf{comparer} la situation des femmes au Cameroun et en Chine (section 6).
\end{itemize}

\begin{aretenir}[L'essentiel de la mise en route]
\begin{itemize}
\item Le thème : \zh{男女平等} (nánnǚ píngděng), l'égalité homme-femme.
\item Le proverbe : \zh{妇女能顶半边天} — « Les femmes portent la moitié du ciel », attribué à Mao Zedong.
\item La clé \zh{女} (nǚ) signale les caractères liés aux femmes ; attention au tréma ü.
\item Stratégie d'écoute : capter le sens global grâce à quelques mots-clés, sans tout traduire.
\end{itemize}
\end{aretenir}

\begin{exercice}[À faire avant la section suivante]
\begin{enumerate}
\item Écris cinq fois le caractère \zh{女} en respectant l'ordre des traits (trois traits : le premier est un trait brisé qui descend, le deuxième une courbe vers la gauche, le troisième l'horizontal).
\item Classe ces mots dans les cinq branches de la carte mentale : \zh{学校} (xuéxiào, école), \zh{妈妈} (māma, maman), \zh{投票} (tóupiào, voter), \zh{公司} (gōngsī, entreprise), \zh{医生} (yīshēng, médecin).
\item En français : rédige trois phrases sur les droits des femmes au Cameroun ; tu les traduiras en chinois en fin de leçon.
\end{enumerate}
\end{exercice}

\begin{corrige}[Exercice « À faire avant la section suivante »]
1. \zh{女} s'écrit en trois traits : trait 1 (撇点, trait brisé : on part en haut à gauche, on descend en oblique puis on repart vers la droite en point), trait 2 (撇, courbe qui descend vers la gauche), trait 3 (横, horizontal de gauche à droite, qui traverse le caractère). Règle générale respectée : de haut en bas, et l'horizontal en dernier.\\
2. Éducation : 学校 ; famille : 妈妈 ; droits : 投票 ; travail : 公司, 医生 (un médecin exerce aussi dans la société).\\
3. Exemple de réponse attendue : « Les femmes camerounaises ont le droit de vote. Les filles vont à l'école comme les garçons. De plus en plus de femmes sont ministres ou chefs d'entreprise. »
\end{corrige}

\coursec{Texte support : lecture et compréhension globale}

\courssub{Transition : de la découverte à la lecture}

Dans la section précédente, nous avons découvert le thème de la leçon à partir d'images et de mots-clés : la \cle{femme} \zh{女人} (\emph{nǚrén}), l'\cle{égalité} \zh{平等} (\emph{píngděng}), le \cle{statut} ou la \cle{position sociale} \zh{地位} (\emph{dìwèi}). Vous avez émis des hypothèses sur le contenu du texte. Il est temps maintenant de vérifier ces hypothèses par une lecture réelle. Nous allons lire un court texte original intitulé \zh{中国妇女的地位} (\emph{Zhōngguó fùnǚ de dìwèi}), « La situation des femmes chinoises », puis le comprendre globalement, phrase par phrase.

\courssub{Le texte support : lecture}

\begin{textebox}[中国妇女的地位 — La situation des femmes chinoises]
现在，中国妇女的地位越来越高。\\
\emph{Xiànzài, Zhōngguó fùnǚ de dìwèi yuè lái yuè gāo.}\\
Aujourd'hui, le statut des femmes chinoises est de plus en plus élevé.\\[4pt]
男女平等是中国的基本国策。\\
\emph{Nánnǚ píngděng shì Zhōngguó de jīběn guócè.}\\
L'égalité homme-femme est une politique d'État fondamentale de la Chine.\\[4pt]
女人和男人一样，可以上学、工作、参加社会活动。\\
\emph{Nǚrén hé nánrén yíyàng, kěyǐ shàngxué, gōngzuò, cānjiā shèhuì huódòng.}\\
Les femmes, comme les hommes, peuvent aller à l'école, travailler et participer aux activités sociales.\\[4pt]
很多妇女当医生、老师、工程师，也有女经理和女领导。\\
\emph{Hěn duō fùnǚ dāng yīshēng, lǎoshī, gōngchéngshī, yě yǒu nǚ jīnglǐ hé nǚ lǐngdǎo.}\\
Beaucoup de femmes sont médecins, enseignantes, ingénieures ; il y a aussi des femmes directrices et des femmes dirigeantes.\\[4pt]
在家里，丈夫和妻子一起做家务、照顾孩子。\\
\emph{Zài jiā lǐ, zhàngfu hé qīzi yìqǐ zuò jiāwù, zhàogù háizi.}\\
À la maison, le mari et la femme font ensemble le ménage et s'occupent des enfants.\\[4pt]
但是，在一些农村，重男轻女的思想还存在。\\
\emph{Dànshì, zài yìxiē nóngcūn, zhòng nán qīng nǚ de sīxiǎng hái cúnzài.}\\
Mais, dans certaines campagnes, la mentalité qui valorise les garçons et méprise les filles existe encore.\\[4pt]
我们应该继续努力，让男女真正平等。\\
\emph{Wǒmen yīnggāi jìxù nǔlì, ràng nánnǚ zhēnzhèng píngděng.}\\
Nous devons continuer à faire des efforts, pour que les hommes et les femmes soient vraiment égaux.
\end{textebox}

\courssub{Lecture à voix haute : chorale puis individuelle}

La lecture se fait en trois temps :

\begin{enumerate}
\item \textbf{Lecture modèle.} Le professeur lit le texte une première fois, lentement, en marquant bien les tons. Les élèves écoutent, doigt sur le texte.
\item \textbf{Lecture chorale.} Toute la classe lit ensemble, phrase par phrase, en imitant le modèle. Insistez sur deux points de prononciation :
\begin{itemize}
\item le \textbf{troisième ton de \zh{女}} (\emph{nǚ}) : la voyelle \emph{ü} (lèvres arrondies comme pour dire « u » français, langue en position de « i ») ; le ton descend puis remonte : \emph{nǚ}. Attention à ne pas dire \emph{nǔ} ;
\item les tons de \zh{平等} (\emph{píngděng}) : deuxième ton montant sur \emph{píng}, troisième ton sur \emph{děng}. Devant un troisième ton, le premier troisième ton se prononce presque comme un deuxième ton : on entend \emph{píng-déng}.
\end{itemize}
\item \textbf{Lecture individuelle.} Quelques élèves lisent à tour de rôle ; la classe corrige les tons. Critère de réussite : chaque syllabe porte son ton exact.
\end{enumerate}

\begin{attention}[Pièges de prononciation pour les francophones]
\begin{itemize}
\item \zh{女} \emph{nǚ} : le son \emph{ü} n'existe pas en français écrit, mais c'est exactement le « u » de « tu ». Ne lisez jamais « nou » ni « niou ».
\item \zh{国} \emph{guó} : deuxième ton montant, pas de « oua » français traînant.
\item \zh{重男轻女} \emph{zhòng nán qīng nǚ} : ici \zh{重} se lit \emph{zhòng} (quatrième ton, « accorder de l'importance »), et non \emph{chóng} (« répéter »). Même caractère, deux lectures selon le sens.
\end{itemize}
\end{attention}

\courssub{Compréhension globale : traduction phrase par phrase}

\begin{methode}[Comprendre un texte chinois en quatre étapes]
\begin{enumerate}
\item \textbf{Lire le titre} : \zh{中国妇女的地位} annonce le sujet (le statut des femmes chinoises) ; on active ce qu'on sait déjà.
\item \textbf{Repérer les mots connus} : souligner les mots déjà appris (\zh{中国}, \zh{女人}, \zh{可以}, \zh{工作}, \zh{家}, \zh{孩子}...).
\item \textbf{Deviner les mots inconnus par le contexte} : par exemple, dans \zh{男女平等是中国的基本国策}, si \zh{国策} est inconnu, le mot \zh{国家} (pays) et le contexte « politique » permettent de deviner « politique d'État ».
\item \textbf{Vérifier au glossaire} : confirmer ou corriger ses hypothèses avec le tableau de vocabulaire (section suivante).
\end{enumerate}
\end{methode}

Appliquons la méthode. Voici la traduction \textbf{littérale} (mot à mot) puis \textbf{naturelle} (bon français) de chaque phrase.

\begin{center}
\begin{tabularx}{\linewidth}{|X|X|X|}
\hline
\rowcolor{popblueL} Phrase chinoise & Traduction littérale & Traduction naturelle \\
\hline
中国妇女的地位越来越高。 & femmes-chinoises-de-statut-de-plus-en-plus-élevé & Le statut des femmes chinoises est de plus en plus élevé. \\
\hline
男女平等是中国的基本国策。 & hommes-femmes-égalité-est-Chine-de-fondamentale-politique-d'État & L'égalité homme-femme est une politique d'État fondamentale de la Chine. \\
\hline
女人和男人一样，可以上学、工作、参加社会活动。 & femmes-et-hommes-pareils-peuvent-étudier-travailler-participer-activités-sociales & Les femmes, comme les hommes, peuvent étudier, travailler et participer à la vie sociale. \\
\hline
很多妇女当医生、老师、工程师，也有女经理和女领导。 & beaucoup-femmes-font-médecin-professeur-ingénieur-aussi-il-y-a-femmes-directrices-et-femmes-dirigeantes & Beaucoup de femmes exercent comme médecins, enseignantes ou ingénieures ; il y a aussi des femmes cadres et dirigeantes. \\
\hline
在家里，丈夫和妻子一起做家务、照顾孩子。 & à-maison-mari-et-épouse-ensemble-font-ménage-s'occupent-enfants & À la maison, mari et femme partagent les tâches ménagères et s'occupent ensemble des enfants. \\
\hline
但是，在一些农村，重男轻女的思想还存在。 & mais-dans-certaines-campagnes-valoriser-garçons-mépriser-filles-de-mentalité-encore-existe & Mais dans certaines campagnes, la préférence pour les garçons existe encore. \\
\hline
我们应该继续努力，让男女真正平等。 & nous-devons-continuer-efforts-faire-que-hommes-femmes-vraiment-égaux & Nous devons poursuivre nos efforts pour une égalité réelle entre hommes et femmes. \\
\hline
\end{tabularx}
\end{center}

\courssub{Observation grammaticale : deux structures-clés du texte}

\begin{exemplebox}[越来越 + adjectif : « de plus en plus... »]
Observez : \zh{地位越来越高} (\emph{dìwèi yuè lái yuè gāo}) = « le statut est de plus en plus élevé ».\\
Structure : \textbf{越来越 + adjectif}. Autres exemples :
\begin{itemize}
\item \zh{天气越来越热。} \emph{Tiānqì yuè lái yuè rè.} « Il fait de plus en plus chaud. »
\item \zh{我的汉语越来越好。} \emph{Wǒ de Hànyǔ yuè lái yuè hǎo.} « Mon chinois est de mieux en mieux. »
\item \zh{雅温得的汽车越来越多。} \emph{Yǎwēndé de qìchē yuè lái yuè duō.} « Il y a de plus en plus de voitures à Yaoundé. »
\end{itemize}
Remarque : après \zh{越来越}, on ne met jamais \zh{很}. On dit \zh{越来越好}, jamais \zh{越来越很好}.
\end{exemplebox}

\begin{attention}[La place de 一样 : une erreur classique des francophones]
En français, on dit « les femmes sont \emph{pareilles aux} hommes » : le mot « pareil » vient d'abord. En chinois, \zh{一样} (\emph{yíyàng}, « pareil ») se place \textbf{après} les deux éléments comparés, reliés par \zh{和} :\\
\zh{女人和男人一样。} \emph{Nǚrén hé nánrén yíyàng.} « Les femmes et les hommes sont pareils. »\\
Structure : \textbf{A + 和 + B + 一样}. Ne dites jamais \zh{一样女人和男人}. Pour ajouter un adjectif : \zh{女人和男人一样重要。} « Les femmes sont aussi importantes que les hommes. »
\end{attention}

\courssub{Repérage des mots connus et hypothèses de sens}

\begin{exoresolu}[Deviner un mot inconnu par le contexte]
Énoncé : dans la phrase \zh{男女平等是中国的基本国策}, vous ne connaissez pas le mot \zh{国策}. En vous aidant du contexte et des caractères, proposez un sens, puis vérifiez.
\tcblower
Solution détaillée :
\begin{enumerate}
\item Le caractère \zh{国} est connu : on le retrouve dans \zh{中国} (Zhōngguó, la Chine) et \zh{国家} (guójiā, le pays, l'État). Il signifie « pays, État ».
\item Le caractère \zh{策} est inconnu, mais la phrase dit que l'égalité « est » (\zh{是}) quelque chose de « fondamental » (\zh{基本}) lié à l'État.
\item Hypothèse : « une règle, une politique de l'État ».
\item Vérification au glossaire : \zh{国策} (\emph{guócè}) = « politique d'État ». Hypothèse confirmée. L'expression entière \zh{基本国策} signifie « politique d'État fondamentale ».
\end{enumerate}
\end{exoresolu}

\begin{center}
\begin{tabularx}{\linewidth}{|X|X|X|X|}
\hline
\rowcolor{popblueL} Mot & Pinyin & Indice dans le contexte & Sens vérifié \\
\hline
地位 & dìwèi & « de plus en plus élevé » & statut, position sociale \\
\hline
国策 & guócè & lié à 国 (État) & politique d'État \\
\hline
参加 & cānjiā & devant « activités sociales » & participer à \\
\hline
照顾 & zhàogù & devant « enfants » & s'occuper de \\
\hline
重男轻女 & zhòng nán qīng nǚ & « mais... existe encore » (idée négative) & valoriser les garçons, mépriser les filles \\
\hline
\end{tabularx}
\end{center}

\courssub{Repères historiques : frise chronologique}

\begin{popfigure}
\begin{tikzpicture}[scale=1]
\draw[line width=2pt, popblue, -{Stealth[length=3mm]}] (0,0) -- (12.5,0);
\foreach \x/\date/\txt/\col in {0.8/1950/{婚姻法\\ Loi sur le mariage :\\ mariage libre,\\ monogamie}/poporange, 6.2/今天 Aujourd'hui/{女医生、女工程师、\\ 女领导\\ Femmes médecins,\\ ingénieures, dirigeantes}/popgreen, 11.4/目标 Objectif/{真正平等\\ L'égalité réelle\\ entre hommes\\ et femmes}/poppurple}{
  \fill[\col] (\x,0) circle (0.12);
  \node[above=0.15cm, font=\bfseries\small, text=\col] at (\x,0) {\date};
  \node[below=0.2cm, align=center, font=\footnotesize] at (\x,0) {\txt};
}
\end{tikzpicture}
\legende{Frise historique : les grandes étapes de l'évolution du statut des femmes en Chine, de la loi sur le mariage de 1950 à l'objectif d'égalité réelle.}
\end{popfigure}

\courssub{Questions de compréhension globale}

\begin{exercice}[Premières questions sur le texte]
Répondez en français, en vous appuyant sur le texte :
\begin{enumerate}
\item Quel est le sujet général du texte ?
\item Selon le texte, quelles activités les femmes chinoises peuvent-elles faire aujourd'hui ?
\item Quel problème existe encore dans certaines campagnes ?
\item Que propose le texte pour l'avenir ?
\end{enumerate}
\end{exercice}

\begin{corrige}[Exercice : questions de compréhension globale]
\begin{enumerate}
\item Le texte parle de la situation (le statut) des femmes en Chine et de l'égalité entre hommes et femmes.
\item Elles peuvent aller à l'école, travailler, participer aux activités sociales ; beaucoup sont médecins, enseignantes, ingénieures, directrices ou dirigeantes.
\item Dans certaines campagnes, la mentalité \zh{重男轻女} (préférer les garçons et mépriser les filles) existe encore.
\item Le texte propose de continuer les efforts (\zh{继续努力}) pour parvenir à une égalité réelle (\zh{真正平等}) entre hommes et femmes.
\end{enumerate}
\end{corrige}

\begin{aretenir}[L'essentiel de cette section]
\begin{itemize}
\item Lire un texte chinois se fait en quatre étapes : titre, mots connus, hypothèses par le contexte, vérification au glossaire.
\item Structure \textbf{越来越 + adjectif} = « de plus en plus... » (sans \zh{很}).
\item Structure \textbf{A + 和 + B + 一样} = « A et B sont pareils » ; \zh{一样} se place après les deux éléments.
\item Idée directrice du texte : le statut des femmes chinoises s'élève, mais l'égalité réelle demande encore des efforts — un débat qui concerne aussi le Cameroun.
\end{itemize}
\end{aretenir}

\coursec{Glossaire du thème en tableau}

Après avoir lu et compris globalement le texte support sur la situation de la femme et l'égalité homme-femme, nous allons maintenant nous arrêter sur chaque mot important. Un texte ne se possède vraiment que si l'on maîtrise son vocabulaire : reconnaître les caractères, les prononcer avec les bons tons, connaître leur nature grammaticale et savoir les employer dans une phrase. C'est l'objectif de cette section. Nous procéderons en quatre étapes : observer les mots dans leur contexte, les classer dans un tableau, travailler la prononciation, puis étudier les emplois des mots clés.

\courssub{Observer avant de mémoriser : les mots dans leur contexte}

Avant de présenter le tableau complet, observons quelques phrases tirées du texte support. Soulignez mentalement les mots nouveaux et essayez de deviner leur sens grâce au contexte :

\begin{exemplebox}[Les mots du thème en contexte]
\zh{现在，妇女的社会地位提高了。} \emph{Xiànzài, fùnǚ de shèhuì dìwèi tígāo le.} « Aujourd'hui, le statut social des femmes s'est élevé. »\\
\zh{男人和妻子一起做家务。} \emph{Nánrén hé qīzi yìqǐ zuò jiāwù.} « Les hommes font le ménage ensemble avec leur épouse. »\\
\zh{但是，重男轻女的思想还存在。} \emph{Dànshì, zhòngnánqīngnǚ de sīxiǎng hái cúnzài.} « Mais la mentalité qui valorise les garçons et méprise les filles existe encore. »\\
\zh{她努力工作，当了医生。} \emph{Tā nǔlì gōngzuò, dāng le yīshēng.} « Elle a travaillé dur et est devenue médecin. »
\end{exemplebox}

On remarque déjà une organisation naturelle : des \cle{noms} (妇女, 地位, 社会), des \cle{verbes} (参加, 当, 照顾, 提高), des \cle{adjectifs et expressions} (平等, 一样, 重男轻女) et des \cle{mots de liaison} (但是, 也, 还, 一起). C'est exactement la structure de notre glossaire. Classer les mots par nature grammaticale est une excellente stratégie de mémorisation : le cerveau retient mieux des familles que des listes désordonnées.

\courssub{Tableau complet du glossaire}

\begin{center}
\begin{tabularx}{\linewidth}{|X|X|X|X|}
\hline
\rowcolor{popblueL} Caractères & Pinyin & Nature & Traduction \\
\hline
妇女 & fùnǚ & nom & femme (terme général, soutenu) \\
\hline
男人 & nánrén & nom & homme \\
\hline
女人 & nǚrén & nom & femme (terme courant) \\
\hline
地位 & dìwèi & nom & statut, position \\
\hline
社会 & shèhuì & nom & société \\
\hline
家庭 & jiātíng & nom & famille \\
\hline
权利 & quánlì & nom & droit \\
\hline
国策 & guócè & nom & politique d'État \\
\hline
思想 & sīxiǎng & nom & mentalité, idée, pensée \\
\hline
家务 & jiāwù & nom & tâches ménagères \\
\hline
丈夫 & zhàngfu & nom & mari (ton neutre sur 夫) \\
\hline
妻子 & qīzi & nom & épouse (ton neutre sur 子) \\
\hline
孩子 & háizi & nom & enfant (ton neutre sur 子) \\
\hline
劳动 & láodòng & nom / verbe & travail ; travailler \\
\hline
\rowcolor{popgreenL} 参加 & cānjiā & verbe & participer à \\
\hline
\rowcolor{popgreenL} 当 & dāng & verbe & exercer le métier de, devenir \\
\hline
\rowcolor{popgreenL} 照顾 & zhàogù & verbe & s'occuper de \\
\hline
\rowcolor{popgreenL} 存在 & cúnzài & verbe & exister \\
\hline
\rowcolor{popgreenL} 努力 & nǔlì & verbe / adjectif & faire des efforts ; appliqué \\
\hline
\rowcolor{popgreenL} 提高 & tígāo & verbe & élever, améliorer \\
\hline
\rowcolor{poporangeL} 平等 & píngděng & adjectif / nom & égal ; égalité \\
\hline
\rowcolor{poporangeL} 一样 & yíyàng & adjectif & pareil, identique \\
\hline
\rowcolor{poporangeL} 真正 & zhēnzhèng & adjectif / adverbe & véritable ; vraiment \\
\hline
\rowcolor{poporangeL} 重男轻女 & zhòngnánqīngnǚ & expression figée & valoriser les garçons et mépriser les filles \\
\hline
但是 & dànshì & conjonction & mais \\
\hline
也 & yě & adverbe & aussi \\
\hline
还 & hái & adverbe & encore, de plus \\
\hline
一起 & yìqǐ & adverbe & ensemble \\
\hline
\end{tabularx}
\end{center}

\begin{popfigure}
\begin{tikzpicture}[
  bloc/.style={rectangle, rounded corners=3pt, align=center, font=\small, inner sep=4pt},
  racine/.style={rectangle, rounded corners=4pt, fill=popink, text=white, align=center, font=\bfseries\small, inner sep=5pt}
]
\node[racine, align=center] (r) at (0,0){男女平等\\ Nánnǚ píngděng};
\node[bloc, fill=popblueL, align=center] (n) at (-4.6,2.4){\textbf{Noms}\\ 妇女 fùnǚ\\ 地位 dìwèi\\ 社会 shèhuì\\ 家庭 jiātíng\\ 权利 quánlì\\ 国策 guócè\\ 思想 sīxiǎng\\ 家务 jiāwù\\ 丈夫 zhàngfu\\ 妻子 qīzi\\ 孩子 háizi\\ 劳动 láodòng};
\node[bloc, fill=popgreenL, align=center] (v) at (0,3.2){\textbf{Verbes}\\ 参加 cānjiā\\ 当 dāng\\ 照顾 zhàogù\\ 存在 cúnzài\\ 努力 nǔlì\\ 提高 tígāo};
\node[bloc, fill=poporangeL, align=center] (a) at (4.6,2.4){\textbf{Adjectifs et expressions}\\ 平等 píngděng\\ 一样 yíyàng\\ 真正 zhēnzhèng\\ 重男轻女\\ zhòngnánqīngnǚ\\ \textbf{Liaison :} 但是 dànshì\\ 也 yě \quad 还 hái \quad 一起 yìqǐ};
\draw[-{Stealth}, thick, popblue] (r) -- (n);
\draw[-{Stealth}, thick, popgreen] (r) -- (v);
\draw[-{Stealth}, thick, poporange] (r) -- (a);
\end{tikzpicture}
\legende{Carte mentale du glossaire : trois blocs colorés — noms (bleu), verbes (vert), adjectifs et expressions (orange).}
\end{popfigure}

\courssub{Prononciation guidée : les tons mot par mot}

En chinois, un mot ne se mémorise jamais sans ses tons. Voici les points délicats de ce glossaire, mot par mot.

\begin{itemize}
\item \zh{妇女} \emph{fùnǚ} : 4e ton + 3e ton. Le 3e ton est bas ; devant une pause ou en fin de mot, il descend puis remonte légèrement. Ne le transformez pas en 2e ton.
\item \zh{男人} \emph{nánrén} : 2e + 2e ton, deux montées successives ; gardez le \emph{-r-} rétroflexe de 儿 fondu dans \emph{nánr}.
\item \zh{地位} \emph{dìwèi} : 4e + 4e ton, deux tons descendants secs et nets.
\item \zh{社会} \emph{shèhuì} : 4e + 4e ton ; attention au \emph{sh} rétroflexe (langue recourbée).
\item \zh{权利} \emph{quánlì} : 2e + 4e ton ; le \emph{qu-} se prononce avec les lèvres arrondies (le \emph{u} après \emph{q} se lit comme \emph{ü}).
\item \zh{妻子} \emph{qīzi} : 1er ton haut et plat, puis \textbf{ton neutre} sur 子, court et léger.
\item \zh{孩子} \emph{háizi} : 2e ton montant, puis \textbf{ton neutre} sur 子.
\item \zh{丈夫} \emph{zhàngfu} : 4e ton, puis ton neutre sur 夫.
\item \zh{努力} \emph{nǔlì} : 3e + 4e ton. Le 3e ton est prononcé bas (demi-3e ton) devant le 4e ton : ne le faites pas monter.
\item \zh{存在} \emph{cúnzài} : 2e + 4e ton ; \emph{c-} est une affriquée aspirée (comme « ts » fortement soufflé).
\item \zh{重男轻女} \emph{zhòngnánqīngnǚ} : 4e, 2e, 1er, 3e ton. Ici 重 se prononce \emph{zhòng} (« lourd, important ») et non \emph{chóng} (« répéter »).
\item \zh{一样} \emph{yíyàng} : par la règle de changement de ton de 一, 一 (yī) devient \emph{yí} devant un 4e ton.
\item \zh{一起} \emph{yìqǐ} : devant un 3e ton, 一 devient \emph{yì} (4e ton).
\end{itemize}

\begin{propriete}[Le changement de ton de 一 yī]
Le caractère \zh{一} « un » porte le 1er ton (\emph{yī}) lorsqu'il est seul ou dans un nombre (十一 \emph{shíyī} « onze »). Mais son ton change selon la syllabe qui suit :
\begin{itemize}
\item Devant un \textbf{4e ton}, 一 devient \textbf{2e ton} : \zh{一样} \emph{yíyàng} « pareil », \zh{一个} \emph{yí ge} « un ».
\item Devant un \textbf{1er, 2e ou 3e ton}, 一 devient \textbf{4e ton} : \zh{一起} \emph{yìqǐ} « ensemble », \zh{一些} \emph{yìxiē} « quelques ».
\end{itemize}
Attention : à l'écrit, on note souvent le ton d'origine dans les dictionnaires, mais à l'oral et dans le pinyin des manuels, on écrit le ton réellement prononcé.
\end{propriete}

\begin{attention}[Le ton neutre de 子 : qīzi, háizi]
Dans \zh{妻子} \emph{qīzi} et \zh{孩子} \emph{háizi}, la deuxième syllabe 子 prend le \cle{ton neutre} : elle se prononce courte, légère, sans contour tonal. Erreur fréquente des francophones : prononcer \emph{qīzǐ} ou \emph{háizǐ} avec un 3e ton plein. Or \zh{子} seul se lit bien \emph{zǐ} (3e ton), mais en fin de mot familier il s'affaiblit. Même phénomène dans \zh{丈夫} \emph{zhàngfu} : 夫 perd son 1er ton. Retenez : \textbf{1re syllabe pleine + 2e syllabe légère}. Autres exemples déjà connus : \zh{妈妈} \emph{māma}, \zh{我们} \emph{wǒmen}.
\end{attention}

\courssub{Étude des mots clés : emplois et comparaison avec le français}

\textbf{1. 当 dāng : « exercer le métier de ».}\par
Le verbe \zh{当} se construit directement devant un nom de métier : Sujet + 当 + profession. Le français exige une périphrase (« exercer le métier de », « travailler comme »), le chinois est plus direct. Ne le confondez pas avec \zh{做} \emph{zuò} « faire » : on dit \zh{当医生} « être médecin » mais \zh{做家务} « faire le ménage ».

\begin{exemplebox}[当 dāng en action]
\zh{她当医生。} \emph{Tā dāng yīshēng.} « Elle exerce le métier de médecin. »\\
\zh{我姐姐当老师。} \emph{Wǒ jiějie dāng lǎoshī.} « Ma grande sœur est enseignante. »\\
\zh{在杜阿拉，很多妇女当商人。} \emph{Zài Dùālālā, hěn duō fùnǚ dāng shāngrén.} « À Douala, beaucoup de femmes exercent le commerce. »
\end{exemplebox}

\textbf{2. 一起 yìqǐ : « ensemble ».}\par
Adverbe placé \textbf{avant le verbe} : Sujet + 一起 + Verbe. En français, « ensemble » se met en fin de phrase ; en chinois, jamais après le verbe.

\begin{exemplebox}[一起 yìqǐ : place dans la phrase]
\zh{我们一起做家务。} \emph{Wǒmen yìqǐ zuò jiāwù.} « Nous faisons le ménage ensemble. »\\
\zh{丈夫和妻子一起照顾孩子。} \emph{Zhàngfu hé qīzi yìqǐ zhàogù háizi.} « Le mari et l'épouse s'occupent ensemble des enfants. »
\end{exemplebox}

\textbf{3. 还 hái et 也 yě : deux adverbes à ne pas confondre.}\par
\zh{也} signifie « aussi » (addition d'un fait identique) ; \zh{还} signifie « encore » (durée) ou « de plus ». Tous deux se placent avant le verbe, jamais en début de phrase comme le « aussi » français.

\begin{center}
\begin{tabularx}{\linewidth}{|X|X|X|}
\hline
\rowcolor{popblueL} Structure & Exemple (caractères + pinyin) & Traduction \\
\hline
Sujet + 也 + verbe & 她也参加劳动。Tā yě cānjiā láodòng. & Elle aussi participe au travail. \\
\hline
Sujet + 还 + verbe & 这种思想还存在。Zhè zhǒng sīxiǎng hái cúnzài. & Cette mentalité existe encore. \\
\hline
Sujet + 一起 + verbe & 我们一起学习。Wǒmen yìqǐ xuéxí. & Nous étudions ensemble. \\
\hline
但是 + phrase & 但是，问题还存在。Dànshì, wèntí hái cúnzài. & Mais le problème existe encore. \\
\hline
\end{tabularx}
\end{center}

\begin{definition}[重男轻女 zhòngnánqīngnǚ]
Expression figée en quatre caractères (de type \emph{chéngyǔ} 成语) : \zh{重} \emph{zhòng} « considérer comme lourd, important », \zh{男} \emph{nán} « homme, garçon », \zh{轻} \emph{qīng} « considérer comme léger », \zh{女} \emph{nǚ} « femme, fille ». Sens littéral : « considérer les hommes comme importants et les femmes comme légères », c'est-à-dire la \textbf{préférence traditionnelle pour les garçons}. Elle s'emploie comme nom ou adjectif : \zh{重男轻女的思想} \emph{zhòngnánqīngnǚ de sīxiǎng} « la mentalité qui privilégie les garçons ».
\end{definition}

\begin{attention}[Faux amis et pièges de traduction]
\begin{itemize}
\item \zh{妇女} \emph{fùnǚ} est un terme soutenu et collectif (« les femmes » en général, dans les textes officiels) ; dans la conversation, on préfère \zh{女人} \emph{nǚrén} ou \zh{女的}. Ne traduisez pas « une femme » par 妇女 dans une phrase familière.
\item \zh{努力} \emph{nǔlì} peut être verbe (\zh{努力学习} « étudier avec effort ») ou adjectif (\zh{她很努力} « elle est très appliquée »). Ce n'est jamais un nom : « ses efforts » se dira \zh{她的努力}.
\item \zh{一样} \emph{yíyàng} exige la structure A + 和 + B + 一样 (+ adjectif) : \zh{女人和男人一样} « les femmes sont comme les hommes ». Ne dites pas « 女人一样男人 ».
\end{itemize}
\end{attention}

\begin{savaistu}[Le pronom 她 tā « elle » : une invention récente]
À l'oral, \zh{他} (il), \zh{她} (elle) et \zh{它} (il/elle pour les choses) se prononcent tous \emph{tā} : impossible de les distinguer ! À l'écrit, \zh{她} « elle » n'existe que depuis le début du XXe siècle : avant, \zh{他} servait pour les deux genres. C'est le poète Liu Bannong qui a popularisé \zh{她} vers 1920, en plein débat sur la place des femmes dans la société chinoise. Le caractère contient la clé de la femme \zh{女} à gauche — la même clé que dans \zh{妇}, \zh{妻}, \zh{好} ou \zh{妈}. Un caractère peut donc porter l'histoire de l'égalité homme-femme !
\end{savaistu}

\courssub{Texte de réinvestissement}

Lisons maintenant un court texte original qui réemploie le glossaire dans un contexte camerounais. Lisez-le d'abord en silence, puis à voix haute en respectant les tons.

\begin{textebox}[La famille de Mme Etoundi à Yaoundé]
\zh{在雅温得，埃通迪夫人是一位老师。她的丈夫在一家公司工作。他们有一个儿子和一个女儿。}\\
\emph{Zài Yǎwēndé, Àitōngdí fūrén shì yí wèi lǎoshī. Tā de zhàngfu zài yì jiā gōngsī gōngzuò. Tāmen yǒu yí ge érzi hé yí ge nǚ'ér.}\\
« À Yaoundé, Mme Etoundi est enseignante. Son mari travaille dans une entreprise. Ils ont un fils et une fille. »\\
\zh{在家里，他们一起做家务：丈夫做饭，妻子洗衣服，孩子也帮忙。埃通迪夫人说：「男人和女人有一样的权利，家庭里也应该平等。」}\\
\emph{Zài jiā lǐ, tāmen yìqǐ zuò jiāwù : zhàngfu zuòfàn, qīzi xǐ yīfu, háizi yě bāngmáng. Àitōngdí fūrén shuō : « Nánrén hé nǚrén yǒu yíyàng de quánlì, jiātíng lǐ yě yīnggāi píngděng. »}\\
« À la maison, ils font le ménage ensemble : le mari cuisine, l'épouse lave le linge, les enfants aident aussi. Mme Etoundi dit : "Les hommes et les femmes ont les mêmes droits, et dans la famille aussi il doit y avoir l'égalité." »\\
\zh{但是，在一些地方，重男轻女的思想还存在。她认为，只有努力教育孩子，才能真正提高妇女的地位。}\\
\emph{Dànshì, zài yìxiē dìfang, zhòngnánqīngnǚ de sīxiǎng hái cúnzài. Tā rènwéi, zhǐyǒu nǔlì jiàoyù háizi, cái néng zhēnzhèng tígāo fùnǚ de dìwèi.}\\
« Mais dans certaines régions, la mentalité qui valorise les garçons existe encore. Elle pense que ce n'est qu'en faisant des efforts pour éduquer les enfants que l'on pourra vraiment élever le statut des femmes. »
\end{textebox}

Petit vocabulaire d'appoint du texte :

\begin{center}
\begin{tabularx}{\linewidth}{|X|X|X|X|}
\hline
\rowcolor{popblueL} Caractères & Pinyin & Nature & Traduction \\
\hline
公司 & gōngsī & nom & entreprise, société \\
\hline
儿子 & érzi & nom & fils \\
\hline
女儿 & nǚ'ér & nom & fille \\
\hline
做饭 & zuòfàn & verbe & cuisiner, faire à manger \\
\hline
洗衣服 & xǐ yīfu & verbe & laver le linge \\
\hline
帮忙 & bāngmáng & verbe & aider, donner un coup de main \\
\hline
认为 & rènwéi & verbe & penser, estimer \\
\hline
应该 & yīnggāi & verbe modal & devoir (il faut) \\
\hline
教育 & jiàoyù & verbe / nom & éduquer ; éducation \\
\hline
只有……才…… & zhǐyǒu... cái... & structure & seulement si... alors... \\
\hline
\end{tabularx}
\end{center}

Questions de vérification rapide : 1) Quel métier exerce Mme Etoundi ? 2) Qui fait le ménage dans la famille ? 3) Que faut-il faire, selon elle, pour élever le statut des femmes ?

\courssub{Exercice résolu et entraînement}

\begin{exoresolu}[Complétez avec le mot du glossaire qui convient]
Énoncé — Complétez chaque phrase avec l'un des mots suivants : \zh{当}, \zh{一起}, \zh{还}, \zh{地位}, \zh{重男轻女}.\\
1. \zh{我妹妹\_\_\_\_护士。} 2. \zh{我们\_\_\_\_做作业。} 3. \zh{这个问题\_\_\_\_存在。} 4. \zh{妇女的社会\_\_\_\_提高了。} 5. \zh{\_\_\_\_的思想是不正确的。}
\tcblower
Solution détaillée — 1. \zh{当} : devant un nom de métier (护士 \emph{hùshi} « infirmière »), seul 当 convient : \zh{我妹妹当护士。} \emph{Wǒ mèimei dāng hùshi.} « Ma petite sœur est infirmière. » 2. \zh{一起} : adverbe placé avant le verbe 做 : \zh{我们一起做作业。} \emph{Wǒmen yìqǐ zuò zuòyè.} « Nous faisons les devoirs ensemble. » 3. \zh{还} : \zh{这个问题还存在。} \emph{Zhège wèntí hái cúnzài.} « Ce problème existe encore » (durée). 4. \zh{地位} : la collocation est 社会地位 « statut social » : \zh{妇女的社会地位提高了。} \emph{Fùnǚ de shèhuì dìwèi tígāo le.} « Le statut social des femmes s'est élevé. » 5. \zh{重男轻女} : seule expression qui qualifie une 思想 : \zh{重男轻女的思想是不正确的。} \emph{Zhòngnánqīngnǚ de sīxiǎng shì bù zhèngquè de.} « La mentalité qui privilégie les garçons n'est pas correcte. »
\end{exoresolu}

\begin{exercice}[À vous de jouer]
1. Traduisez en français : \zh{妻子和丈夫一起参加社会工作。} / \zh{真正的平等还不存在。}\\
2. Traduisez en chinois (avec pinyin) : « Les femmes ont les mêmes droits que les hommes. » (utilisez \zh{一样}) ; « Elle s'occupe des enfants et fait aussi les tâches ménagères. »\\
3. Classez ces mots dans les trois colonnes noms / verbes / adjectifs-expressions : \zh{权利}, \zh{照顾}, \zh{平等}, \zh{家务}, \zh{提高}, \zh{真正}.\\
4. Prononciation : indiquez les tons de chaque syllabe dans \zh{妇女}, \zh{努力}, \zh{一样}, \zh{孩子}, puis lisez ces mots à voix haute.
\end{exercice}

\begin{corrige}[Exercice « À vous de jouer »]
1. « L'épouse et le mari participent ensemble au travail social. » / « La véritable égalité n'existe pas encore. »\\
2. \zh{女人和男人有一样的权利。} \emph{Nǚrén hé nánrén yǒu yíyàng de quánlì.} ; \zh{她照顾孩子，也做家务。} \emph{Tā zhàogù háizi, yě zuò jiāwù.}\\
3. Noms : 权利, 家务. Verbes : 照顾, 提高. Adjectifs et expressions : 平等, 真正.\\
4. \zh{妇女} \emph{fùnǚ} : 4e + 3e ton ; \zh{努力} \emph{nǔlì} : 3e + 4e ton ; \zh{一样} \emph{yíyàng} : 2e + 4e ton (changement de ton de 一) ; \zh{孩子} \emph{háizi} : 2e ton + ton neutre.
\end{corrige}

\begin{aretenir}[L'essentiel du glossaire]
\begin{itemize}
\item Trois familles de mots : \textbf{noms} (妇女, 地位, 社会, 家庭, 权利, 国策, 思想, 家务, 丈夫, 妻子, 孩子, 劳动), \textbf{verbes} (参加, 当, 照顾, 存在, 努力, 提高), \textbf{adjectifs et expressions} (平等, 一样, 真正, 重男轻女), plus les mots-outils 但是, 也, 还, 一起.
\item Tons piégeux : \emph{qīzi}, \emph{háizi}, \emph{zhàngfu} avec ton neutre ; \emph{nǔlì} 3e + 4e ton ; \emph{yíyàng} et \emph{yìqǐ} avec changement de ton de 一.
\item Collocations à retenir : 社会地位 « statut social », 做家务 « faire le ménage », 照顾孩子 « s'occuper des enfants », 提高地位 « élever le statut », 参加劳动 « participer au travail », 重男轻女的思想 « la mentalité de préférence pour les garçons ».
\item Structures clés : Sujet + 当 + métier ; Sujet + 一起 + verbe ; A + 和 + B + 一样 ; 只有……才…… « seulement si... alors... ».
\end{itemize}
\end{aretenir}

\coursec{Clés (radicaux) et ordre des traits}

Dans la section précédente, nous avons dressé le glossaire du thème « la situation de la femme et l'égalité homme-femme » : \zh{妇女} \emph{fùnǚ} « femmes », \zh{平等} \emph{píngděng} « égalité », \zh{家庭} \emph{jiātíng} « famille », \zh{丈夫} \emph{zhàngfu} « mari », \zh{妻子} \emph{qīzi} « épouse »… Vous avez sans doute remarqué que certains caractères se ressemblent ou partagent un même petit dessin. Ce n'est pas un hasard ! Les caractères chinois ne sont pas des dessins arbitraires : ils sont construits à partir de \cle{clés} (aussi appelées \cle{radicaux}), qui donnent un indice de sens, et de parties phonétiques, qui donnent un indice de prononciation. Comprendre ce mécanisme, c'est transformer l'apprentissage par cœur en apprentissage intelligent. C'est l'objet de cette section.

\courssub{Qu'est-ce qu'une clé (radical) ?}

\begin{definition}[La clé ou radical]
Une \cle{clé} (部首, \emph{bùshǒu}) est un élément graphique qui se répète dans de nombreux caractères et qui indique généralement la \textbf{famille de sens} du caractère. C'est aussi l'élément sous lequel le caractère est classé dans les dictionnaires chinois. Un caractère est souvent composé d'une clé (le sens) et d'une partie phonétique (la prononciation approximative).
\end{definition}

Observez d'abord ces caractères de notre glossaire :

\begin{exemplebox}[Observation : un même dessin, un même champ de sens]
\zh{她} \emph{tā} « elle » \quad \zh{妇} \emph{fù} « femme (mariée) » \quad \zh{妻} \emph{qī} « épouse » \quad \zh{好} \emph{hǎo} « bon » \quad \zh{姓} \emph{xìng} « nom de famille »\\[4pt]
Tous contiennent l'élément \zh{女} \emph{nǚ} « femme ». Et tous ont un rapport avec la femme ou le féminin !
\end{exemplebox}

\begin{propriete}[La clé 女 (\emph{nǚ}) : l'indice de la féminité]
La clé \zh{女} « femme » indique très souvent un lien avec la femme, le féminin ou la parenté féminine :
\begin{itemize}
\item \zh{妈} \emph{mā} : maman
\item \zh{姐} \emph{jiě} : grande sœur
\item \zh{妹} \emph{mèi} : petite sœur
\item \zh{妇} \emph{fù} : femme mariée
\item \zh{她} \emph{tā} : elle (pronom féminin)
\item \zh{姓} \emph{xìng} : nom de famille (hérité autrefois de la mère !)
\end{itemize}
Attention toutefois : la clé donne un \emph{indice} de sens, pas une définition exacte. \zh{好} \emph{hǎo} « bon » contient \zh{女} mais ne veut pas dire « féminin » : c'est une \emph{association d'idées} (voir le schéma ci-dessous).
\end{propriete}

\courssub{Quatre clés essentielles de notre thème}

\textbf{1. La clé 人 / 亻 (\emph{rén}) : la personne}\par
Sous sa forme complète, \zh{人} \emph{rén} signifie « personne, homme (au sens humain) ». Placée à gauche d'un caractère, elle se comprime en \zh{亻} (appelé « personne debout »). Elle indique un lien avec la personne humaine :

\begin{itemize}
\item \zh{他} \emph{tā} : il, lui (\zh{亻} + \zh{也})
\item \zh{们} \emph{men} : suffixe du pluriel des personnes → \zh{我们} \emph{wǒmen} « nous », \zh{他们} \emph{tāmen} « ils »
\item \zh{作} \emph{zuò} : faire, travailler (une activité humaine)
\item \zh{位} \emph{wèi} : classificateur poli pour les personnes → \zh{一位女士} \emph{yí wèi nǚshì} « une dame »
\end{itemize}

\textbf{2. La clé 力 (\emph{lì}) : la force}\par
\zh{力} dessine un bras musclé : c'est la force physique. On la retrouve dans :
\begin{itemize}
\item \zh{男} \emph{nán} : homme (masculin) = \zh{田} \emph{tián} « champ » + \zh{力} \emph{lì} « force » : celui qui met sa force aux champs. Image de la société agraire ancienne !
\item \zh{努} \emph{nǔ} : faire des efforts → \zh{努力} \emph{nǔlì} « s'efforcer, travailler dur »
\item \zh{动} \emph{dòng} : bouger, mouvement → \zh{劳动} \emph láodòng « travail »
\end{itemize}

\textbf{3. La clé 宀 (\emph{mián}) : le toit}\par
\zh{宀} représente un toit de maison. Elle indique un lien avec l'habitation et la vie domestique :
\begin{itemize}
\item \zh{家} \emph{jiā} : maison, famille (un toit qui abrite \zh{豕}, le cochon : la richesse du foyer ancien)
\item \zh{安} \emph{ān} : paix, tranquillité = \zh{宀} toit + \zh{女} femme : une femme sous le toit, c'est la paix du foyer.
\end{itemize}

\begin{popfigure}
\begin{tikzpicture}[font=\small, node distance=0.5cm]
% 好
\node[draw=popblue, fill=popblueL, rounded corners, minimum width=1.1cm, minimum height=0.9cm] (nv) at (0,0) {\zh{女}};
\node[below=0.05cm of nv, text=popdark] {femme};
\node at (1.5,0) {\textbf{+}};
\node[draw=popblue, fill=popblueL, rounded corners, minimum width=1.1cm, minimum height=0.9cm] (zi) at (2.4,0) {\zh{子}};
\node[below=0.05cm of zi, text=popdark] {enfant};
\node at (3.9,0) {\textbf{→}};
\node[draw=popgreen, fill=popgreenL, rounded corners, minimum width=1.3cm, minimum height=0.9cm] (hao) at (5.3,0) {\zh{好}};
\node[below=0.05cm of hao, text=popdark] {\emph{hǎo} « bon »};
% 安
\node[draw=poporange, fill=poporangeL, rounded corners, minimum width=1.1cm, minimum height=0.9cm] (mian) at (0,-2) {\zh{宀}};
\node[below=0.05cm of mian, text=popdark] {toit};
\node at (1.5,-2) {\textbf{+}};
\node[draw=poporange, fill=poporangeL, rounded corners, minimum width=1.1cm, minimum height=0.9cm] (nv2) at (2.4,-2) {\zh{女}};
\node[below=0.05cm of nv2, text=popdark] {femme};
\node at (3.9,-2) {\textbf{→}};
\node[draw=popgreen, fill=popgreenL, rounded corners, minimum width=1.3cm, minimum height=0.9cm] (an) at (5.3,-2) {\zh{安}};
\node[below=0.05cm of an, text=popdark] {\emph{ān} « paix »};
% 男
\node[draw=poppurple, fill=poppurpleL, rounded corners, minimum width=1.1cm, minimum height=0.9cm] (tian) at (0,-4) {\zh{田}};
\node[below=0.05cm of tian, text=popdark] {champ};
\node at (1.5,-4) {\textbf{+}};
\node[draw=poppurple, fill=poppurpleL, rounded corners, minimum width=1.1cm, minimum height=0.9cm] (li) at (2.4,-4) {\zh{力}};
\node[below=0.05cm of li, text=popdark] {force};
\node at (3.9,-4) {\textbf{→}};
\node[draw=popgreen, fill=popgreenL, rounded corners, minimum width=1.3cm, minimum height=0.9cm] (nan) at (5.3,-4) {\zh{男}};
\node[below=0.05cm of nan, text=popdark] {\emph{nán} « homme »};
\end{tikzpicture}
\legende{Décomposition de trois caractères composés : l'association des sens crée le sens global.}
\end{popfigure}

\begin{aretenir}[Tableau récapitulatif des clés du thème]
\begin{center}
\begin{tabularx}{\linewidth}{|X|X|X|X|}\hline
\rowcolor{popblueL} Clé & Sens de la clé & Exemples & Indice donné \\ \hline
\zh{女} \emph{nǚ} & femme & \zh{妇} \zh{妻} \zh{她} \zh{好} \zh{姓} \zh{妈} \zh{姐} \zh{妹} & féminité, parenté féminine \\ \hline
\zh{人} / \zh{亻} \emph{rén} & personne & \zh{他} \zh{们} \zh{作} \zh{位} & personne humaine \\ \hline
\zh{力} \emph{lì} & force & \zh{男} \zh{努} \zh{动} & force, effort, action \\ \hline
\zh{宀} \emph{mián} & toit & \zh{家} \zh{安} & maison, vie domestique \\ \hline
\end{tabularx}
\end{center}
\end{aretenir}

\begin{savaistu}[La clé 女 et les préjugés anciens]
Certains caractères à clé \zh{女} ont un sens péjoratif : \zh{妒} \emph{dù} « jalousie », \zh{奸} \emph{jiān} « traître, perfide », \zh{嫌} \emph{xián} « suspect ». Pourquoi associer la femme à ces défauts ? Ces caractères ont été créés dans une société profondément patriarcale, et ils témoignent des préjugés de leur époque. Aujourd'hui, un débat existe en Chine : faut-il réformer ces caractères ? Pour l'instant, l'écriture reste inchangée, mais la discussion montre que la langue porte l'histoire d'une société — exactement le sujet de notre texte support sur l'évolution de la condition féminine !
\end{savaistu}

\courssub{Les règles d'ordre des traits}

Écrire un caractère chinois, ce n'est pas dessiner dans le désordre : chaque trait s'écrit dans un ordre précis, codifié. Respecter cet ordre rend l'écriture plus rapide, plus belle et plus facile à mémoriser.

\begin{propriete}[Les quatre règles fondamentales de l'ordre des traits]
\begin{enumerate}
\item \textbf{De haut en bas} : on écrit d'abord la partie supérieure du caractère. Exemple : \zh{男} (\zh{田} en haut, puis \zh{力} en bas).
\item \textbf{De gauche à droite} : on écrit d'abord la partie gauche. Exemple : \zh{好} (\zh{女} à gauche, puis \zh{子} à droite).
\item \textbf{L'extérieur avant l'intérieur} : pour les caractères avec une enceinte, on trace d'abord le contour, puis le contenu (et on « ferme la porte » en dernier).
\item \textbf{L'horizontal avant le vertical} : quand un trait horizontal et un trait vertical se croisent, l'horizontal passe d'abord. Exemple : \zh{十} \emph{shí} « dix » : d'abord \zh{一}, puis \zh{丨}.
\end{enumerate}
\end{propriete}

\courssub{Ordre des traits démontré, caractère par caractère}

\textbf{1. \zh{女} \emph{nǚ} « femme » — 3 traits}\par
Trait 1 : \zh{撇点} \emph{piědiǎn}, un mouvement unique qui descend en oblique vers la gauche puis repart en point vers la droite (c'est le trait le plus difficile : il se fait \emph{sans lever le stylo}). Trait 2 : \zh{撇} \emph{piě}, une oblique qui descend vers la gauche en traversant le premier trait. Trait 3 : \zh{横} \emph{héng}, l'horizontal, tracé de gauche à droite en dernier.

\textbf{2. \zh{好} \emph{hǎo} « bon » — 6 traits}\par
Règle appliquée : gauche avant droite. On écrit d'abord \zh{女} (3 traits, comme ci-dessus), puis \zh{子} (3 traits) : trait 4, le crochet horizontal \zh{了} en haut ; trait 5, le vertical courbé \zh{丿} ; trait 6, l'horizontal final \zh{一}.

\textbf{3. \zh{男} \emph{nán} « homme » — 7 traits}\par
Règle appliquée : haut avant bas. On écrit d'abord \zh{田} (5 traits) : le vertical gauche \zh{丨}, le pli horizontal-vertical \zh{┐}, puis l'horizontal \zh{一} et le vertical \zh{丨} intérieurs (l'horizontal avant le vertical !), puis l'horizontal du bas qui « ferme » \zh{一}. Ensuite \zh{力} (2 traits) : le pli \zh{┐}, puis l'oblique \zh{丿}.

\textbf{4. \zh{妻} \emph{qī} « épouse » — 8 traits}\par
Règle appliquée : haut avant bas. Partie supérieure (5 traits) : les trois horizontaux \zh{一} de haut en bas, le vertical \zh{丨} qui les traverse, puis l'horizontal long \zh{一}. Partie inférieure : \zh{女} (3 traits) écrit comme appris ci-dessus.

\textbf{5. \zh{家} \emph{jiā} « famille, maison » — 10 traits}\par
Règle appliquée : haut avant bas, extérieur avant intérieur. D'abord le toit \zh{宀} (3 traits : le point \zh{丶}, le point vertical gauche \zh{丿}, le crochet horizontal \zh{乛}), puis le corps \zh{豕} (7 traits) sous le toit.

\begin{methode}[Comment mémoriser durablement un caractère]
\begin{enumerate}
\item \textbf{Identifier la clé} : quel indice de sens ? (\zh{女} → féminin ; \zh{宀} → maison ; \zh{力} → force.)
\item \textbf{Identifier la partie phonétique} : donne-t-elle une piste de prononciation ? Exemple : \zh{妈} \emph{mā} contient \zh{马} \emph{mǎ} « cheval » — le sens vient de \zh{女}, le son vient de \zh{马}.
\item \textbf{Se raconter l'histoire du caractère} : \zh{安} = « une femme sous le toit, c'est la paix ». Une image mentale vaut dix répétitions !
\item \textbf{Écrire en comptant les traits à voix haute} : « un, deux, trois… » en respectant l'ordre, 5 fois minimum par caractère.
\item \textbf{Relire le lendemain} : réécrire le caractère de mémoire, puis vérifier.
\end{enumerate}
\end{methode}

\begin{attention}[Pièges fréquents des francophones]
\begin{itemize}
\item \textbf{\zh{女} s'écrit en 3 traits, pas 4 !} Le premier trait \zh{撇点} \emph{piědiǎn} est un \emph{seul} mouvement continu. Si vous le coupez en deux, le caractère est mal formé et compte 4 traits : c'est une faute classique à l'examen.
\item \textbf{Ne pas confondre \zh{妻} \emph{qī} « épouse » et \zh{妾} \emph{qiè} « concubine »} : les deux contiennent \zh{女} en bas, mais la partie supérieure diffère. \zh{妾} est un mot ancien et péjoratif, à ne jamais employer pour parler d'une épouse aujourd'hui !
\item \textbf{Ne pas confondre \zh{他} \emph{tā} « il » et \zh{她} \emph{tā} « elle »} : même prononciation, mais clés différentes (\zh{亻} personne / \zh{女} femme). À l'écrit, le choix de la clé change le sens ; à l'oral, le contexte suffit.
\item \textbf{Respecter « haut avant bas » dans \zh{男}} : beaucoup d'élèves écrivent \zh{力} avant de finir \zh{田}. Erreur ! On termine toujours le haut d'abord.
\end{itemize}
\end{attention}

\begin{exoresolu}[Identifier les clés et compter les traits]
Énoncé : pour chaque caractère, indique la clé, le nombre de traits et l'indice de sens : (a) \zh{妹} (b) \zh{努} (c) \zh{安} (d) \zh{们}.
\tcblower
Solution détaillée :
\begin{itemize}
\item (a) \zh{妹} \emph{mèi} « petite sœur » : clé \zh{女} (féminité) ; 8 traits (\zh{女} = 3 + \zh{未} = 5) ; la partie \zh{未} \emph{wèi} donne une piste phonétique (m/w proches à l'origine).
\item (b) \zh{努} \emph{nǔ} « effort » : clé \zh{力} (force) ; 7 traits (\zh{奴} = 5 + \zh{力} = 2) ; le sens vient de la force, le son de \zh{奴} \emph{nú}.
\item (c) \zh{安} \emph{ān} « paix » : clé \zh{宀} (toit) ; 6 traits (\zh{宀} = 3 + \zh{女} = 3) ; association d'idées : la femme sous le toit.
\item (d) \zh{们} \emph{men} « pluriel des personnes » : clé \zh{亻} (personne) ; 5 traits (\zh{亻} = 2 + \zh{门} = 3) ; la partie \zh{门} \emph{mén} donne le son.
\end{itemize}
\end{exoresolu}

\begin{exercice}[Atelier d'écriture]
\begin{enumerate}
\item Trace chaque caractère \textbf{5 fois} sur ton cahier, en comptant les traits à voix haute et en respectant l'ordre : \zh{女} \emph{nǚ} (3 traits), \zh{好} \emph{hǎo} (6), \zh{男} \emph{nán} (7), \zh{妻} \emph{qī} (8), \zh{家} \emph{jiā} (10).
\item Décompose les caractères suivants en clé + autre partie, et explique le sens comme dans le schéma de la leçon : \zh{妈} \emph{mā}, \zh{姓} \emph{xìng}, \zh{动} \emph{dòng}, \zh{位} \emph{wèi}.
\item Trouve dans le glossaire de la section précédente deux autres caractères contenant la clé \zh{女} et explique leur lien avec le thème de la leçon.
\end{enumerate}
\end{exercice}

\begin{corrige}[Exercice « Atelier d'écriture »]
\begin{enumerate}
\item Vérification en classe : le professeur contrôle l'ordre des traits au tableau. Points de vigilance : le premier trait de \zh{女} en un seul mouvement ; \zh{田} terminé avant \zh{力} dans \zh{男} ; le toit \zh{宀} avant le corps dans \zh{家}.
\item \zh{妈} \emph{mā} = \zh{女} (sens : femme) + \zh{马} \emph{mǎ} (phonétique) → « maman ». \zh{姓} \emph{xìng} = \zh{女} (sens) + \zh{生} \emph{shēng} (phonétique) → « nom de famille ». \zh{动} \emph{dòng} = \zh{云} \emph{yún} (phonétique) + \zh{力} (sens : force) → « bouger ». \zh{位} \emph{wèi} = \zh{亻} (sens : personne) + \zh{立} \emph{lì} (phonétique) → « place, classificateur de personnes ».
\item Par exemple : \zh{妇} \emph{fù} (dans \zh{妇女} « les femmes ») et \zh{她} \emph{tā} « elle » : tous deux renvoient directement au sujet du texte, la place des femmes dans la société.
\end{enumerate}
\end{corrige}

Vous savez désormais « lire dans » les caractères : derrière chaque dessin se cachent une clé de sens, souvent une piste de prononciation, et parfois toute une histoire de la société chinoise. Cette compétence va nous servir immédiatement dans la section suivante, où nous reprendrons le texte support pour une compréhension détaillée, phrase par phrase, et où nous étudierons les collocations importantes du thème de l'égalité homme-femme.

\coursec{Compréhension détaillée et collocations}

Nous savons maintenant reconnaître les clés (radicaux) et tracer dans le bon ordre les caractères importants du thème, comme \zh{女} (clé « femme »), \zh{男} (clé « champ » 田 + « force » 力) ou \zh{家} (clé « toit » 宀). Il est temps de revenir sur le texte support pour le comprendre \cle{en détail} : répondre à des questions précises, dégager les idées principales, puis réutiliser les \cle{collocations} (associations de mots qui vont naturellement ensemble) dans nos propres phrases. C'est l'étape qui transforme une lecture passive en compétence active — et c'est exactement ce que l'examen attend de vous en compréhension de l'écrit.

\courssub{Questions de compréhension sur le texte}

\begin{methode}[Comment répondre à une question de compréhension ?]
Voici la démarche en quatre étapes, valable pour tout texte chinois au lycée :
\begin{enumerate}
\item \textbf{Lire la question d'abord} : souligner le mot-clé de la question (par exemple \zh{基本国策} « politique de base de l'État »).
\item \textbf{Repérer ce mot-clé dans le texte} : le chinois est une langue très régulière, la réponse se trouve presque toujours dans la phrase qui contient le même mot.
\item \textbf{Recopier la phrase concernée} et vérifier qu'elle répond bien à la question posée.
\item \textbf{Reformuler brièvement} : répondre par une phrase complète et courte, en reprenant les mots du texte, sans tout recopier.
\end{enumerate}
En français comme en chinois, une réponse complète vaut toujours mieux qu'un mot isolé : on écrira \zh{中国的基本国策是男女平等。} plutôt que simplement \zh{男女平等}。
\end{methode}

\begin{exemplebox}[Questions en français, avec les réponses attendues]
\begin{enumerate}
\item \textbf{Quelle est la politique de base de la Chine ?} — L'égalité entre les hommes et les femmes (\zh{男女平等}) est une politique de base de l'État chinois.
\item \textbf{Que peuvent faire les femmes ?} — Elles peuvent aller à l'école, travailler et participer aux activités sociales.
\item \textbf{Quels métiers exercent-elles ?} — Elles sont enseignantes, médecins, ingénieures, ouvrières, commerçantes, et certaines occupent des postes de responsabilité.
\item \textbf{Que font les couples à la maison ?} — Le mari et l'épouse partagent les tâches : ils font le ménage et s'occupent des enfants ensemble.
\item \textbf{Quel problème existe encore dans les campagnes ?} — Dans certaines campagnes, la mentalité de « préférence pour les garçons » (\zh{重男轻女的思想}) n'a pas complètement disparu ; il faut continuer les efforts.
\end{enumerate}
\end{exemplebox}

\begin{exemplebox}[Les mêmes questions en chinois, avec réponses guidées]
\zh{中国的基本国策是什么？} \emph{Zhōngguó de jīběn guócè shì shénme ?} « Quelle est la politique de base de la Chine ? » \\
→ \zh{中国的基本国策是男女平等。} \emph{Zhōngguó de jīběn guócè shì nánnǚ píngděng.} « C'est l'égalité homme-femme. » \\[4pt]
\zh{妇女可以做什么？} \emph{Fùnǚ kěyǐ zuò shénme ?} « Que peuvent faire les femmes ? » \\
→ \zh{她们可以上学、工作、参加社会活动。} \emph{Tāmen kěyǐ shàngxué, gōngzuò, cānjiā shèhuì huódòng.} « Elles peuvent aller à l'école, travailler et participer aux activités sociales. » \\[4pt]
\zh{妇女做什么工作？} \emph{Fùnǚ zuò shénme gōngzuò ?} « Quels métiers les femmes exercent-elles ? » \\
→ \zh{她们是老师、医生、工程师，也有女经理。} \emph{Tāmen shì lǎoshī, yīshēng, gōngchéngshī, yě yǒu nǚ jīnglǐ.} « Elles sont enseignantes, médecins, ingénieures, et il y a aussi des femmes cadres. » \\[4pt]
\zh{在家里，夫妻做什么？} \emph{Zài jiāli, fūqī zuò shénme ?} « À la maison, que font les couples ? » \\
→ \zh{他们一起做家务、照顾孩子。} \emph{Tāmen yìqǐ zuò jiāwù, zhàogù háizi.} « Ils font le ménage et s'occupent des enfants ensemble. » \\[4pt]
\zh{农村还有什么问题？} \emph{Nóngcūn hái yǒu shénme wèntí ?} « Quel problème existe encore dans les campagnes ? » \\
→ \zh{一些地方还有重男轻女的思想。} \emph{Yìxiē dìfang hái yǒu zhòngnánqīngnǚ de sīxiǎng.} « Dans certains endroits, la mentalité de préférence pour les garçons existe encore. »
\end{exemplebox}

\begin{attention}[La particule interrogative \zh{什么}]
En chinois, on ne change pas l'ordre des mots pour poser une question : le mot interrogatif \zh{什么} (shénme, « quoi / que ») reste \textbf{à la place de la réponse}, contrairement au français qui place « que » en tête. Comparez : \\
\zh{妇女可以做什么？} (littéralement : « Les femmes peuvent faire quoi ? ») = « Que peuvent faire les femmes ? » \\
Erreur fréquente des francophones : écrire \zh{什么妇女可以做？} par calque du français. C'est incorrect : le sujet reste en tête de phrase. \\
Même logique avec les autres mots interrogatifs : \zh{谁} (shéi, « qui »), \zh{哪儿} (nǎr, « où »), \zh{什么时候} (shénme shíhou, « quand ») restent tous à la place du mot qu'ils remplacent.
\end{attention}

\begin{attention}[Deux adverbes du texte à ne pas confondre : \zh{还} et \zh{也}]
\zh{还} (hái, 2\textsuperscript{e} ton) signifie « encore, toujours » : \zh{农村还有问题。} \emph{Nóngcūn hái yǒu wèntí.} « Les campagnes ont encore des problèmes. » \\
\zh{也} (yě) signifie « aussi » : \zh{爸爸也照顾孩子。} \emph{Bàba yě zhàogù háizi.} « Papa aussi s'occupe des enfants. » \\
Les deux se placent \textbf{avant le verbe}, après le sujet. Ne confondez pas leurs tons : hái (montant) / yě (descendant-montant).
\end{attention}

\courssub{Les collocations essentielles du thème}

Une \cle{collocation} est un couple de mots qui s'emploie presque toujours ensemble : en français on dit « faire le ménage » et non « fabriquer le ménage ». En chinois, apprendre le verbe avec son complément habituel évite les contresens et fait gagner un temps précieux à l'écrit comme à l'oral.

\begin{center}
\begin{tabularx}{\linewidth}{|X|X|X|X|}
\hline
\rowcolor{popblueL} Caractères & Pinyin & Nature & Traduction \\
\hline
\zh{提高地位} & tígāo dìwèi & verbe + nom & élever le statut \\
\hline
\zh{男女平等} & nánnǚ píngděng & expression nominale & égalité homme-femme \\
\hline
\zh{参加社会活动} & cānjiā shèhuì huódòng & verbe + complément & participer aux activités sociales \\
\hline
\zh{做家务} & zuò jiāwù & verbe + complément & faire le ménage \\
\hline
\zh{照顾孩子} & zhàogù háizi & verbe + complément & s'occuper des enfants \\
\hline
\zh{重男轻女的思想} & zhòngnánqīngnǚ de sīxiǎng & expression nominale & mentalité de préférence pour les garçons \\
\hline
\end{tabularx}
\end{center}

\begin{exemplebox}[Chaque collocation dans une phrase personnelle]
\zh{教育可以提高妇女的地位。} \emph{Jiàoyù kěyǐ tígāo fùnǚ de dìwèi.} « L'éducation peut élever le statut des femmes. » \\[4pt]
\zh{男女平等是中国的基本国策。} \emph{Nánnǚ píngděng shì Zhōngguó de jīběn guócè.} « L'égalité homme-femme est une politique de base de la Chine. » \\[4pt]
\zh{我妈妈经常参加社会活动。} \emph{Wǒ māma jīngcháng cānjiā shèhuì huódòng.} « Ma mère participe souvent aux activités sociales. » \\[4pt]
\zh{在家里，丈夫和妻子一起做家务。} \emph{Zài jiāli, zhàngfu hé qīzi yìqǐ zuò jiāwù.} « À la maison, le mari et l'épouse font le ménage ensemble. » \\[4pt]
\zh{爸爸也照顾孩子。} \emph{Bàba yě zhàogù háizi.} « Papa aussi s'occupe des enfants. » \\[4pt]
\zh{我们要反对重男轻女的思想。} \emph{Wǒmen yào fǎnduì zhòngnánqīngnǚ de sīxiǎng.} « Nous devons combattre la mentalité de préférence pour les garçons. »
\end{exemplebox}

\begin{attention}[\zh{参加} : participer à, mais pas n'importe comment]
\zh{参加} (cānjiā) s'emploie pour les \textbf{activités, réunions, organisations et événements collectifs} : \zh{参加会议} (assister à une réunion), \zh{参加比赛} (participer à une compétition), \zh{参加社会活动}. \\
Deux pièges pour les francophones :
\begin{itemize}
\item « Passer un examen » se dit \zh{参加考试} (cānjiā kǎoshì) : ici \zh{参加} convient, car il s'agit de « se présenter à » une épreuve collective.
\item En revanche, « fréquenter une école » ne se dit pas \zh{参加学校} : on dit \zh{上学} (aller à l'école). Retenez : \zh{参加} + activité ponctuelle ou collective, jamais + lieu.
\end{itemize}
\end{attention}

\courssub{Exprimer l'égalité : la structure avec \zh{一样}}

Observons d'abord deux phrases :

\begin{exemplebox}[Observation]
\zh{女人和男人一样聪明。} \emph{Nǚrén hé nánrén yíyàng cōngming.} « Les femmes sont aussi intelligentes que les hommes. » \\[4pt]
\zh{女儿和儿子一样重要。} \emph{Nǚ'ér hé érzi yíyàng zhòngyào.} « La fille est aussi importante que le fils. »
\end{exemplebox}

Que remarque-t-on ? Le mot \zh{一样} (yíyàng, « pareil, identique ») se place \textbf{après} le couple comparé relié par \zh{和} (hé, « et/avec »), et \textbf{avant} l'adjectif. C'est l'équivalent chinois de la comparaison d'égalité française « aussi ... que ». Remarquez aussi le pinyin : \zh{一} seul est au premier ton (yī), mais il passe au deuxième ton (yí) devant un quatrième ton — d'où \emph{yíyàng}.

\begin{propriete}[La comparaison d'égalité : A 和 B 一样 (+ adjectif)]
\textbf{Structure :} A + \zh{和} + B + \zh{一样} (+ adjectif) = « A est aussi ... que B ». \\
\textbf{Emploi :} pour affirmer que deux personnes ou choses sont au même niveau sur un point (intelligence, importance, travail, droits). \\
\textbf{Négation :} on place \zh{不} devant \zh{一样} : A + \zh{和} + B + \zh{不一样} = « A n'est pas comme B / A et B ne sont pas pareils ». \\
\textbf{Variante :} on peut remplacer \zh{和} par \zh{跟} (gēn) ou \zh{与} (yǔ, plus écrit) sans changer le sens. \\
\textbf{Sans adjectif :} \zh{A 和 B 一样} seul signifie « A et B sont pareils » : \zh{我的看法和你的一样。} \emph{Wǒ de kànfǎ hé nǐ de yíyàng.} « Mon avis est le même que le tien. »
\end{propriete}

\begin{center}
\begin{tabularx}{\linewidth}{|X|X|X|}
\hline
\rowcolor{popblueL} Structure & Exemple (caractères + pinyin) & Traduction \\
\hline
A 和 B 一样 + adj. & \zh{女人和男人一样聪明。} Nǚrén hé nánrén yíyàng cōngming. & Les femmes sont aussi intelligentes que les hommes. \\
\hline
A 和 B 一样 + adj. & \zh{女孩和男孩一样能上学。} Nǚhái hé nánhái yíyàng néng shàngxué. & Les filles peuvent aller à l'école comme les garçons. \\
\hline
A 和 B 不一样 & \zh{以前和现在不一样。} Yǐqián hé xiànzài bù yíyàng. & Autrefois et aujourd'hui, ce n'est pas pareil. \\
\hline
A 跟 B 一样 + adj. & \zh{妻子跟丈夫一样工作。} Qīzi gēn zhàngfu yíyàng gōngzuò. & L'épouse travaille comme le mari. \\
\hline
\end{tabularx}
\end{center}

\begin{attention}[Ordre des mots : le piège du calque français]
Le français dit « aussi intelligent \emph{que} les hommes » : le mot de comparaison vient \emph{après} l'adjectif. En chinois, c'est l'inverse : \zh{一样} vient \emph{avant} l'adjectif. Ne dites jamais \zh{女人聪明和男人一样} : la seule possibilité est \zh{女人和男人一样聪明}. Retenez le schéma : \textbf{A 和 B 一样 + adjectif}, sans rien intervertir. \\
Autre piège : ne confondez pas avec la comparaison de supériorité \zh{比} (A 比 B + adjectif = « A est plus ... que B »). \zh{一样} exprime l'\textbf{égalité}, \zh{比} exprime la \textbf{supériorité} : \zh{女人和男人一样聪明} (égalité) ≠ \zh{女人比男人聪明} (supériorité).
\end{attention}

\courssub{Carte de compréhension du texte}

\begin{popfigure}
\begin{tikzpicture}[
  node distance=0.55cm and 1.1cm,
  theme/.style={rectangle, rounded corners, draw=popdark, fill=popblueL, align=center, font=\small\bfseries, minimum width=4.2cm, minimum height=0.9cm},
  idea/.style={rectangle, rounded corners, draw=popblue, fill=popgreenL, align=center, font=\small, minimum width=3.4cm, minimum height=1.1cm},
  prob/.style={rectangle, rounded corners, draw=poporange, fill=poporangeL, align=center, font=\small, minimum width=3.4cm, minimum height=1.1cm},
  concl/.style={rectangle, rounded corners, draw=poppurple, fill=poppurpleL, align=center, font=\small\bfseries, minimum width=4.6cm, minimum height=0.9cm},
  fleche/.style={-{Stealth[length=2.5mm]}, thick, draw=popmuted}
]
\node[theme, align=center] (t){Thème : 男女平等\\ \emph{l'égalité homme-femme}};
\node[idea, below left=0.7cm and -0.6cm of t, align=center] (i1){Progrès :\\妇女的地位提高了\\ \emph{le statut s'est élevé}};
\node[idea, below=0.7cm of t, align=center] (i2){Égalité légale :\\基本国策\\ \emph{politique de base}};
\node[idea, below right=0.7cm and -0.6cm of t, align=center] (i3){Partage des tâches :\\一起做家务\\ \emph{ménage ensemble}};
\node[prob, below=2.6cm of i2, align=center] (p){Problème restant : 农村\\重男轻女的思想\\ \emph{préférence pour les garçons}};
\node[concl, below=0.6cm of p, align=center] (c){Conclusion : 继续努力\\ \emph{continuer les efforts}};
\draw[fleche] (t) -- (i1);
\draw[fleche] (t) -- (i2);
\draw[fleche] (t) -- (i3);
\draw[fleche] (i1.south) |- (p.west);
\draw[fleche] (i2.south) -- (p.north);
\draw[fleche] (i3.south) |- (p.east);
\draw[fleche] (p) -- (c);
\end{tikzpicture}
\legende{Organigramme de compréhension : du thème aux idées principales, au problème restant et à la conclusion.}
\end{popfigure}

\courssub{Reformulation : résumer le texte en trois phrases}

Savoir résumer est une compétence clé de l'expression écrite en Terminale. La règle : une phrase par idée principale, avec les mots du texte, sans exemple ni détail. Notez au passage l'emploi de \zh{了} dans \zh{地位提高了} : il marque ici un \textbf{changement d'état} (le statut, autrefois bas, s'est élevé), et non le passé composé français.

\begin{exoresolu}[Résume le texte en trois phrases simples en chinois]
Énoncé : à l'aide de l'organigramme ci-dessus, écris un résumé du texte en trois phrases : (1) la situation actuelle, (2) le partage des tâches, (3) le problème qui reste.
\tcblower
Solution détaillée : on reprend les trois étages de l'organigramme et on choisit les collocations apprises. \\[4pt]
1. Situation actuelle : \zh{现在，中国妇女的地位提高了，男女平等是基本国策。} \emph{Xiànzài, Zhōngguó fùnǚ de dìwèi tígāo le, nánnǚ píngděng shì jīběn guócè.} « Aujourd'hui, le statut des femmes chinoises s'est élevé et l'égalité homme-femme est une politique de base. » \\[4pt]
2. Partage des tâches : \zh{在家里，丈夫和妻子一起做家务、照顾孩子。} \emph{Zài jiāli, zhàngfu hé qīzi yìqǐ zuò jiāwù, zhàogù háizi.} « À la maison, le mari et l'épouse font le ménage et s'occupent des enfants ensemble. » \\[4pt]
3. Problème restant : \zh{但是，一些农村还有重男轻女的思想，我们要继续努力。} \emph{Dànshì, yìxiē nóngcūn hái yǒu zhòngnánqīngnǚ de sīxiǎng, wǒmen yào jìxù nǔlì.} « Mais dans certaines campagnes, la mentalité de préférence pour les garçons existe encore ; nous devons continuer les efforts. »
\end{exoresolu}

\begin{exoresolu}[Traduction guidée (version) : du français vers le chinois]
Énoncé : traduis en chinois la phrase suivante, en utilisant la structure de comparaison d'égalité : « Au Cameroun aussi, les femmes sont aussi importantes que les hommes dans la famille. » (Cameroun = \zh{喀麦隆} Kāmàilóng ; dans la famille = \zh{在家里}).
\tcblower
Solution détaillée, étape par étape : \\
1. « Au Cameroun aussi » → complément de lieu en tête + adverbe \zh{也} avant le verbe : \zh{在喀麦隆，女人也……} \\
2. « sont aussi importantes que les hommes » → structure A 和 B 一样 + adjectif : \zh{和男人一样重要} \\
3. « dans la famille » → complément de lieu avant le verbe : \zh{在家里} \\
Phrase complète : \zh{在喀麦隆，女人在家里也和男人一样重要。} \emph{Zài Kāmàilóng, nǚrén zài jiāli yě hé nánrén yíyàng zhòngyào.} \\
Vérification : sujet (\zh{女人}) → adverbe (\zh{也}) → comparaison (\zh{和男人一样}) → adjectif (\zh{重要}). L'ordre est correct.
\end{exoresolu}

\begin{exercice}[À toi de jouer]
\begin{enumerate}
\item Réponds en chinois par une phrase complète : \zh{妇女可以做什么工作？}
\item Traduis en chinois avec \zh{一样} : « Les filles sont aussi travailleuses que les garçons. » (travailleur = \zh{努力} nǔlì)
\item Construis une phrase personnelle avec la collocation \zh{参加社会活动}.
\item Complète : \zh{在家里，丈夫和妻子一起\_\_\_\_\_\_、\_\_\_\_\_\_孩子。}
\item Traduis en français : \zh{以前和现在不一样，现在妇女的地位提高了。}
\end{enumerate}
\end{exercice}

\begin{corrige}[Exercice « À toi de jouer »]
1. \zh{她们可以做老师、医生、工程师。} \emph{Tāmen kěyǐ zuò lǎoshī, yīshēng, gōngchéngshī.} « Elles peuvent être enseignantes, médecins, ingénieures. » (Toute réponse citant des métiers du texte est acceptée.) \\
2. \zh{女孩和男孩一样努力。} \emph{Nǚhái hé nánhái yíyàng nǔlì.} Attention à l'ordre : \zh{一样} avant l'adjectif. \\
3. Exemple : \zh{我们班的女生喜欢参加社会活动。} \emph{Wǒmen bān de nǚshēng xǐhuan cānjiā shèhuì huódòng.} « Les filles de notre classe aiment participer aux activités sociales. » \\
4. \zh{在家里，丈夫和妻子一起做家务、照顾孩子。} \emph{Zài jiāli, zhàngfu hé qīzi yìqǐ zuò jiāwù, zhàogù háizi.} \\
5. « Autrefois et aujourd'hui, ce n'est pas pareil : maintenant, le statut des femmes s'est élevé. » (Noter \zh{了} qui marque le changement d'état.)
\end{corrige}

\begin{aretenir}[L'essentiel de la section]
\begin{itemize}
\item Pour répondre à une question de compréhension : repérer le mot-clé, retrouver la phrase du texte, reformuler en phrase complète.
\item Six collocations à connaître par cœur : \zh{提高地位}, \zh{男女平等}, \zh{参加社会活动}, \zh{做家务}, \zh{照顾孩子}, \zh{重男轻女的思想}.
\item Comparaison d'égalité : \textbf{A 和 B 一样 + adjectif} ; négation : \zh{不一样} ; à ne pas confondre avec la supériorité \zh{比}.
\item \zh{参加} s'emploie avec les activités et réunions (\zh{参加考试} = passer un examen), jamais avec un lieu.
\item \zh{还} (hái, « encore ») et \zh{也} (yě, « aussi ») se placent avant le verbe ; \zh{了} peut marquer un changement d'état (\zh{地位提高了}).
\item Un résumé = une phrase par idée principale, avec les mots du texte.
\end{itemize}
\end{aretenir}

\coursec{Production guidée et ouverture Cameroun-Chine}

Après avoir compris le texte en détail et mémorisé les collocations essentielles (\zh{提高地位} \emph{tígāo dìwèi} « élever le statut », \zh{男女平等} \emph{nánnǚ píngděng} « égalité homme-femme », \zh{分担家务} \emph{fēndān jiāwù} « partager les tâches ménagères »), il est temps de passer à la \cle{production} : c'est vous qui allez maintenant écrire et parler en chinois sur la situation de la femme au Cameroun. C'est l'étape qui transforme un vocabulaire reconnu en un vocabulaire réellement acquis — et c'est exactement ce type de tâche qui vous sera demandé à l'examen (expression écrite guidée et traduction).

\courssub{Production écrite guidée : cinq phrases sur la femme camerounaise}

\textbf{Principe de la production guidée.}\par

En production guidée, on ne part pas d'une page blanche : on part d'une \cle{trame} (un texte à trous) dont la structure est déjà correcte. Votre travail consiste à choisir les bons mots du glossaire, à les placer au bon endroit et à vérifier trois choses : les \cle{tons}, l'\cle{ordre des mots} et l'écriture des \cle{caractères}. Observez d'abord la trame, lisez-la à voix haute, puis seulement après, complétez-la.

\begin{textebox}[Trame de production — 喀麦隆妇女的情况]
现在，\_\_\_\_妇女的地位越来越高。她们可以\_\_\_\_、\_\_\_\_、\_\_\_\_。但是，在一些地方，\_\_\_\_还存在。我们应该\_\_\_\_。\\
\emph{Xiànzài, \_\_\_\_ fùnǚ de dìwèi yuèláiyuè gāo. Tāmen kěyǐ \_\_\_\_, \_\_\_\_, \_\_\_\_. Dànshì, zài yìxiē dìfang, \_\_\_\_ hái cúnzài. Wǒmen yīnggāi \_\_\_\_.}\\
« Aujourd'hui, le statut des femmes \_\_\_\_ est de plus en plus élevé. Elles peuvent \_\_\_\_, \_\_\_\_, \_\_\_\_. Mais, dans certains endroits, \_\_\_\_ existe encore. Nous devons \_\_\_\_. »
\end{textebox}

\textbf{Vocabulaire d'appoint pour compléter la trame.}\par

\begin{center}
\begin{tabularx}{\linewidth}{|X|X|X|X|}
\hline
\rowcolor{popblueL} Caractères & Pinyin & Nature & Traduction \\
\hline
喀麦隆 & Kāmàilóng & nom propre & le Cameroun \\
\hline
农村 & nóngcūn & nom & la campagne, le monde rural \\
\hline
城市 & chéngshì & nom & la ville \\
\hline
上学 & shàngxué & verbe & aller à l'école \\
\hline
工作 & gōngzuò & verbe / nom & travailler ; le travail \\
\hline
参加 & cānjiā & verbe & participer à \\
\hline
歧视 & qíshì & verbe / nom & discriminer ; la discrimination \\
\hline
保护 & bǎohù & verbe & protéger \\
\hline
权利 & quánlì & nom & le droit, les droits \\
\hline
\end{tabularx}
\end{center}

\begin{methode}[Comment réussir une production guidée en cinq étapes]
\begin{enumerate}
\item \textbf{Lire} la trame entière sans écrire, pour comprendre le sens global et repérer la nature du mot attendu à chaque trou (nom ? verbe ? adjectif ?).
\item \textbf{Choisir} dans le glossaire les mots qui conviennent au sens et à la grammaire : après \zh{可以} il faut un verbe ; après \zh{在一些地方} il faut un nom ou une phrase complète.
\item \textbf{Rédiger} les cinq phrases en caractères, en recopiant soigneusement chaque caractère (ordre des traits : haut → bas, gauche → droite, horizontal avant vertical).
\item \textbf{Ajouter le pinyin} avec les marques de tons au-dessus des voyelles porteuses de ton (ā á ǎ à).
\item \textbf{Relire} en vérifiant : ordre Sujet + Verbe + Complément ; place de \zh{在} avant le complément de lieu ; présence de \zh{的} entre le nom et son complément du nom (\zh{妇女的地位} \emph{fùnǚ de dìwèi} « le statut des femmes »).
\end{enumerate}
\end{methode}

\begin{exemplebox}[Un modèle de production complétée]
现在，喀麦隆妇女的地位越来越高。她们可以上学、工作、参加政治活动。但是，在一些农村地区，歧视还存在。我们应该保护妇女的权利。\\
\emph{Xiànzài, Kāmàilóng fùnǚ de dìwèi yuèláiyuè gāo. Tāmen kěyǐ shàngxué, gōngzuò, cānjiā zhèngzhì huódòng. Dànshì, zài yìxiē nóngcūn dìqū, qíshì hái cúnzài. Wǒmen yīnggāi bǎohù fùnǚ de quánlì.}\\
« Aujourd'hui, le statut des femmes camerounaises est de plus en plus élevé. Elles peuvent aller à l'école, travailler, participer aux activités politiques. Mais, dans certaines régions rurales, la discrimination existe encore. Nous devons protéger les droits des femmes. »
\end{exemplebox}

\begin{attention}[Trois pièges fréquents des francophones dans cette trame]
\begin{itemize}
\item \textbf{La place du complément de lieu} : en français on dit « la discrimination existe encore \emph{dans certaines campagnes} » (lieu à la fin). En chinois, le lieu avec \zh{在} se place \emph{avant} le verbe : \zh{在一些地方，歧视还存在} et jamais \zh{歧视还存在在一些地方}.
\item \textbf{La série de verbes} : pour énumérer trois actions après \zh{可以}, on sépare par le signe d'énumération chinois 、 (et non par une virgule) : \zh{上学、工作、参加政治活动}.
\item \textbf{Les tons de 喀麦隆} : Kā (premier ton) — mài (quatrième ton) — lóng (deuxième ton). Beaucoup d'élèves prononcent tout au premier ton : entraînez-vous à dire le nom de votre pays correctement !
\end{itemize}
\end{attention}

\courssub{Échange et correction collective}

Une fois vos cinq phrases rédigées, la classe procède à une \cle{correction collective} : deux ou trois élèves écrivent leur production au tableau, les autres comparent avec la leur. La vérification porte sur trois points précis, dans cet ordre :

\begin{enumerate}
\item \textbf{Les caractères} : sont-ils complets et bien formés ? (Vérifier en particulier \zh{歧视}, \zh{地位}, \zh{权利}, qui contiennent beaucoup de traits.)
\item \textbf{L'ordre des mots} : Sujet + (complément de temps / de lieu) + Verbe + Complément. Le mot \zh{还} (\emph{hái} « encore ») se place juste avant le verbe : \zh{还存在} « existe encore ».
\item \textbf{Les tons} : lecture à voix haute phrase par phrase ; la classe corrige la prononciation, notamment \zh{妇女} fùnǚ (quatrième + troisième ton) et \zh{平等} píngděng (deuxième + troisième ton).
\end{enumerate}

\begin{exoresolu}[Corriger une phrase d'élève]
Énoncé : un élève a écrit : \zh{妇女可以工作在城市。} (\emph{Fùnǚ kěyǐ gōngzuò zài chéngshì.}) Cette phrase est-elle correcte ? Si non, corrigez-la et expliquez.
\tcblower
Solution : la phrase est \textbf{incorrecte}. En chinois, le complément de lieu introduit par \zh{在} se place \emph{avant} le verbe, et non après comme en français (« travailler en ville »). La phrase correcte est : \zh{妇女可以在城市工作。} (\emph{Fùnǚ kěyǐ zài chéngshì gōngzuò.}) « Les femmes peuvent travailler en ville. » Structure : Sujet (\zh{妇女}) + \zh{可以} + lieu (\zh{在城市}) + verbe (\zh{工作}). Attention : si le verbe est \zh{在} lui-même (« être à, se trouver à »), le lieu vient après : \zh{她在城市。} (\emph{Tā zài chéngshì.}) « Elle est en ville. »
\end{exoresolu}

\courssub{Mini-débat oral guidé : 男女平等重要吗？}

Passons maintenant à l'oral. Question du débat : \zh{男女平等重要吗？} (\emph{Nánnǚ píngděng zhòngyào ma?}) « L'égalité homme-femme est-elle importante ? » Pour donner votre avis, vous disposez de trois formules :

\begin{center}
\begin{tabularx}{\linewidth}{|X|X|X|}
\hline
\rowcolor{popblueL} Formule & Exemple en caractères + pinyin & Traduction \\
\hline
我觉得 + phrase & 我觉得男女平等很重要。Wǒ juéde nánnǚ píngděng hěn zhòngyào. & Je pense que l'égalité homme-femme est très importante. \\
\hline
我同意 & 我同意你的看法。Wǒ tóngyì nǐ de kànfǎ. & Je suis d'accord avec ton point de vue. \\
\hline
我不同意，因为…… & 我不同意，因为农村还有很多问题。Wǒ bù tóngyì, yīnwèi nóngcūn hái yǒu hěn duō wèntí. & Je ne suis pas d'accord, parce qu'à la campagne il y a encore beaucoup de problèmes. \\
\hline
\end{tabularx}
\end{center}

\begin{attention}[觉得 ou 想 ?]
\zh{觉得} (\emph{juéde}) exprime une \textbf{opinion personnelle}, un jugement : « je pense que, je trouve que ». Il est toujours suivi d'une \textbf{phrase complète} : \zh{我觉得男女平等很重要。} Ne le confondez pas avec \zh{想} (\emph{xiǎng}), qui signifie « vouloir » devant un verbe (\zh{我想学汉语。} \emph{Wǒ xiǎng xué Hànyǔ.} « Je veux apprendre le chinois. ») ou « penser à ». Dire \zh{我觉得学汉语} pour « je veux apprendre le chinois » est une erreur : employez \zh{我想学汉语}.
\end{attention}

\begin{experience}[Débat en deux camps]
La classe se divise en deux groupes. Le groupe A défend l'importance de l'égalité (\zh{男女平等很重要。} « L'égalité homme-femme est très importante. »), le groupe B nuance ou discute les difficultés (\zh{但是，在农村还有很多问题。} « Mais, à la campagne, il y a encore beaucoup de problèmes. »). Chaque élève prend la parole au moins une fois en commençant par \zh{我觉得...}, puis réagit à un camarade avec \zh{我同意} ou \zh{我不同意，因为...}. Le professeur note au tableau les arguments en chinois : \zh{教育} \emph{jiàoyù} (éducation), \zh{工作} \emph{gōngzuò} (travail), \zh{家庭} \emph{jiātíng} (famille), \zh{传统} \emph{chuántǒng} (tradition).
\end{experience}

\courssub{Ouverture Cameroun-Chine : regard comparatif}

Pour clore la leçon, comparons les deux pays que vous étudiez ensemble depuis la Première. Le schéma suivant, à recopier et à commenter oralement, résume quelques points de comparaison factuels.

\begin{popfigure}
\begin{tikzpicture}[font=\small]
\draw[fill=popblueL, draw=popblue] (0,0) rectangle (5.2,1);
\draw[fill=poporangeL, draw=poporange] (5.4,0) rectangle (10.6,1);
\node at (2.6,0.5) {\textbf{中国 Zhōngguó (Chine)}};
\node at (8,0.5) {\textbf{喀麦隆 Kāmàilóng (Cameroun)}};
\draw[fill=popcream, draw=popline] (0,-1.2) rectangle (5.2,0);
\draw[fill=popcream, draw=popline] (5.4,-1.2) rectangle (10.6,0);
\node[align=center] at (2.6,-0.6) {女孩都上学\\ (scolarisation des filles)};
\node[align=center] at (8,-0.6) {女孩上学越来越多\\ (progrès de la scolarisation)};
\draw[fill=popcream, draw=popline] (0,-2.4) rectangle (5.2,-1.2);
\draw[fill=popcream, draw=popline] (5.4,-2.4) rectangle (10.6,-1.2);
\node[align=center] at (2.6,-1.8) {妇女有很多工作机会\\ (accès aux métiers)};
\node[align=center] at (8,-1.8) {妇女经商、从政\\ (commerce et politique)};
\draw[fill=popcream, draw=popline] (0,-3.6) rectangle (5.2,-2.4);
\draw[fill=popcream, draw=popline] (5.4,-3.6) rectangle (10.6,-2.4);
\node[align=center] at (2.6,-3) {城市里分担家务\\ (partage des tâches en ville)};
\node[align=center] at (8,-3) {农村家务多是妇女做\\ (tâches surtout rurales)};
\draw[fill=popcream, draw=popline] (0,-4.8) rectangle (5.2,-3.6);
\draw[fill=popcream, draw=popline] (5.4,-4.8) rectangle (10.6,-3.6);
\node[align=center] at (2.6,-4.2) {挑战：传统观念\\ (défis : traditions)};
\node[align=center] at (8,-4.2) {挑战：早婚、贫困\\ (défis : mariages précoces)};
\end{tikzpicture}
\legende{Tableau comparatif : la situation des femmes en Chine et au Cameroun (éducation, métiers, partage des tâches, défis).}
\end{popfigure}

\begin{savaistu}[Le Cameroun, la Chine et la CEDAW]
La \textbf{CEDAW} (Convention sur l'élimination de toutes les formes de discrimination à l'égard des femmes) est le grand traité international des droits des femmes. La \textbf{Chine} l'a ratifiée en \textbf{1980}, et le \textbf{Cameroun} en \textbf{1994}. En chinois, on dit souvent que les femmes \zh{能顶半边天} (\emph{néng dǐng bànbiāntiān}) : « peuvent soutenir la moitié du ciel ». L'expression \zh{半边天} (« la moitié du ciel ») désigne aujourd'hui les femmes elles-mêmes : l'employer dans une copie de traduction ou d'expression écrite est un \textbf{atout à l'examen} qui montre une vraie culture chinoise.
\end{savaistu}

\courssub{Synthèse de la leçon : carte mentale à compléter}

Pour fixer durablement le vocabulaire du thème, complétez cette carte mentale en y ajoutant au moins deux mots par branche, tirés du glossaire de la leçon.

\begin{popfigure}
\begin{tikzpicture}[mindmap, grow cyclic, every node/.style={concept, align=center, font=\small},
root concept/.append style={concept color=popblue, font=\small\bfseries},
level 1/.append style={level distance=3.4cm, sibling angle=90}]
\node[root concept, text=white, align=center]{妇女\\ fùnǚ\\ la femme}
child[concept color=poporange] { node[align=center]{地位\\ dìwèi\\ statut} }
child[concept color=popgreen] { node[align=center]{平等\\ píngděng\\ égalité} }
child[concept color=poppurple] { node[align=center]{权利\\ quánlì\\ droits} }
child[concept color=poppink] { node[align=center]{挑战\\ tiǎozhàn\\ défis} };
\end{tikzpicture}
\legende{Carte mentale du thème « la femme et l'égalité » : à compléter avec le vocabulaire de la leçon (ex. 教育, 工作, 家务, 歧视, 保护...).}
\end{popfigure}

\begin{aretenir}[Ce qu'il faut retenir de cette section]
\begin{itemize}
\item Une production guidée se réussit en cinq étapes : lire, choisir, rédiger, ajouter le pinyin, relire.
\item Le complément de lieu avec \zh{在} se place \textbf{avant} le verbe ; les énumérations se séparent par le signe 、.
\item Pour débattre : \zh{我觉得 + phrase} (opinion), \zh{我同意} / \zh{我不同意，因为...}.
\item Expression d'examen : \zh{半边天} (\emph{bànbiāntiān}), « la moitié du ciel », pour parler des femmes avec élégance.
\end{itemize}
\end{aretenir}

\begin{exercice}[À faire à la maison]
Rédigez votre propre version de la trame en parlant de votre ville ou village (Yaoundé, Douala, Buea...) : cinq phrases en caractères avec pinyin, en utilisant au moins une fois \zh{城市} ou \zh{农村}, une fois \zh{但是}, et une fois \zh{我们应该}. Préparez-vous à la lire à voix haute au prochain cours.
\end{exercice}

\newpage

\coursec{Activité d'intégration}

\courssub{Objectifs de l'activité}

À l'issue de cette activité d'intégration, l'élève sera capable de :
\begin{itemize}
\item \cle{réinvestir} le vocabulaire du thème « situation de la femme et égalité homme-femme » dans une production écrite authentique et créative ;
\item \cle{écrire correctement} les caractères du glossaire (ordre des traits, clés) et les prononcer avec les bons tons ;
\item \cle{adapter} des phrases du texte support à un nouveau contexte de communication ;
\item \cle{produire à l'oral} une courte présentation avec la structure d'opinion \zh{我们觉得...} (wǒmen juéde... « nous pensons que... ») ;
\item \cle{comparer} de manière simple et respectueuse la situation de la femme au Cameroun et en Chine.
\end{itemize}

\courssub{Situation}

Le club « Chinois et citoyenneté » de ton lycée de Yaoundé prépare la \cle{Journée internationale des droits des femmes} (8 mars). La principale a demandé au club de décorer le tableau d'affichage du lycée avec des affiches de sensibilisation \textbf{bilingues chinois-français} sur le thème de l'égalité entre les hommes et les femmes. La meilleure affiche sera reproduite et envoyée à l'école partenaire du lycée à Chengdu (Chine), dans le cadre du programme d'échange scolaire.

Ton groupe a été choisi pour représenter la classe de Terminale A. Pour vous inspirer, la professeure vous remet le message suivant, envoyé par les correspondants chinois de l'école partenaire :

\begin{textebox}[Message de l'école partenaire de Chengdu]
亲爱的喀麦隆朋友们：\\
Qīn'ài de Kāmàilóng péngyoumen :\\
« Chers amis camerounais, »\\
我们听说你们要制作关于男女平等的海报，太好了！\\
Wǒmen tīngshuō nǐmen yào zhìzuò guānyú nánnǚ píngděng de hǎibào, tài hǎo le !\\
« Nous avons entendu dire que vous allez créer une affiche sur l'égalité homme-femme, c'est formidable ! »\\
在中国，人们常说“妇女能顶半边天”。\\
Zài Zhōngguó, rénmen cháng shuō “fùnǚ néng dǐng bàn biān tiān”.\\
« En Chine, on dit souvent : “Les femmes peuvent soutenir la moitié du ciel.” »\\
我们很想知道喀麦隆的情况。祝你们成功！\\
Wǒmen hěn xiǎng zhīdào Kāmàilóng de qíngkuàng. Zhù nǐmen chénggōng !\\
« Nous aimerions beaucoup connaître la situation au Cameroun. Nous vous souhaitons du succès ! »
\end{textebox}

\courssub{Consignes}

Travail en groupes de 3 ou 4 élèves. Durée totale : 1 h 30 (préparation) + 30 min (présentations et vote).

\begin{enumerate}
\item \textbf{Relecture du texte support} (10 min). Relisez le texte de la leçon sur la situation de la femme. Surlignez \textbf{trois phrases} que vous pourriez reprendre ou adapter pour votre affiche. Notez-les avec leur pinyin et leur traduction.
\item \textbf{Choix du slogan} (10 min). Composez un slogan court et percutant en chinois, par exemple \zh{男女平等，人人有责} (nánnǚ píngděng, rénrén yǒu zé) « Égalité homme-femme, affaire de tous ». Vérifiez les tons de chaque syllabe à voix haute.
\item \textbf{Sélection du vocabulaire} (15 min). Choisissez \textbf{au moins 8 mots du glossaire} de la leçon (par exemple \zh{平等}, \zh{权利}, \zh{教育}, \zh{工作}, \zh{尊重}, \zh{社会}, \zh{家庭}, \zh{发展}). Écrivez chaque mot en respectant l'ordre des traits et les clés étudiées, puis ajoutez le pinyin avec les tons et la traduction française.
\item \textbf{Rédaction des phrases} (20 min). Rédigez \textbf{trois phrases} tirées ou adaptées du texte support (ex. : \zh{女人和男人有一样的权利。} Nǚrén hé nánrén yǒu yíyàng de quánlì. « Les femmes et les hommes ont les mêmes droits. »). Attention à l'ordre des mots : Sujet + Verbe + Complément.
\item \textbf{Comparaison Chine-Cameroun} (10 min). Ajoutez \textbf{une phrase} comparant les deux pays, en utilisant \zh{和...一样} ou \zh{在...，...} (ex. : \zh{在中国和在喀麦隆，女孩都能上学。} Zài Zhōngguó hé zài Kāmàilóng, nǚhái dōu néng shàngxué. « En Chine comme au Cameroun, les filles peuvent aller à l'école. »).
\item \textbf{Mise en page de l'affiche} (15 min). Disposez le slogan en haut, les phrases au centre, le vocabulaire en encadré. Soignez l'écriture des caractères : haut avant bas, gauche avant droite, horizontal avant vertical.
\item \textbf{Présentation orale} (1 min par groupe). Présentez votre affiche en commençant par \zh{我们觉得...} (wǒmen juéde...). Chaque membre du groupe doit prononcer au moins une phrase en chinois.
\item \textbf{Vote de la classe} (10 min). La classe vote pour la meilleure affiche selon les critères : exactitude des caractères, correction des tons à l'oral, richesse du vocabulaire, clarté du message.
\end{enumerate}

\courssub{Ressources à mobiliser}

\begin{aretenir}[Boîte à outils de l'affiche]
\begin{itemize}
\item \textbf{Structure d'opinion} : \zh{我们觉得 + phrase} (wǒmen juéde...) « nous pensons que... ». Ex. : \zh{我们觉得男女应该平等。} Wǒmen juéde nánnǚ yīnggāi píngděng. « Nous pensons que les hommes et les femmes doivent être égaux. »
\item \textbf{Comparaison} : A \zh{和} B \zh{一样} + adjectif (A hé B yíyàng...) « A et B sont pareillement... » ; \zh{在} + lieu « dans/en + lieu ».
\item \textbf{Verbes modaux utiles} : \zh{能} néng « pouvoir », \zh{应该} yīnggāi « devoir », \zh{要} yào « devoir, vouloir ».
\item \textbf{Lexique du thème} : \zh{男女平等} nánnǚ píngděng « égalité homme-femme », \zh{权利} quánlì « droit », \zh{尊重} zūnzhòng « respecter », \zh{教育} jiàoyù « éducation », \zh{机会} jīhuì « chance, opportunité », \zh{责任} zérèn « responsabilité ».
\item \textbf{Proverbe chinois} : \zh{妇女能顶半边天。} Fùnǚ néng dǐng bàn biān tiān. « Les femmes peuvent soutenir la moitié du ciel. »
\end{itemize}
\end{aretenir}

\begin{attention}[Pièges fréquents]
\begin{itemize}
\item Ne traduis pas mot à mot « égalité homme-femme » : l'ordre chinois est \zh{男女平等} (homme-femme-égalité), jamais \zh{平等男女}.
\item \zh{一样} yíyàng s'emploie avec \zh{和} : on dit \zh{女人和男人一样} et non \zh{女人一样男人}.
\item Attention aux tons : \zh{nǚ} (女, 3e ton, ü avec tréma) ne se prononce pas comme \zh{nú} ; \zh{quánlì} (权利, 2e-4e ton) ne doit pas devenir \zh{quānlì}.
\item Le caractère \zh{女} (3 traits : le trait brisé d'abord, puis le jeté, puis l'horizontal) est la clé de \zh{妇}, \zh{好}, \zh{妈} : respecte l'ordre des traits.
\end{itemize}
\end{attention}

\courssub{Proposition de production et corrigé commenté}

\begin{exoresolu}[Modèle d'affiche et de présentation orale]
Énoncé : rédigez l'affiche complète du groupe (slogan, 3 phrases, 8 mots du glossaire, 1 comparaison) et le texte de la présentation orale d'une minute.
\tcblower
\textbf{Affiche modèle :}\\[4pt]
Slogan : \zh{男女平等，人人有责！} Nánnǚ píngděng, rénrén yǒu zé ! « Égalité homme-femme, affaire de tous ! »\\[4pt]
Phrase 1 : \zh{女人和男人有一样的权利。} Nǚrén hé nánrén yǒu yíyàng de quánlì. « Les femmes et les hommes ont les mêmes droits. »\\
Phrase 2 : \zh{女孩和男孩都应该上学。} Nǚhái hé nánhái dōu yīnggāi shàngxué. « Les filles et les garçons doivent tous aller à l'école. »\\
Phrase 3 : \zh{我们要尊重每一个人。} Wǒmen yào zūnzhòng měi yí ge rén. « Nous devons respecter chaque personne. »\\[4pt]
Comparaison : \zh{在中国和在喀麦隆，女人都能工作、学习。} Zài Zhōngguó hé zài Kāmàilóng, nǚrén dōu néng gōngzuò, xuéxí. « En Chine comme au Cameroun, les femmes peuvent travailler et étudier. »\\[4pt]
Encadré vocabulaire (8 mots) : \zh{平等} píngděng « égalité » ; \zh{权利} quánlì « droit » ; \zh{教育} jiàoyù « éducation » ; \zh{工作} gōngzuò « travail » ; \zh{尊重} zūnzhòng « respecter » ; \zh{社会} shèhuì « société » ; \zh{家庭} jiātíng « famille » ; \zh{机会} jīhuì « opportunité ».\\[6pt]
\textbf{Présentation orale modèle :}\\
\zh{大家好！我们觉得男女平等很重要。女人和男人有一样的权利，女孩和男孩都应该上学。在中国和在喀麦隆，女人都能工作、学习。男女平等，人人有责！谢谢大家！}\\
Dàjiā hǎo ! Wǒmen juéde nánnǚ píngděng hěn zhòngyào. Nǚrén hé nánrén yǒu yíyàng de quánlì, nǚhái hé nánhái dōu yīnggāi shàngxué. Zài Zhōngguó hé zài Kāmàilóng, nǚrén dōu néng gōngzuò, xuéxí. Nánnǚ píngděng, rénrén yǒu zé ! Xièxie dàjiā !\\
« Bonjour à tous ! Nous pensons que l'égalité homme-femme est très importante. Les femmes et les hommes ont les mêmes droits, les filles et les garçons doivent tous aller à l'école. En Chine comme au Cameroun, les femmes peuvent travailler et étudier. Égalité homme-femme, affaire de tous ! Merci à tous ! »\\[6pt]
\textbf{Commentaire des choix de langue :}
\begin{itemize}
\item Le slogan reprend la structure en deux membres de quatre caractères, typique des slogans chinois : rythme équilibré, facile à mémoriser.
\item La phrase 1 applique la comparaison avec \zh{和...一样} ; la phrase 2 utilise le modal \zh{应该} et l'adverbe \zh{都} placé avant le verbe, conformément à l'ordre des mots chinois.
\item La comparaison Chine-Cameroun répète \zh{在} devant chaque pays, ce qui est plus clair et plus correct qu'une seule occurrence.
\item La présentation orale s'ouvre par \zh{我们觉得}, structure d'opinion exigée, et se clôt par la formule de politesse \zh{谢谢大家}, adaptée à un exposé devant la classe.
\end{itemize}
\end{exoresolu}

\courssub{Bilan de l'activité}

Cette activité t'a permis de mobiliser l'ensemble des savoir-faire de la leçon dans une tâche authentique et concrète :
\begin{itemize}
\item \textbf{compréhension de l'écrit} : relire le texte support et le message des correspondants pour y puiser des idées et des phrases ;
\item \textbf{lexique} : sélectionner, écrire et prononcer correctement au moins huit mots du glossaire, en respectant clés, traits et tons ;
\item \textbf{grammaire} : employer les structures \zh{我们觉得...}, \zh{和...一样}, les verbes modaux \zh{能}, \zh{应该}, \zh{要} et l'adverbe \zh{都} ;
\item \textbf{production écrite et orale} : concevoir un document réel destiné à un public réel (le lycée et l'école partenaire de Chengdu) et le présenter en une minute ;
\item \textbf{interculturalité} : formuler une comparaison simple et factuelle entre le Cameroun et la Chine sur la place des femmes dans la société.
\end{itemize}

\begin{savaistu}[Le proverbe de « la moitié du ciel »]
L'expression \zh{妇女能顶半边天} (fùnǚ néng dǐng bàn biān tiān, « les femmes peuvent soutenir la moitié du ciel ») est devenue en Chine le symbole de la participation des femmes à la vie économique et sociale. Au Cameroun aussi, les femmes jouent un rôle central : commerçantes au marché Mokolo ou au marché central de Douala, enseignantes, médecins, ministres. Dire \zh{男女平等} dans les deux langues, c'est reconnaître que le développement d'un pays a besoin de « tout le ciel » : les hommes \emph{et} les femmes.
\end{savaistu}

\newpage

\coursec{Méthodes et exercices résolus}

\courssub{Méthodes et exercices résolus}

\begin{methode}[Lire et comprendre un court texte chinois sur le thème]
Pour réussir la compréhension de l'écrit au Bac, suis toujours la même démarche en cinq étapes.
\begin{enumerate}
\item \textbf{Lecture silencieuse globale} : lis le texte une première fois sans dictionnaire. Repère le thème (ici : la situation de la femme, l'égalité homme-femme) et les mots déjà connus.
\item \textbf{Repérage des mots-clés du thème} : souligne les mots du champ lexical : \zh{妇女} (fùnǚ, femmes), \zh{平等} (píngděng, égalité), \zh{权利} (quánlì, droits), \zh{工作} (gōngzuò, travail), \zh{家庭} (jiātíng, famille).
\item \textbf{Découpage des phrases} : identifie pour chaque phrase le sujet, le verbe, le complément. En chinois, l'ordre est presque toujours Sujet + (temps/lieu) + Verbe + Complément.
\item \textbf{Lecture des questions avant la relecture} : lis les questions, puis relis le texte en cherchant uniquement les réponses. Les réponses apparaissent presque toujours dans l'ordre du texte.
\item \textbf{Rédaction de la réponse} : réponds par une phrase complète en chinois, en reprenant les mots du texte. Ne traduis jamais mot à mot depuis le français.
\end{enumerate}
Réflexe à retenir : en chinois, une information de temps ou de lieu se place \textbf{avant le verbe}, jamais à la fin comme en français.
\end{methode}

\begin{textebox}[Texte support : 李老师的家庭]
李老师是喀麦隆的一位女老师。她在雅温得的一所中学教汉语。\\
Lǐ lǎoshī shì Kāmàilóng de yí wèi nǚ lǎoshī. Tā zài Yǎwēndé de yì suǒ zhōngxué jiāo Hànyǔ.\\
« Madame Li est une enseignante camerounaise. Elle enseigne le chinois dans un lycée de Yaoundé. »

她和丈夫都在工作，也一起做家务。他们的女儿和儿子都上学。\\
Tā hé zhàngfu dōu zài gōngzuò, yě yìqǐ zuò jiāwù. Tāmen de nǚ'ér hé érzi dōu shàngxué.\\
« Son mari et elle travaillent tous les deux, et font aussi le ménage ensemble. Leur fille et leur fils vont tous deux à l'école. »

李老师说：「女人和男人一样，都有学习和工作的权利。」\\
Lǐ lǎoshī shuō : « Nǚrén hé nánrén yíyàng, dōu yǒu xuéxí hé gōngzuò de quánlì. »\\
« Madame Li dit : les femmes sont comme les hommes, elles ont toutes le droit d'étudier et de travailler. »
\end{textebox}

\begin{exoresolu}[Compréhension de texte : 李老师的家庭]
Énoncé. Relis le texte support \zh{李老师的家庭} ci-dessus, puis réponds aux questions en chinois par des phrases complètes.

Questions :
\begin{enumerate}
\item 李老师做什么工作？(Lǐ lǎoshī zuò shénme gōngzuò ?)
\item 谁做家务？(Shéi zuò jiāwù ?)
\item 李老师的女儿上学吗？(Lǐ lǎoshī de nǚ'ér shàngxué ma ?)
\item 女人有什么权利？(Nǚrén yǒu shénme quánlì ?)
\end{enumerate}
\tcblower
Solution détaillée.

\textbf{Question 1.} La question porte sur le métier. On cherche dans le texte le mot \zh{工作} ou le nom de métier. Première phrase : \zh{李老师是喀麦隆的一位女老师} — « Madame Li est une enseignante ». Attention au classificateur : \zh{一位女老师} (\zh{位} est le classificateur respectueux des personnes).

Réponse rédigée : \zh{她是一位老师，在雅温得的一所中学教汉语。} (Tā shì yí wèi lǎoshī, zài Yǎwēndé de yì suǒ zhōngxué jiāo Hànyǔ.) « Elle est enseignante, elle enseigne le chinois dans un lycée de Yaoundé. »

\textbf{Question 2.} Le mot interrogatif \zh{谁} (shéi, qui) demande une personne. Deuxième phrase : \zh{她和丈夫都在工作，也一起做家务} — le sujet du ménage est « elle et son mari ». Le mot \zh{一起} (yìqǐ, ensemble) est essentiel : il exprime le partage des tâches.

Réponse rédigée : \zh{李老师和她的丈夫一起做家务。} (Lǐ lǎoshī hé tā de zhàngfu yìqǐ zuò jiāwù.) « Madame Li et son mari font le ménage ensemble. »

\textbf{Question 3.} Question fermée en \zh{吗}. Le texte dit : \zh{他们的女儿和儿子都上学} — la fille \textbf{et} le fils vont à l'école (\zh{都} = tous les deux). C'est un point central du thème : la fille a accès à l'éducation comme le garçon. Pour répondre « oui » à une question en \zh{吗}, on reprend le verbe de la question : \zh{上学。}

Réponse rédigée : \zh{上学。她的女儿和儿子都上学。} (Shàngxué. Tā de nǚ'ér hé érzi dōu shàngxué.) « Oui. Sa fille et son fils vont tous les deux à l'école. »

\textbf{Question 4.} On reprend la citation finale : \zh{都有学习和工作的权利}. Le verbe \zh{有} (yǒu, avoir) + \zh{权利} (quánlì, droit).

Réponse rédigée : \zh{女人和男人一样，都有学习和工作的权利。} (Nǚrén hé nánrén yíyàng, dōu yǒu xuéxí hé gōngzuò de quánlì.) « Les femmes, comme les hommes, ont le droit d'étudier et de travailler. »

Conclusion : chaque réponse reprend les mots exacts du texte, avec une phrase complète (sujet + verbe). C'est ce que le correcteur attend.
\end{exoresolu}

\begin{methode}[Reconnaître et prononcer le vocabulaire du thème]
Pour mémoriser durablement le vocabulaire de l'égalité homme-femme :
\begin{enumerate}
\item \textbf{Apprends les mots par familles de caractères.} La clé 女 (nǚ, femme) revient dans \zh{妇女} (femmes), \zh{女儿} (fille), \zh{她} (elle), \zh{妈妈} (maman). Repérer la clé aide à deviner le sens.
\item \textbf{Vérifie toujours le ton.} \zh{女} se dit nǚ (3\textsuperscript{e} ton), avec ü : les deux points restent après n. Ne confonds pas nǚ (女, femme) et nǔ (努, s'efforcer).
\item \textbf{Apprends les mots avec leur verbe habituel (collocation)} : on dit \zh{有权利} (avoir un droit), \zh{做家务} (faire le ménage), \zh{上学} (aller à l'école), \zh{教汉语} (enseigner le chinois).
\item \textbf{Associe chaque mot à son contraire ou son pendant} : \zh{男人} nánrén / \zh{女人} nǚrén ; \zh{儿子} érzi (fils) / \zh{女儿} nǚ'ér (fille) ; \zh{丈夫} zhàngfu (mari) / \zh{妻子} qīzi (épouse).
\item \textbf{Teste-toi dans les deux sens} : caractères → français, puis français → caractères + pinyin. Au Bac, les deux sens sont exigés.
\end{enumerate}
\end{methode}

\begin{exoresolu}[Vocabulaire : complète les phrases]
Énoncé. Complète chaque phrase avec le mot qui convient, choisi dans la liste : \zh{平等} (píngděng, égalité) — \zh{权利} (quánlì, droit) — \zh{家务} (jiāwù, tâches ménagères) — \zh{女儿} (nǚ'ér, fille) — \zh{丈夫} (zhàngfu, mari).

\begin{enumerate}
\item 男孩和女孩都有上学的 \_\_\_\_\_\_ 。(Les garçons et les filles ont le droit d'aller à l'école.)
\item 男人和女人应该 \_\_\_\_\_\_ 。(Les hommes et les femmes doivent être égaux.)
\item 她的 \_\_\_\_\_\_ 在银行工作。(Son mari travaille dans une banque.)
\item 周末，我们全家一起做 \_\_\_\_\_\_ 。(Le week-end, toute la famille fait le ménage ensemble.)
\item 他们的 \_\_\_\_\_\_ 今年六岁，上小学了。(Leur fille a six ans, elle va à l'école primaire.)
\end{enumerate}
\tcblower
Solution détaillée.

\textbf{1.} \zh{权利}. Structure : \zh{有} + nom = « avoir ». « Avoir le droit de faire » = \zh{有……的权利}. Phrase complète : \zh{男孩和女孩都有上学的权利。} (Nánhái hé nǚhái dōu yǒu shàngxué de quánlì.)

\textbf{2.} \zh{平等}. Ici on cherche un adjectif après \zh{应该} (yīnggāi, devoir). \zh{平等} est un adjectif : « être égal ». Phrase : \zh{男人和女人应该平等。} (Nánrén hé nǚrén yīnggāi píngděng.)

\textbf{3.} \zh{丈夫}. Le possessif \zh{她的} (tā de) appelle un nom de personne ; « mari » est le seul qui convient avec « travaille dans une banque ». Phrase : \zh{她的丈夫在银行工作。} (Tā de zhàngfu zài yínháng gōngzuò.) Remarque : le complément de lieu \zh{在银行} se place avant le verbe \zh{工作}.

\textbf{4.} \zh{家务}. Collocation obligatoire : \zh{做家务} (zuò jiāwù, faire le ménage). On ne dit jamais *\zh{做家庭}. Phrase : \zh{周末，我们全家一起做家务。} (Zhōumò, wǒmen quánjiā yìqǐ zuò jiāwù.)

\textbf{5.} \zh{女儿}. Indice : « leur fille, six ans ». \zh{女儿} = fille (de quelqu'un) ; \zh{儿子} = fils. Phrase : \zh{他们的女儿今年六岁，上小学了。} (Tāmen de nǚ'ér jīnnián liù suì, shàng xiǎoxué le.)

Conclusion : pour choisir le bon mot, appuie-toi sur la grammaire de la phrase (verbe qui précède, possessif, adjectif attendu) et sur les collocations apprises.
\end{exoresolu}

\begin{methode}[Employer les structures du thème : 一样, 都, 有……的权利]
Trois structures dominent les textes sur l'égalité homme-femme. Retiens leurs schémas.
\begin{enumerate}
\item \textbf{Comparaison d'égalité : A + 和 + B + 一样 (+ adjectif).} \zh{女人和男人一样} = « les femmes sont comme les hommes ». Pour nier : \zh{不一样} (pas pareil).
\item \textbf{都 (dōu, tous)} : se place \textbf{après le sujet pluriel et avant le verbe}. \zh{他们都工作} = « ils travaillent tous ». Jamais en début de phrase, jamais après le verbe.
\item \textbf{有 + (verbe + 的) + 权利} : « avoir le droit de faire ». \zh{有学习的权利} = « avoir le droit d'étudier ». Le petit mot \zh{的} relie le verbe au nom \zh{权利}.
\end{enumerate}
Réflexe de rédaction : pour écrire une phrase sur l'égalité, combine les trois : \zh{女儿和儿子一样，都有上学的权利。} (Nǚ'ér hé érzi yíyàng, dōu yǒu shàngxué de quánlì.) « La fille comme le fils ont tous le droit d'aller à l'école. »
\end{methode}

\begin{exoresolu}[Grammaire : traduis en chinois]
Énoncé. Traduis les phrases suivantes en chinois (caractères obligatoires, pinyin apprécié).
\begin{enumerate}
\item Les femmes et les hommes sont égaux.
\item Les filles et les garçons vont tous à l'école.
\item Elle a le droit de travailler.
\item Mon mari et moi faisons le ménage ensemble.
\end{enumerate}
\tcblower
Solution détaillée.

\textbf{1.} Structure A + 和 + B + 一样. « Femmes » = \zh{女人}, « hommes » = \zh{男人}, « égaux » = \zh{一样} suffit (on peut ajouter \zh{平等}). Réponse : \zh{女人和男人一样（平等）。} (Nǚrén hé nánrén yíyàng (píngděng).) Erreur à éviter : ne mets pas \zh{是} devant \zh{一样} ; en chinois, les adjectifs font office de verbe et n'ont pas besoin de \zh{是}.

\textbf{2.} « Tous » = \zh{都}, placé après le sujet et avant le verbe \zh{上学}. Réponse : \zh{女孩和男孩都上学。} (Nǚhái hé nánhái dōu shàngxué.) Piège : ne traduis pas « tous » par un mot en fin de phrase comme en français.

\textbf{3.} Structure 有 + verbe + 的 + 权利. Réponse : \zh{她有工作的权利。} (Tā yǒu gōngzuò de quánlì.) Le \zh{的} est obligatoire entre le verbe \zh{工作} et le nom \zh{权利}.

\textbf{4.} « Mon mari et moi » = \zh{我和丈夫} ; « ensemble » = \zh{一起}, placé avant le verbe ; « faire le ménage » = \zh{做家务}. Réponse : \zh{我和丈夫一起做家务。} (Wǒ hé zhàngfu yìqǐ zuò jiāwù.) Remarque : en chinois, on place volontiers \zh{我} en premier dans ce type d'énumération, contrairement à la politesse française (« mon mari et moi »).

Conclusion : chaque phrase applique une structure apprise ; vérifie toujours la place de \zh{都}, de \zh{一起} et de \zh{的}.
\end{exoresolu}

\begin{methode}[Répondre à des questions de compréhension et rédiger une phrase personnelle]
Au Bac, on te demande souvent de passer de la compréhension à une courte production personnelle.
\begin{enumerate}
\item \textbf{Repère le type de question} : \zh{谁} (qui), \zh{什么} (quoi), \zh{哪儿} (où), \zh{为什么} (pourquoi), \zh{吗} (oui/non).
\item \textbf{Reprends la structure de la question dans la réponse} : à \zh{她在哪儿工作？} on répond \zh{她在……工作。} en ne changeant que l'information demandée.
\item \textbf{Pour une question ouverte} (et toi ? ta famille ?), réutilise le vocabulaire du texte : \zh{在我们家，爸爸妈妈都工作，也一起做家务。} (Zài wǒmen jiā, bàba māma dōu gōngzuò, yě yìqǐ zuò jiāwù.) « Chez nous, papa et maman travaillent tous les deux et font aussi le ménage ensemble. »
\item \textbf{Pour donner ton avis} sur l'égalité, utilise : \zh{我觉得……} (wǒ juéde, je pense que) + phrase. Exemple : \zh{我觉得女人和男人应该平等。} (Wǒ juéde nǚrén hé nánrén yīnggāi píngděng.) « Je pense que les femmes et les hommes doivent être égaux. »
\item \textbf{Relis-toi} : vérifie les tons des mots-clés, la place de \zh{都}/\zh{一起}, et le classificateur devant les personnes (\zh{一个女儿}, \zh{一位老师}).
\end{enumerate}
\end{methode}

\begin{exoresolu}[Production guidée : et dans ta famille ?]
Énoncé. Réponds en chinois, par deux ou trois phrases complètes, à la question : 在你们家，谁做家务？你觉得男人和女人应该平等吗？(Zài nǐmen jiā, shéi zuò jiāwù ? Nǐ juéde nánrén hé nǚrén yīnggāi píngděng ma ?) (Dans ta famille, qui fait le ménage ? Penses-tu que les hommes et les femmes doivent être égaux ?)
\tcblower
Solution détaillée.

\textbf{Étape 1 — Analyse de la question.} Deux questions : (a) \zh{谁做家务} demande une personne ; (b) \zh{你觉得……吗} demande un avis personnel, avec \zh{应该} (devoir) + adjectif \zh{平等}.

\textbf{Étape 2 — Choix du vocabulaire.} Famille : \zh{爸爸} (papa), \zh{妈妈} (maman), \zh{我} (je/moi). Verbes : \zh{做家务}, \zh{工作}, \zh{一起}. Opinion : \zh{我觉得} + \zh{应该平等} ; on peut enrichir avec \zh{都有学习和工作的权利}.

\textbf{Étape 3 — Rédaction.} Modèle de réponse :

\zh{在我们家，爸爸和妈妈一起做家务。妈妈做饭，爸爸洗碗，我也帮忙。我觉得女人和男人应该平等，因为大家都有学习和工作的权利。}

(Zài wǒmen jiā, bàba hé māma yìqǐ zuò jiāwù. Māma zuòfàn, bàba xǐ wǎn, wǒ yě bāngmáng. Wǒ juéde nǚrén hé nánrén yīnggāi píngděng, yīnwèi dàjiā dōu yǒu xuéxí hé gōngzuò de quánlì.)

« Chez nous, papa et maman font le ménage ensemble. Maman cuisine, papa lave la vaisselle, et moi j'aide aussi. Je pense que les femmes et les hommes doivent être égaux, parce que tout le monde a le droit d'étudier et de travailler. »

\textbf{Étape 4 — Vérification.} Le complément de lieu \zh{在我们家} est bien en tête de phrase ; \zh{一起} et \zh{也} sont bien avant les verbes ; \zh{因为} (yīnwèi, parce que) introduit correctement la justification ; la collocation \zh{有……的权利} est respectée avec \zh{的}.

Conclusion : une bonne réponse personnelle = vocabulaire du texte + structure d'opinion \zh{我觉得} + une justification avec \zh{因为}.
\end{exoresolu}

\begin{methode}[Écrire correctement les caractères du thème : clés et ordre des traits]
Pour ne pas perdre de points en écriture :
\begin{enumerate}
\item \textbf{Identifie la clé (radical)} avant d'écrire : 女 (femme) dans \zh{好}, \zh{她}, \zh{妇} ; 子 (enfant) dans \zh{孩}, \zh{学} ; 宀 (toit) dans \zh{家}, \zh{室}.
\item \textbf{Respecte l'ordre général des traits} : de haut en bas, de gauche à droite, l'horizontal avant le vertical, l'extérieur avant l'intérieur, et on ferme la boîte en dernier.
\item \textbf{Cas de 女 (3 traits)} : d'abord le trait brisé (descendant oblique puis horizontal-crochet), ensuite le trait oblique tombant à gauche, enfin l'horizontal. Attention : à gauche d'un caractère composé (comme dans \zh{好}), 女 s'écrit plus étroit et son dernier trait horizontal devient un trait montant.
\item \textbf{Cas de 家 (10 traits)} : d'abord le toit 宀 (point, trait vertical gauche, crochet horizontal), puis le cochon 豕 en dessous. Sens étymologique : « un cochon sous le toit » = la maison, la famille.
\item \textbf{Entraîne-toi par paires} de caractères proches : \zh{女} nǚ / \zh{母} mǔ ; \zh{子} zi / \zh{了} le ; \zh{她} tā / \zh{他} tā (même prononciation, clés différentes : 女 pour « elle », 亻 pour « il »).
\end{enumerate}
\end{methode}

\begin{exoresolu}[Écriture : distingue les caractères]
Énoncé. Pour chaque paire, donne le pinyin, la traduction et la clé de chaque caractère, puis écris une phrase avec le premier élément : (a) 她 / 他 ; (b) 女儿 / 儿子.
\tcblower
Solution détaillée.

\textbf{(a) 她 / 他.}
\zh{她} : tā, « elle » ; clé 女 (femme) à gauche + 也 (yě, aussi) à droite qui donne le son.
\zh{他} : tā, « il » ; clé 亻 (homme, variante de 人) à gauche + 也 à droite.
Même prononciation tā, sens différent selon la clé : c'est la clé qui porte le sens.
Phrase avec \zh{她} : \zh{她是我的老师。} (Tā shì wǒ de lǎoshī.) « Elle est mon professeur. »

\textbf{(b) 女儿 / 儿子.}
\zh{女儿} : nǚ'ér, « fille » ; composé de 女 (femme) + 儿 (enfant). Attention à l'apostrophe dans le pinyin nǚ'ér : elle sépare les deux syllabes pour éviter la lecture fautive « nüér ».
\zh{儿子} : érzi, « fils » ; 儿 + 子 (enfant). Le 子 final se prononce au ton neutre (zi).
Phrase avec \zh{女儿} : \zh{他们的女儿今年六岁。} (Tāmen de nǚ'ér jīnnián liù suì.) « Leur fille a six ans cette année. »

Conclusion : la clé 女 signale le féminin ; la clé 亻 signale une personne (souvent masculine ou neutre). Apprendre les caractères par clés évite les confusions à l'écrit.
\end{exoresolu}

\begin{attention}[Les 5 erreurs qui coûtent des points au Bac]
\begin{enumerate}
\item \textbf{Oublier le ton ou le ü de 女 (nǚ).} Écrire « nu » au lieu de nǚ fausse le mot entier. Après n et l, les deux points du ü restent : nǚ, lǜ. Après j, q, x, ils disparaissent mais le son reste ü (qù, xué).
\item \textbf{Confondre 她 (elle) et 他 (il).} Même son tā, clés différentes. Au Bac, écrire \zh{他} pour parler d'une femme est une faute de caractère. Retiens : clé 女 = féminin.
\item \textbf{Mal placer 都 et 一起.} Ces deux mots se placent \textbf{après le sujet et avant le verbe} : \zh{他们都上学}, \zh{我们一起做家务}. Le calque du français « ils vont tous à l'école » → *\zh{他们上学都} est faux.
\item \textbf{Oublier le 的 dans 有……的权利.} On écrit \zh{有学习的权利}, jamais *\zh{有学习权利}. Le \zh{的} relie le verbe au nom ; son oubli est sanctionné comme faute de structure.
\item \textbf{Négliger le classificateur.} « Une enseignante » = \zh{一位老师} (classificateur respectueux \zh{位}), « une fille » = \zh{一个女儿}. Utiliser \zh{个} partout, ou aucun classificateur, fait perdre des points en expression écrite.
\end{enumerate}
\end{attention}

\newpage

\coursec{Exercices}

\courssub{Je vérifie mes connaissances}

\begin{exercice}[QCM : le sens des mots]
Entoure la bonne réponse.
\begin{enumerate}
\item \zh{平等} (píngděng) signifie : a) la liberté \quad b) l'égalité \quad c) la paix
\item \zh{重男轻女} (zhòng nán qīng nǚ) signifie : a) respecter les anciens \quad b) valoriser les garçons et mépriser les filles \quad c) aimer toute la famille
\item \zh{地位} (dìwèi) signifie : a) la place, le statut \quad b) la terre \quad c) le travail
\item \zh{参加} (cānjiā) signifie : a) quitter \quad b) regarder \quad c) participer à
\end{enumerate}
\end{exercice}

\begin{exercice}[Vrai ou faux ?]
Réponds par vrai (V) ou faux (F) et justifie en une phrase en français.
\begin{enumerate}
\item \zh{妇女} (fùnǚ) désigne les enfants de la famille.
\item \zh{权利} (quánlì) signifie « les droits ».
\item \zh{照顾} (zhàogù) signifie « s'occuper de ».
\item \zh{越来越} (yuè lái yuè) exprime une diminution progressive.
\end{enumerate}
\end{exercice}

\begin{exercice}[Vocabulaire : nature des mots]
Classe les douze mots suivants dans un tableau à trois colonnes (nom / verbe / adjectif) et donne la traduction française de chacun : \zh{女人}、\zh{平等}、\zh{工作}、\zh{社会}、\zh{照顾}、\zh{重要}、\zh{参加}、\zh{家庭}、\zh{进步}、\zh{思想}、\zh{尊重}、\zh{幸福}。
\end{exercice}

\begin{exercice}[Associe caractères et pinyin]
Relie chaque caractère à son pinyin, puis donne sa traduction.
\begin{center}
\begin{tabularx}{\linewidth}{|X|X|}\hline
\rowcolor{popblueL} Caractères & Pinyin \\ \hline
1. 女 & a) jiā \\ \hline
2. 好 & b) nǚ \\ \hline
3. 男 & c) qī \\ \hline
4. 家 & d) hǎo \\ \hline
5. 妻 & e) nán \\ \hline
\end{tabularx}
\end{center}
\end{exercice}

\courssub{Je m'entraîne}

\begin{exercice}[Traduction : version]
Traduis en français les phrases suivantes.
\begin{enumerate}
\item \zh{在家里，丈夫和妻子一起照顾孩子。但是，重男轻女的思想还存在。}\\
\emph{Zài jiā li, zhàngfu hé qīzi yìqǐ zhàogù háizi. Dànshì, zhòng nán qīng nǚ de sīxiǎng hái cúnzài.}
\item \zh{现在，越来越多的女人上大学、做工作。}\\
\emph{Xiànzài, yuè lái yuè duō de nǚrén shàng dàxué, zuò gōngzuò.}
\end{enumerate}
\end{exercice}

\begin{exercice}[Traduction : thème]
Traduis en chinois (caractères et pinyin).
\begin{enumerate}
\item Les femmes peuvent étudier, travailler et participer aux activités sociales.
\item À la maison, le mari et la femme sont égaux.
\item La situation des femmes camerounaises change de plus en plus vite.
\end{enumerate}
\end{exercice}

\begin{exercice}[Texte à trous]
Complète le texte avec les mots de la liste : \zh{平等}、\zh{照顾}、\zh{参加}、\zh{地位}、\zh{越来越}。

\zh{今天，女人的社会（1）提高了。她们可以（2）社会活动，也可以工作。在家里，丈夫和妻子（3）孩子。男人和女人是（4）的。女人的生活（5）好。}
\end{exercice}

\begin{exercice}[Remets les mots dans l'ordre]
Remets les mots dans l'ordre pour faire une phrase correcte, puis traduis-la.
\begin{enumerate}
\item \zh{女人 / 越来越 / 的 / 地位 / 提高 / 了 / 社会}
\item \zh{一起 / 丈夫 / 妻子 / 和 / 家务 / 做}
\item \zh{重男轻女 / 思想 / 的 / 还 / 存在}
\end{enumerate}
\end{exercice}

\courssub{Je me prépare au Bac}

\begin{exercice}[Compréhension écrite — type Bac]
Lis le texte suivant, puis réponds aux questions en français.

\begin{textebox}[Une femme médecin chinoise]
\zh{王芳是北京的一位医生。她今年四十岁，有一个幸福的家庭。}\\
\emph{Wáng Fāng shì Běijīng de yí wèi yīshēng. Tā jīnnián sìshí suì, yǒu yí gè xìngfú de jiātíng.}\\
« Wang Fang est médecin à Beijing. Elle a quarante ans et une famille heureuse. »\\[0.3em]
\zh{二十年前，她上大学学习医学。那时候，学医学的女人不太多。}\\
\emph{Èrshí nián qián, tā shàng dàxué xuéxí yīxué. Nà shíhou, xué yīxué de nǚrén bú tài duō.}\\
« Il y a vingt ans, elle est entrée à l'université pour étudier la médecine. À cette époque, les femmes qui étudiaient la médecine n'étaient pas très nombreuses. »\\[0.3em]
\zh{现在，她在医院工作，也参加很多社会活动。}\\
\emph{Xiànzài, tā zài yīyuàn gōngzuò, yě cānjiā hěn duō shèhuì huódòng.}\\
« Maintenant, elle travaille à l'hôpital et participe aussi à de nombreuses activités sociales. »\\[0.3em]
\zh{在家里，她和丈夫一起照顾孩子、做家务。}\\
\emph{Zài jiā li, tā hé zhàngfu yìqǐ zhàogù háizi, zuò jiāwù.}\\
« À la maison, elle s'occupe des enfants et fait le ménage avec son mari. »\\[0.3em]
\zh{她说：男人和女人是平等的，女人的地位越来越高了。}\\
\emph{Tā shuō: nánrén hé nǚrén shì píngděng de, nǚrén de dìwèi yuè lái yuè gāo le.}\\
« Elle dit : les hommes et les femmes sont égaux, et le statut des femmes est de plus en plus élevé. »
\end{textebox}

Questions :
\begin{enumerate}
\item Vrai ou faux ? Justifie avec une phrase du texte : a) Wang Fang est institutrice. b) Il y a vingt ans, beaucoup de femmes étudiaient la médecine. c) Le mari de Wang Fang l'aide à la maison.
\item Quel est le métier de Wang Fang et où travaille-t-elle ?
\item Que fait Wang Fang en dehors de son travail à l'hôpital ?
\item Question de langue : relève dans le texte la structure \zh{越来越} et explique son sens.
\item Question de langue : traduis en français la dernière phrase du texte.
\end{enumerate}
\end{exercice}

\begin{exercice}[Thème et version — type Bac]
\textbf{A. Version.} Traduis en français le texte suivant.\\
\zh{在中国和喀麦隆，女人的生活越来越好了。她们可以学习、工作，也可以参加社会活动。但是，在一些地方，重男轻女的思想还存在。我们应该尊重女人，因为男人和女人是平等的。}\\[0.3em]
\textbf{B. Thème.} Traduis en chinois (caractères et pinyin) :
\begin{enumerate}
\item Les femmes peuvent étudier, travailler et participer aux activités sociales.
\item Le mari et la femme s'occupent ensemble des enfants.
\item La pensée qui valorise les garçons et méprise les filles existe encore.
\end{enumerate}
\end{exercice}

\begin{exercice}[Expression écrite guidée — type Bac]
Sujet : \textbf{« La femme camerounaise aujourd'hui »} (\zh{今天的喀麦隆女人}).

Rédige un texte de \textbf{5 à 7 phrases en chinois} (environ 60 à 80 caractères). Consignes :
\begin{enumerate}
\item utilise au moins \textbf{5 mots du glossaire} de la leçon (par exemple : \zh{平等}、\zh{地位}、\zh{参加}、\zh{照顾}、\zh{尊重}、\zh{社会}…) ;
\item emploie au moins une fois la structure \zh{越来越} ;
\item présente la situation actuelle des femmes au Cameroun (études, travail, famille) et donne ton avis en une phrase finale avec \zh{我觉得}…
\end{enumerate}
N'oublie pas la ponctuation chinoise ：，。？！
\end{exercice}

\newpage

\coursec{Corrigés des exercices}

\begin{corrige}[Exercice 1 — QCM : le sens des mots]
\begin{enumerate}
\item \textbf{b) l'égalité}. \zh{平等} (píngděng) : 平 = « plat, égal », 等 = « égal, rang ». Exemple : \zh{男人和女人是平等的。} \emph{Nánrén hé nǚrén shì píngděng de.} « Les hommes et les femmes sont égaux. »
\item \textbf{b) valoriser les garçons et mépriser les filles}. \zh{重} (zhòng) = « accorder de l'importance », \zh{轻} (qīng) = « prendre à la légère, mépriser ». C'est une expression figée en quatre caractères.
\item \textbf{a) la place, le statut}. \zh{地位} (dìwèi) : \zh{社会地位} « le statut social ». Attention : « la terre » se dit \zh{土地} (tǔdì).
\item \textbf{c) participer à}. \zh{参加} (cānjiā) + nom d'activité : \zh{参加社会活动} « participer aux activités sociales ».
\end{enumerate}
\end{corrige}

\begin{corrige}[Exercice 2 — Vrai ou faux ?]
\begin{enumerate}
\item \textbf{F}. \zh{妇女} (fùnǚ) désigne les femmes (adultes), pas les enfants ; « les enfants » se dit \zh{孩子} (háizi).
\item \textbf{V}. \zh{权利} (quánlì) signifie bien « les droits » : \zh{女人的权利} « les droits des femmes ».
\item \textbf{V}. \zh{照顾} (zhàogù) signifie « s'occuper de, prendre soin de » : \zh{照顾孩子} « s'occuper des enfants ».
\item \textbf{F}. \zh{越来越} + adjectif exprime une \emph{évolution progressive}, le plus souvent une augmentation : \zh{越来越好} « de mieux en mieux ». La diminution ou l'augmentation dépend de l'adjectif qui suit, pas de la structure elle-même.
\end{enumerate}
\end{corrige}

\begin{corrige}[Exercice 3 — Vocabulaire : nature des mots]
\begin{center}
\begin{tabularx}{\linewidth}{|X|X|X|}\hline
\rowcolor{popblueL} Nom & Verbe & Adjectif \\ \hline
\zh{女人} nǚrén — la femme & \zh{工作} gōngzuò — travailler (aussi nom : le travail) & \zh{平等} píngděng — égal (aussi nom : l'égalité) \\ \hline
\zh{社会} shèhuì — la société & \zh{照顾} zhàogù — s'occuper de & \zh{重要} zhòngyào — important \\ \hline
\zh{家庭} jiātíng — la famille & \zh{参加} cānjiā — participer à & \zh{进步} jìnbù — progressif (aussi nom/verbe : progrès, progresser) \\ \hline
\zh{思想} sīxiǎng — la pensée, l'idéologie & \zh{尊重} zūnzhòng — respecter & \zh{幸福} xìngfú — heureux \\ \hline
\end{tabularx}
\end{center}
\textbf{Remarque méthodologique :} en chinois, beaucoup de mots ont une double fonction (nom \emph{et} verbe, ou nom \emph{et} adjectif) : c'est le contexte de la phrase qui décide. Ainsi \zh{工作} peut être un verbe (\zh{她在医院工作。} « Elle travaille à l'hôpital. ») ou un nom (\zh{她的工作很好。} « Son travail est bien. »).
\end{corrige}

\begin{corrige}[Exercice 4 — Associe caractères et pinyin]
\begin{enumerate}
\item \zh{女} — \textbf{b) nǚ} : « femme, féminin ». Clé : 女 (la femme), 3 traits.
\item \zh{好} — \textbf{d) hǎo} : « bon, bien ». Composé de 女 (femme) + 子 (enfant) : une femme avec son enfant, image du « bien ».
\item \zh{男} — \textbf{e) nán} : « homme, masculin ». Composé de 田 (champ) + 力 (force) : la force au champ.
\item \zh{家} — \textbf{a) jiā} : « famille, maison ». Clé : 宀 (toit) ; sous le toit, 豕 (le cochon), signe de richesse domestique.
\item \zh{妻} — \textbf{c) qī} : « épouse ». Attention à ne pas confondre avec \zh{妾} ni avec la prononciation de \zh{七} (qī, « sept »), homophone au ton près identique : le contexte tranche.
\end{enumerate}
\end{corrige}

\begin{corrige}[Exercice 5 — Traduction : version]
\begin{enumerate}
\item \zh{在家里，丈夫和妻子一起照顾孩子。但是，重男轻女的思想还存在。}\\
« À la maison, le mari et la femme s'occupent ensemble des enfants. Mais la pensée qui valorise les garçons et méprise les filles existe encore. »\\
\textbf{Points de langue :} \zh{一起} (yìqǐ, « ensemble ») se place \emph{avant} le verbe ; \zh{还} (hái) signifie ici « encore » (et non « aussi ») ; \zh{存在} (cúnzài) = « exister ».
\item \zh{现在，越来越多的女人上大学、做工作。}\\
« Maintenant, de plus en plus de femmes entrent à l'université et travaillent. »\\
\textbf{Points de langue :} \zh{越来越} + adjectif (ici \zh{多} « nombreux ») = « de plus en plus » ; \zh{上大学} = « faire des études universitaires » (litt. « monter à l'université ») ; la virgule chinoise \zh{、} énumère deux actions coordonnées.
\end{enumerate}
\end{corrige}

\begin{corrige}[Exercice 6 — Traduction : thème]
\begin{enumerate}
\item \zh{女人可以学习、工作，也可以参加社会活动。}\\
\emph{Nǚrén kěyǐ xuéxí, gōngzuò, yě kěyǐ cānjiā shèhuì huódòng.}\\
\textbf{Justification :} le modal \zh{可以} (« pouvoir », permission) précède le verbe ; \zh{也} (« aussi ») se place avant \zh{可以} dans la deuxième proposition.
\item \zh{在家里，丈夫和妻子是平等的。}\\
\emph{Zài jiā li, zhàngfu hé qīzi shì píngděng de.}\\
\textbf{Justification :} adjectif employé en attribut après \zh{是}…\zh{的} : structure Sujet + 是 + adjectif + 的. On peut aussi dire \zh{丈夫和妻子平等。} sans \zh{是}…\zh{的}, mais la structure avec \zh{是}…\zh{的} est plus sûre au Bac.
\item \zh{喀麦隆女人的情况变化得越来越快。} (ou plus simplement : \zh{喀麦隆女人的生活变化越来越快。})\\
\emph{Kāmàilóng nǚrén de qíngkuàng biànhuà de yuè lái yuè kuài.}\\
\textbf{Justification :} « la situation » = \zh{情况} (qíngkuàng) ; avec un verbe d'action comme \zh{变化}, on utilise le complément de degré introduit par \zh{得} : Verbe + 得 + 越来越 + adjectif. \textbf{Piège :} ne pas écrire \zh{得} à la place de \zh{的} : \zh{的} suit le nom (\zh{喀麦隆女人的情况}), \zh{得} suit le verbe.
\end{enumerate}
\end{corrige}

\begin{corrige}[Exercice 7 — Texte à trous]
Texte complété :\\
\zh{今天，女人的社会（1）\textbf{地位}提高了。她们可以（2）\textbf{参加}社会活动，也可以工作。在家里，丈夫和妻子（3）\textbf{照顾}孩子。男人和女人是（4）\textbf{平等}的。女人的生活（5）\textbf{越来越}好。}\\[0.3em]
Traduction : « Aujourd'hui, le statut social des femmes s'est élevé. Elles peuvent participer aux activités sociales et aussi travailler. À la maison, le mari et la femme s'occupent des enfants. Les hommes et les femmes sont égaux. La vie des femmes est de mieux en mieux. »\\[0.3em]
\textbf{Justifications :}
\begin{enumerate}
\item \zh{社会地位} : collocation figée « statut social » ; \zh{提高} (tígāo, « s'élever ») exige un nom comme \zh{地位} ou \zh{水平}.
\item \zh{参加} + activité : seul verbe de la liste qui se construit avec \zh{活动}.
\item \zh{照顾孩子} : collocation « s'occuper des enfants ».
\item \zh{是平等的} : structure 是 + adjectif + 的 ; \zh{平等} est le seul adjectif de la liste.
\item \zh{越来越好} : structure « de mieux en mieux » ; \zh{越来越} se place immédiatement avant l'adjectif.
\end{enumerate}
\end{corrige}

\begin{corrige}[Exercice 8 — Remets les mots dans l'ordre]
\begin{enumerate}
\item \zh{女人的社会地位越来越提高了。}\\
\emph{Nǚrén de shèhuì dìwèi yuè lái yuè tígāo le.}\\
« Le statut social des femmes s'est élevé de plus en plus. »\\
\textbf{Ordre :} sujet (\zh{女人的社会地位}) + \zh{越来越} + verbe + \zh{了} final (changement acquis).
\item \zh{丈夫和妻子一起做家务。}\\
\emph{Zhàngfu hé qīzi yìqǐ zuò jiāwù.}\\
« Le mari et la femme font le ménage ensemble. »\\
\textbf{Ordre :} sujet composé avec \zh{和} + \zh{一起} (avant le verbe) + verbe + complément d'objet.
\item \zh{重男轻女的思想还存在。}\\
\emph{Zhòng nán qīng nǚ de sīxiǎng hái cúnzài.}\\
« La pensée qui valorise les garçons et méprise les filles existe encore. »\\
\textbf{Ordre :} sujet (\zh{重男轻女的思想}, avec \zh{的} de subordination) + adverbe \zh{还} + verbe \zh{存在}.
\end{enumerate}
\textbf{Rappel :} en chinois, les adverbes (\zh{一起}、\zh{还}、\zh{也}) se placent toujours \emph{avant} le verbe, jamais en fin de phrase comme en français.
\end{corrige}

\begin{corrige}[Exercice 9 — Compréhension écrite — type Bac]
\textbf{Barème indicatif (sur 10 points) :} question 1 : 3 pts (1 pt par item, justification comprise) ; question 2 : 2 pts ; question 3 : 2 pts ; question 4 : 1,5 pt ; question 5 : 1,5 pt.

\begin{enumerate}
\item a) \textbf{Faux}. Le texte dit : \zh{王芳是北京的一位医生。} « Wang Fang est médecin à Beijing », et non institutrice.\\
b) \textbf{Faux}. Le texte dit : \zh{那时候，学医学的女人不太多。} « À cette époque, les femmes qui étudiaient la médecine n'étaient pas très nombreuses. »\\
c) \textbf{Vrai}. Le texte dit : \zh{在家里，她和丈夫一起照顾孩子、做家务。} « À la maison, elle s'occupe des enfants et fait le ménage avec son mari » : le mari participe donc aux tâches.
\item Wang Fang est \textbf{médecin} (\zh{医生}) et elle travaille \textbf{à l'hôpital} (\zh{她在医院工作。}), à Beijing.
\item En dehors de l'hôpital, elle \textbf{participe à de nombreuses activités sociales} (\zh{参加很多社会活动}) et, à la maison, elle \textbf{s'occupe des enfants et fait le ménage avec son mari}.
\item La structure relevée est \zh{女人的地位越来越高了} (\emph{nǚrén de dìwèi yuè lái yuè gāo le}). \zh{越来越} + adjectif signifie « de plus en plus… » : ici, « le statut des femmes est de plus en plus élevé ». Le \zh{了} final marque le changement de situation.
\item \zh{她说：男人和女人是平等的，女人的地位越来越高了。}\\
« Elle dit : les hommes et les femmes sont égaux, et le statut des femmes est de plus en plus élevé. »
\end{enumerate}
\textbf{Points de vigilance :} ne pas traduire \zh{不太多} par « pas du tout » mais par « pas très nombreuses » ; \zh{一位医生} : le classificateur des personnes respectées est \zh{位} (wèi), plus poli que \zh{个}.
\end{corrige}

\begin{corrige}[Exercice 10 — Thème et version — type Bac]
\textbf{Barème indicatif (sur 10 points) :} version 4 pts ; thème 6 pts (2 pts par phrase : 1 pt exactitude des mots, 1 pt ordre des mots et mots-outils).

\textbf{A. Version.}\\
\zh{在中国和喀麦隆，女人的生活越来越好了。她们可以学习、工作，也可以参加社会活动。但是，在一些地方，重男轻女的思想还存在。我们应该尊重女人，因为男人和女人是平等的。}\\
« En Chine et au Cameroun, la vie des femmes est de mieux en mieux. Elles peuvent étudier, travailler et aussi participer aux activités sociales. Mais dans certains endroits, la pensée qui valorise les garçons et méprise les filles existe encore. Nous devons respecter les femmes, parce que les hommes et les femmes sont égaux. »\\
\textbf{Points de vigilance :} \zh{一些} = « quelques, certains » ; \zh{应该} (yīnggāi) = « devoir » (conseil moral) ; \zh{因为} = « parce que ».

\textbf{B. Thème.}
\begin{enumerate}
\item \zh{女人可以学习、工作，也可以参加社会活动。}\\
\emph{Nǚrén kěyǐ xuéxí, gōngzuò, yě kěyǐ cānjiā shèhuì huódòng.}
\item \zh{丈夫和妻子一起照顾孩子。}\\
\emph{Zhàngfu hé qīzi yìqǐ zhàogù háizi.}\\
(« ensemble » = \zh{一起}, placé avant le verbe.)
\item \zh{重男轻女的思想还存在。}\\
\emph{Zhòng nán qīng nǚ de sīxiǎng hái cúnzài.}\\
(« existe encore » = \zh{还存在} : adverbe \zh{还} avant le verbe.)
\end{enumerate}
\textbf{Points de vigilance :} ne pas oublier le \zh{的} entre \zh{重男轻女} et \zh{思想} ; ne pas placer \zh{一起} ou \zh{还} après le verbe ; ponctuation chinoise obligatoire (\zh{、} pour l'énumération, \zh{。} en fin de phrase).
\end{corrige}

\begin{corrige}[Exercice 11 — Expression écrite guidée — type Bac]
\textbf{Barème indicatif (sur 10 points) :} respect des consignes (5 mots du glossaire, structure \zh{越来越}, avis final avec \zh{我觉得}) : 3 pts ; correction grammaticale (ordre des mots, \zh{的}/\zh{了}) : 3 pts ; vocabulaire et caractères : 2 pts ; ponctuation chinoise et présentation : 2 pts.

\textbf{Proposition de rédaction modèle :}\\[0.3em]
\zh{今天的喀麦隆女人}\\
\zh{现在，喀麦隆女人的生活越来越好了。很多女人上大学，也有好工作。她们参加社会活动，女人的社会地位提高了。在家里，丈夫和妻子一起照顾孩子、做家务。但是，在一些地方，重男轻女的思想还存在。我觉得男人和女人是平等的，我们应该尊重女人。}\\[0.3em]
\emph{Xiànzài, Kāmàilóng nǚrén de shēnghuó yuè lái yuè hǎo le. Hěn duō nǚrén shàng dàxué, yě yǒu hǎo gōngzuò. Tāmen cānjiā shèhuì huódòng, nǚrén de shèhuì dìwèi tígāo le. Zài jiā li, zhàngfu hé qīzi yìqǐ zhàogù háizi, zuò jiāwù. Dànshì, zài yìxiē dìfang, zhòng nán qīng nǚ de sīxiǎng hái cúnzài. Wǒ juéde nánrén hé nǚrén shì píngděng de, wǒmen yīnggāi zūnzhòng nǚrén.}\\[0.3em]
Traduction : « Aujourd'hui, la vie des femmes camerounaises est de mieux en mieux. Beaucoup de femmes font des études universitaires et ont de bons emplois. Elles participent aux activités sociales et leur statut social s'est élevé. À la maison, le mari et la femme s'occupent ensemble des enfants et font le ménage. Mais dans certains endroits, la pensée qui valorise les garçons et méprise les filles existe encore. Je pense que les hommes et les femmes sont égaux et que nous devons respecter les femmes. »\\[0.3em]
\textbf{Vérification des consignes :} 7 phrases ; mots du glossaire utilisés : \zh{社会}、\zh{地位}、\zh{参加}、\zh{照顾}、\zh{平等}、\zh{尊重} (6 mots) ; structure \zh{越来越好} employée ; avis final avec \zh{我觉得}.

\textbf{Points de vigilance :}
\begin{itemize}
\item \zh{越来越} se place juste avant l'adjectif, avec \zh{了} final pour marquer le changement : \zh{越来越好了}.
\item \zh{我觉得} + phrase complète, sans \zh{是} après \zh{觉得}.
\item Ne pas écrire \zh{女人的地位提高了} sans \zh{了} : ici \zh{了} marque le résultat acquis.
\item Ponctuation : \zh{、} entre deux verbes énumérés, \zh{，} entre les propositions, \zh{。} en fin de phrase.
\end{itemize}
\end{corrige}

\newpage

\coursec{Fiche bilan}

\begin{popfigure}
\begin{tikzpicture}[mindmap, grow cyclic, every node/.style={concept, text=white, align=center, font=\small}, root concept/.append style={concept color=popink, font=\bfseries}, level 1/.append style={level distance=3.4cm, sibling angle=51.4}, level 2/.append style={level distance=2.2cm, sibling angle=45, font=\scriptsize}]
\node[root concept, align=center]{妇女与平等\\ \emph{Femmes et égalité}}
  child[concept color=popblue]{ node[align=center]{家庭\\ jiātíng\\ Famille} child{ node[align=center]{丈夫 zhàngfu\\ 妻子 qīzi} } child{ node[align=center]{家务 jiāwù\\ 孩子 háizi} } }
  child[concept color=poporange]{ node[align=center]{工作\\ gōngzuò\\ Travail} child{ node[align=center]{职业 zhíyè\\ 工资 gōngzī} } }
  child[concept color=popgreen]{ node[align=center]{教育\\ jiàoyù\\ Éducation} child{ node[align=center]{上学 shàngxué\\ 大学 dàxué} } }
  child[concept color=poppurple]{ node[align=center]{权利\\ quánlì\\ Droits} child{ node[align=center]{平等 píngděng\\ 机会 jīhuì} } }
  child[concept color=poppink]{ node[align=center]{社会\\ shèhuì\\ Société} child{ node[align=center]{地位 dìwèi\\ 发展 fāzhǎn} } }
  child[concept color=popgold]{ node[align=center]{观念\\ guānniàn\\ Mentalités} child{ node[align=center]{以前 yǐqián\\ 现在 xiànzài} } }
  child[concept color=popblue!70]{ node[align=center]{结构\\ Structures} child{ node[align=center]{越来越…\\ 不但…而且…} } };
\end{tikzpicture}
\legende{Carte mentale de la leçon : la situation de la femme et l'égalité homme-femme}
\end{popfigure}

\begin{bilan}
\textbf{1. Lexique clé à connaître par cœur}
\begin{itemize}
  \item \zh{妇女} fùnǚ : la femme (registre soutenu) ; \zh{男女平等} nánnǚ píngděng : égalité homme-femme
  \item \zh{地位} dìwèi : statut, position ; \zh{权利} quánlì : droit ; \zh{机会} jīhuì : occasion, chance
  \item \zh{工作} gōngzuò : travailler / travail ; \zh{职业} zhíyè : métier, profession ; \zh{工资} gōngzī : salaire
  \item \zh{家务} jiāwù : tâches ménagères ; \zh{照顾} zhàogù : s'occuper de ; \zh{教育} jiàoyù : éducation
  \item \zh{社会} shèhuì : société ; \zh{发展} fāzhǎn : développement ; \zh{改变} gǎibiàn : changer
  \item \zh{以前} yǐqián : avant, autrefois ; \zh{现在} xiànzài : maintenant ; \zh{越来越} yuè lái yuè : de plus en plus
\end{itemize}

\textbf{2. Règles de grammaire à connaître par cœur}

\begin{center}
\begin{tabularx}{\linewidth}{|X|X|X|}
\hline
\rowcolor{popblueL} Structure & Exemple (caractères + pinyin) & Traduction \\
\hline
越来越 + adjectif/verbe & 女性的地位越来越高。Nǚxìng de dìwèi yuè lái yuè gāo. & Le statut des femmes est de plus en plus élevé. \\
\hline
不但 A，而且 B & 她不但工作，而且照顾家庭。Tā búdàn gōngzuò, érqiě zhàogù jiātíng. & Non seulement elle travaille, mais elle s'occupe aussi de la famille. \\
\hline
和 / 跟 A 一样 + adj. & 女人和男人一样重要。Nǚrén hé nánrén yíyàng zhòngyào. & Les femmes sont aussi importantes que les hommes. \\
\hline
Sujet + 在 + lieu + verbe & 她在学校学习。Tā zài xuéxiào xuéxí. & Elle étudie à l'école. \\
\hline
时间 + 以前 / 以后 & 三十年以前 sānshí nián yǐqián & il y a trente ans \\
\hline
\end{tabularx}
\end{center}

\textbf{3. Méthodes express}
\begin{itemize}
  \item \emph{Lire un texte} : 1) lire le titre et deviner le thème ; 2) repérer les mots connus (\zh{妇女}, \zh{平等}, \zh{工作}) ; 3) découper en phrases (，。) ; 4) répondre aux questions en recopiant la structure de la phrase du texte.
  \item \emph{Mémoriser le vocabulaire} : apprendre par familles de mots (\zh{工作} → \zh{工人} → \zh{工资}) et par paires d'opposés (\zh{以前}/\zh{现在}, \zh{男}/\zh{女}).
  \item \emph{Écrire les caractères} : identifier d'abord la clé (女 pour les caractères liés aux femmes : \zh{妇}, \zh{妈}, \zh{姐}, \zh{妹}), puis écrire de haut en bas, de gauche à droite, l'horizontal avant le vertical.
  \item \emph{Répondre à une question de compréhension} : reprendre les mots de la question, répondre par une phrase complète sujet + verbe, jamais par un seul mot.
\end{itemize}
\end{bilan}

\begin{aretenir}[Checklist avant le Bac]
\begin{itemize}
  \item[$\square$] Je sais lire et prononcer (avec les bons tons) tout le vocabulaire du thème : \zh{妇女}, \zh{平等}, \zh{地位}, \zh{权利}, \zh{机会}, \zh{家务}, \zh{职业}.
  \item[$\square$] Je connais la clé 女 (femme) et je sais écrire \zh{妇}, \zh{妈}, \zh{姐}, \zh{妹} dans le bon ordre des traits.
  \item[$\square$] Je sais employer la structure \zh{越来越} + adjectif dans une phrase complète.
  \item[$\square$] Je sais construire une phrase avec \zh{不但…而且…} sans changer l'ordre des mots.
  \item[$\square$] Je sais comparer avec \zh{和…一样} : A + 和 + B + 一样 + adjectif.
  \item[$\square$] Je place correctement le complément de lieu avec \zh{在} avant le verbe.
  \item[$\square$] Je distingue \zh{以前} (avant) et \zh{以后} (après) placés après l'expression de temps.
  \item[$\square$] Je sais répondre aux questions de compréhension par des phrases complètes en chinois.
  \item[$\square$] Je peux résumer le texte support en trois phrases simples avec mes propres mots.
  \item[$\square$] Je peux comparer oralement la situation des femmes au Cameroun et en Chine avec deux exemples concrets.
  \item[$\square$] Je fais attention aux tons : \zh{mā} (妈, maman) n'est pas \zh{mǎ} (马, cheval).
  \item[$\square$] Je relis mes phrases : sujet + verbe + complément, ponctuation chinoise ，。 correcte.
\end{itemize}
\end{aretenir}

\findecours

\end{document}
